% ============================================================
%  Dispensa di Algorithmic Game Theory
%  Università di Pisa – Laurea Magistrale in Informatica
%  A.A. 2025/26  –  L1–L24 (febbraio – maggio 2026)
% ============================================================
\documentclass[12pt,a4paper,twoside]{book}

\usepackage[utf8]{inputenc}
\usepackage[T1]{fontenc}
% Italian language support (babel)
\usepackage[italian]{babel}
\usepackage{amsmath,amssymb,amsthm,mathtools}
\usepackage{geometry}
\usepackage[colorlinks=true,linkcolor=blue!60!black,citecolor=green!50!black]{hyperref}
\usepackage{enumitem}
\usepackage{booktabs}
\usepackage{array}
\usepackage{float}
\usepackage{xcolor}
\usepackage[most]{tcolorbox}
\usepackage{fancyhdr}
\usepackage{microtype}
\usepackage{lmodern}

\geometry{a4paper, top=3cm, bottom=3cm, inner=3cm, outer=2.5cm, headheight=15pt}

% ---- Theorem environments ----
\theoremstyle{plain}
\newtheorem{theorem}{Teorema}[chapter]
\newtheorem{lemma}[theorem]{Lemma}
\newtheorem{proposition}[theorem]{Proposizione}
\newtheorem{corollary}[theorem]{Corollario}

\theoremstyle{definition}
\newtheorem{definition}[theorem]{Definizione}
\newtheorem{example}[theorem]{Esempio}
\newtheorem{remark}[theorem]{Osservazione}

% ---- Coloured boxes ----
\tcbuselibrary{breakable,skins}
\newtcolorbox{defbox}[1]{
  enhanced, breakable,
  colback=blue!4!white, colframe=blue!50!black,
  fonttitle=\bfseries, title={#1}, left=4pt, right=4pt}
\newtcolorbox{thmbox}[1]{
  enhanced, breakable,
  colback=green!4!white, colframe=green!55!black,
  fonttitle=\bfseries, title={#1}, left=4pt, right=4pt}
\newtcolorbox{exbox}[1]{
  enhanced, breakable,
  colback=orange!4!white, colframe=orange!65!black,
  fonttitle=\bfseries, title={#1}, left=4pt, right=4pt}

% ---- Headers ----
\pagestyle{fancy}
\fancyhf{}
\fancyhead[LE,RO]{\thepage}
\fancyhead[RE]{\nouppercase{\leftmark}}
\fancyhead[LO]{\nouppercase{\rightmark}}
\renewcommand{\headrulewidth}{0.4pt}

% ---- Macros ----
\newcommand{\NE}{\mathrm{NE}}
\newcommand{\BR}{\mathrm{BR}}
\newcommand{\argmax}{\operatorname*{arg\,max}}
\newcommand{\argmin}{\operatorname*{arg\,min}}
\newcommand{\Sh}{\mathrm{Sh}}
\DeclareMathOperator{\supp}{supp}

% ============================================================
\begin{document}

\begin{titlepage}
\centering
\vspace*{3cm}
{\LARGE\bfseries Dispensa di\\[0.4em] Algorithmic Game Theory\par}
\vspace{1.5cm}
{\large Università di Pisa\\
Laurea Magistrale in Informatica\\
A.A.\ 2025/26\par}
\vspace{1cm}
{\normalsize Sintesi delle lezioni L1--L24\\
(Febbraio -- Maggio 2026)\par}
\vspace{0.8cm}
{\small Fonti: slide del corso, \textit{AGT-notes2425}, appunti personali.\par}
\vfill
{\small Settembre 2026 (rev.\ 2)}
\end{titlepage}

\frontmatter
\tableofcontents

\mainmatter

% ============================================================
\chapter{Introduzione alla Teoria dei Giochi (L1)}
% ============================================================

\section{Che cos'è la teoria dei giochi}

La \textbf{teoria dei giochi} (game theory) è la branca della matematica e dell'economia che studia le \emph{interazioni strategiche}: situazioni in cui il risultato per ciascun agente dipende non solo dalle proprie scelte, ma anche dalle scelte altrui. Il termine ``gioco'' è usato in senso ampio e include conflitti politici, aste, mercati, protocolli di rete, algoritmi distribuiti e persino biologia evolutiva.

Il punto di partenza concettuale è semplice: un \textbf{agente razionale} cerca di massimizzare la propria utilità, ma sa che gli altri agenti fanno lo stesso. L'obiettivo della teoria dei giochi è prevedere l'esito di queste interazioni e valutarne l'efficienza.

\section{Breve storia}

\begin{itemize}[leftmargin=2em]
  \item \textbf{1944} -- Von Neumann e Morgenstern pubblicano \emph{Theory of Games and Economic Behavior}, fondando la disciplina moderna e introducendo la teoria dell'utilità attesa.
  \item \textbf{1950} -- John Nash introduce il concetto di \emph{equilibrio di Nash} nella sua tesi di dottorato (Princeton).
  \item \textbf{1994} -- Premio Nobel per l'economia a \textbf{Nash}, \textbf{Harsanyi} e \textbf{Selten} per i contributi alla teoria dei giochi non cooperativi.
  \item \textbf{2005} -- Premio Nobel a \textbf{Aumann} e \textbf{Schelling} per la teoria dei giochi cooperativi e la teoria dell'equilibrio correlated.
  \item \textbf{2012} -- Premio Nobel a \textbf{Shapley} e \textbf{Roth} per la teoria dei mercati e dei matching stabili.
  \item \textbf{2014} -- Premio Nobel a \textbf{Tirole} per l'analisi del potere di mercato e della regolamentazione.
  \item \textbf{2020} -- Premio Nobel a \textbf{Milgrom} e \textbf{Wilson} per la teoria delle aste e i meccanismi di scoperta dei prezzi.
\end{itemize}

\section{Teorema di Zermelo}

\begin{thmbox}{Teorema di Zermelo (1913)}
Ogni gioco finito a due giocatori, a somma zero, con informazione perfetta, senza turni simultanei e senza pareggi (o con pareggio come esito ammissibile), ammette una strategia vincente per uno dei due giocatori (o, nel caso con pareggio, almeno una strategia che garantisca il pareggio).

In termini moderni: ogni gioco finito con informazione perfetta ha un \emph{equilibrio in strategie pure subgame-perfect}.
\end{thmbox}

La dimostrazione si basa sull'induzione all'indietro (\emph{backward induction}): si analizza il gioco a partire dagli stati terminali e si propaga la valutazione ottima verso la radice dell'albero. Poiché l'albero è finito, il processo termina sempre.

\begin{example}[Gioco di Chomp]
Chomp è un gioco combinatoriale su una tavoletta di cioccolato $m\times n$. I giocatori si alternano a ``mordere'' un quadratino (e tutti i quadratini in alto e a destra di esso). Chi è costretto a mangiare il quadratino in basso a sinistra (velenoso) perde. Per il teorema di Zermelo, il primo giocatore ha sempre una strategia vincente per $m,n>1$ (argomento di \emph{strategy stealing}): se il secondo giocatore avesse una strategia vincente, il primo potrebbe ``rubarla'' giocando subito la mossa che riduce la posizione a quella vincente per il secondo -- contraddizione.
\end{example}

\section{Classificazione dei giochi}

\begin{defbox}{Classificazione}
I giochi si classificano lungo diverse dimensioni:
\begin{itemize}
  \item \textbf{Cooperativi vs.\ non cooperativi}: i giocatori possono formare accordi vincolanti (cooperativi) oppure no (non cooperativi). Questo corso tratta principalmente i giochi non cooperativi.
  \item \textbf{A somma zero vs.\ a somma non nulla}: se $\sum_i u_i(s)=0$ per ogni profilo $s$, il gioco è a somma zero (guadagno di uno = perdita dell'altro).
  \item \textbf{Informazione perfetta vs.\ imperfetta}: nella perfetta ogni giocatore conosce la storia completa del gioco; nell'imperfetta esistono \emph{information sets} che raggruppano nodi indistinguibili.
  \item \textbf{Statici vs.\ dinamici}: nei giochi statici (o in forma normale) i giocatori scelgono simultaneamente; nei dinamici (o in forma estesa) le mosse si alternano nel tempo.
  \item \textbf{Finiti vs.\ infiniti}: il numero di strategie (o di turni) è finito o infinito.
\end{itemize}
\end{defbox}

\section*{In sintesi per l'orale}
\addcontentsline{toc}{section}{In sintesi per l'orale (L1)}
\begin{itemize}
  \item Teoria dei giochi $=$ studio di decisioni \emph{interdipendenti}: il risultato di ognuno dipende anche dalle scelte degli altri. Nasce con von Neumann--Morgenstern (1944), Nash (1950), poi Selten, Harsanyi, Aumann, Shapley.
  \item \textbf{Teorema di Zermelo (1913)}: in un gioco finito a due giocatori, a somma zero e a informazione perfetta, o uno dei due ha una strategia vincente, oppure entrambi possono forzare almeno il pareggio. È il primo risultato di ``determinatezza'' e l'antenato della backward induction.
  \item Classificazione da avere in testa: cooperativi contro non cooperativi (si possono firmare accordi vincolanti?), simultanei contro sequenziali, informazione perfetta contro imperfetta, somma zero contro somma generale, finiti contro continui.
  \item Chiusura tipica: ``il corso segue esattamente questa classificazione, prima i non cooperativi L1--L13, poi bargaining e matching, poi i cooperativi TU''.
\end{itemize}

% ============================================================
\chapter{Giochi in Forma Normale (L2)}
% ============================================================

\section{Forma normale}

\begin{defbox}{Gioco in Forma Normale}
Un \textbf{gioco in forma normale} (o \emph{strategic form game}) è una tripla
\[
  G = (N,\, (S_i)_{i\in N},\, (u_i)_{i\in N})
\]
dove:
\begin{itemize}
  \item $N = \{1,\dots,n\}$ è l'insieme finito dei \textbf{giocatori};
  \item $S_i$ è l'insieme (finito o infinito) delle \textbf{strategie} del giocatore $i$;
  \item $u_i : S_1\times\cdots\times S_n \to \mathbb{R}$ è la funzione di \textbf{utilità} (payoff) del giocatore $i$.
\end{itemize}
Si denota $S = \prod_{i\in N} S_i$ lo spazio dei \textbf{profili di strategia} e $S_{-i} = \prod_{j\ne i} S_j$.
\end{defbox}

La notazione $s = (s_i, s_{-i})$ separa la mossa del giocatore $i$ da quelle degli altri.

\section{Forma estesa}

La \textbf{forma estesa} rappresenta giochi sequenziali mediante un albero orientato in cui:
\begin{itemize}
  \item ogni nodo interno appartiene a un giocatore (o alla ``natura'');
  \item ogni arco rappresenta una mossa;
  \item ogni foglia è associata a un vettore di payoff.
\end{itemize}
Gli \emph{information sets} raggruppano nodi indistinguibili per il giocatore che deve muovere.

\section{Esempi classici}

\subsection{Dilemma del Prigioniero}

Due sospettati vengono interrogati separatamente. Ciascuno può \emph{cooperare} (C, non confessare) o \emph{defezionare} (D, confessare).

\begin{center}
\begin{tabular}{c|cc}
I/II & C & D \\\hline
C & $-2,-2$ & $-7,\phantom{-}0$ \\
D & $\phantom{-}0,-7$ & $-5,-5$
\end{tabular}
\end{center}

Confessare è \emph{strettamente dominante} per entrambi i giocatori. L'unico NE è $(D,D)$, ma $(C,C)$ sarebbe socialmente ottimale. Il dilemma del prigioniero cattura il conflitto tra razionalità individuale ed efficienza collettiva.

\subsection{Battle of Sexes (BoS)}

Una coppia deve scegliere tra \emph{football} (F) e \emph{danza} (D). Lui preferisce F, lei D, ma entrambi preferiscono stare insieme.

\begin{center}
\begin{tabular}{c|cc}
lui/lei & F & D \\\hline
F & $2,1$ & $0,0$ \\
D & $0,0$ & $1,2$
\end{tabular}
\end{center}

Ci sono due NE puri: $(F,F)$ e $(D,D)$.

\subsection{Hawk-Dove}

Due animali si contendono una risorsa di valore $V=6$. Chi sceglie \emph{Hawk} (H) combatte; chi sceglie \emph{Dove} (D) retrocede. Il costo di un combattimento tra due Hawk è $C=8$ (diviso equamente).

\begin{center}
\begin{tabular}{c|cc}
I/II & H & D \\\hline
H & $-1,-1$ & $4,0$ \\
D & $0,4$ & $2,2$
\end{tabular}
\end{center}

\subsection{Rock-Paper-Scissors (RPS)}

Gioco a somma zero $3\times 3$: Sasso (R) batte Forbici (S), Forbici batte Carta (P), Carta batte Sasso. Non esiste NE in strategie pure.

\subsection{Lizard-Spock}

Estensione $5\times 5$ (Sasso, Carta, Forbici, Lucertola, Spock) con le stesse proprietà di simmetria. Anche qui l'unico NE è la distribuzione uniforme $(1/5,\dots,1/5)$.

\section{Sequential Allocation}

Due giocatori si alternano (inizia il giocatore 1) nel selezionare oggetti da un insieme ordinato $\{1,\dots,m\}$. Ogni giocatore preferisce il valore più alto disponibile. La forma estesa è un albero binario di profondità $m$.

\section{Duopolio di Cournot con bene indivisibile}

Due imprese producono quantità $x_1,x_2\in\{0,1,\dots,6\}$ di un bene omogeneo. La funzione di utilità è
\[
  u_i(x_1,x_2) = x_i \cdot \max\{T-(x_1+x_2),\,0\} - c\,x_i
\]
con $T=10$ (prezzo di riserva) e $c=3$ (costo unitario). La matrice dei payoff $7\times 7$ si ottiene calcolando $u_i$ per ogni coppia $(x_1,x_2)$. Si noti che $u_i > 0$ richiede $x_i(T-x_1-x_2-c)>0$, cioè la domanda residua deve superare il costo.

\section*{In sintesi per l'orale}
\addcontentsline{toc}{section}{In sintesi per l'orale (L2)}
\begin{itemize}
  \item Forma normale $G=(N,(S_i),(u_i))$ con $u_i:S_1\times\cdots\times S_n\to\mathbb{R}$: l'utilità dipende dal \emph{profilo completo}, non solo dalla propria strategia. È l'errore formale più comune.
  \item Forma estesa $=$ albero con nodi assegnati ai giocatori (o alla natura), archi $=$ mosse, foglie $=$ vettori di payoff, insiemi di informazione $=$ nodi indistinguibili per chi muove.
  \item \textbf{Strategia pura in forma estesa} $=$ \emph{funzione} che assegna un'azione a ogni insieme di informazione controllato dal giocatore. Quindi $|S_i|=\prod$ (numero di azioni in ciascun suo insieme di informazione): un \emph{prodotto}, non la somma degli archi. Esempio: due nodi con due azioni ciascuno danno 4 strategie pure.
  \item Esempi canonici da saper scrivere in 20 secondi: dilemma del prigioniero (dominanza, inefficienza), BoS (due equilibri, coordinamento), Hawk-Dove (anti-coordinamento), RPS (nessun equilibrio puro).
  \item Esempi ``a numeri'': sequential allocation (albero binario di profondità $m$) e duopolio di Cournot con bene indivisibile, $u_i(x_1,x_2)=x_i\max\{T-x_1-x_2,0\}-cx_i$ con $T=10$, $c=3$.
\end{itemize}

% ============================================================
\chapter{Utilità Attesa e Nash Equilibrium (L3)}
% ============================================================

\section{Lotterie e utilità attesa}

\begin{defbox}{Lotteria}
Una \textbf{lotteria} su un insieme finito $O$ di esiti è una distribuzione di probabilità $p\in\Delta(O)$, cioè $p(o)\ge 0$ e $\sum_{o\in O}p(o)=1$.
\end{defbox}

\begin{thmbox}{Teorema di von Neumann--Morgenstern (1944)}
Sotto quattro assiomi (completezza, transitività, continuità, indipendenza), esistono valori reali $u(o)$ per ogni esito $o$ tali che, per ogni coppia di lotterie $p,q$:
\[
  p \succcurlyeq q \iff \mathbb{E}_p[u] \ge \mathbb{E}_q[u]
\]
cioè le preferenze si rappresentano massimizzando l'\textbf{utilità attesa} $\sum_o p(o)\,u(o)$.
\end{thmbox}

Il teorema giustifica l'uso di funzioni di payoff a valori reali nei giochi: i giocatori si comportano come se massimizzassero un'utilità attesa.

\section{Nash Equilibrium}

\begin{defbox}{Nash Equilibrium (NE)}
Dato un gioco $G=(N,(S_i),(u_i))$, un profilo $s^* = (s^*_1,\dots,s^*_n)\in S$ è un \textbf{equilibrio di Nash} se per ogni giocatore $i\in N$:
\[
  s^*_i \in \argmax_{s_i\in S_i} u_i(s_i,\, s^*_{-i})
\]
Equivalentemente, nessun giocatore può aumentare il proprio payoff deviando unilateralmente: per ogni $i$ e ogni $s_i\in S_i$,
\[
  u_i(s^*_i,\, s^*_{-i}) \ge u_i(s_i,\, s^*_{-i}).
\]
\end{defbox}

L'intuizione è quella di un \emph{punto fisso}: a Nash equilibrium ogni giocatore sta già giocando la best response alle strategie degli avversari, per cui nessuno ha incentivo a cambiare.

\subsection{NE negli esempi}

\paragraph{Dilemma del Prigioniero.} L'unico NE è $(D,D)$: $D$ domina strettamente $C$ per entrambi i giocatori.

\paragraph{Battle of Sexes.} I due NE puri sono $(F,F)$ e $(D,D)$. Esiste anche un NE misto (vedi Capitolo~\ref{ch:mixed}).

\paragraph{Hawk-Dove.} Esistono due NE puri: $(H,D)$ e $(D,H)$. Esiste un NE misto con probabilità $p^*=V/C=3/4$ su $H$.

\paragraph{RPS.} Non esistono NE puri. L'unico NE misto è la distribuzione uniforme $(1/3,1/3,1/3)$ per entrambi i giocatori.

\section{Giochi zero-sum simmetrici}

In un gioco a \emph{due giocatori a somma zero}, $u_2(s) = -u_1(s)$ per ogni $s$. Un gioco è \textbf{simmetrico} se la matrice di payoff del giocatore 1 è uguale alla trasposta di quella del giocatore 2. In un gioco zero-sum simmetrico l'equilibrio misto assegna la stessa distribuzione a entrambi i giocatori, come nel caso di RPS.

\section*{In sintesi per l'orale}
\addcontentsline{toc}{section}{In sintesi per l'orale (L3)}
\begin{itemize}
  \item \textbf{vNM}: sotto completezza, transitività, continuità e indipendenza le preferenze sulle lotterie si rappresentano con l'utilità attesa $\sum_o p(o)u(o)$. La $u$ è unica a meno di \emph{trasformazioni affini positive} $au+b$, $a>0$: serve perché le strategie miste trasformano i payoff in lotterie.
  \item \textbf{NE}: $u_i(s_i^*,s_{-i}^*)\ge u_i(s_i,s_{-i}^*)$ per ogni $i$ e ogni $s_i$. Disuguaglianza \emph{debole} e deviazione \emph{unilaterale}: entrambe le parole vanno dette.
  \item Lettura da punto fisso: al NE ognuno sta già giocando una best response, cioè $s^*\in\BR(s^*)$. È il ponte verso il teorema di esistenza (L6).
  \item NE degli esempi: PD $(D,D)$; BoS $(F,F)$ e $(D,D)$ più il misto; Hawk-Dove i due asimmetrici più il misto; RPS solo l'uniforme.
  \item Somma zero simmetrico ($A=-A^{\mathsf T}$): il valore è $V=0$ e i due giocatori usano la stessa strategia ottima.
\end{itemize}

% ============================================================
\chapter{Dominanza Strategica (L4)}
% ============================================================

\section{Calcolo brute-force degli NE}

Per un gioco finito $2\times m$ o $n\times m$ si costruiscono le \emph{liste di best response}: per ogni strategia $s_j$ del giocatore $j$, si individuano tutte le best response del giocatore $i$. Un profilo $(s_i^*,s_j^*)$ è NE se e solo se $s_i^*\in\BR_i(s_j^*)$ e $s_j^*\in\BR_j(s_i^*)$: si costruisce la lista $L_1$ dei profili dove $i$ risponde ottimamente, la lista $L_2$ dove $j$ risponde ottimamente, e si calcola $L_1\cap L_2$.

\section{Strategie dominanti}

\begin{defbox}{Dominanza Stretta e Debole}
La strategia $s_i\in S_i$ \textbf{domina strettamente} $t_i\in S_i$ se
\[
  u_i(s_i,\, s_{-i}) > u_i(t_i,\, s_{-i})\quad \forall\,s_{-i}\in S_{-i}.
\]
La domina \textbf{debolmente} se la disuguaglianza è $\ge$ con almeno uno stretto.

Una strategia è \textbf{dominante} (strettamente) se domina strettamente tutte le altre.
\end{defbox}

\begin{remark}
Se esiste una strategia dominante per ogni giocatore, il profilo di strategie dominanti è l'unico NE del gioco.
\end{remark}

\begin{example}[PD]
Nel Dilemma del Prigioniero, la strategia $D$ domina strettamente $C$ per entrambi: $0>-2$ (se l'avversario coopera) e $-5>-7$ (se l'avversario defeziona). Dunque $(D,D)$ è l'unico NE.
\end{example}

\section{Eliminazione iterata di strategie strettamente dominate (IESDS)}

\begin{defbox}{IESDS}
Dato il gioco $G^0=G$, si definisce la sequenza di giochi ridotti $G^0\supseteq G^1\supseteq\cdots$:
\[
  S_i^0 = S_i,\quad S_i^{k+1} = S_i^k \setminus \{s_i\in S_i^k : \exists\, t_i\in S_i^k \text{ che domina strettamente } s_i \text{ in } G^k\}.
\]
Il processo termina quando $S_i^{k+1}=S_i^k$ per ogni $i$. Il profilo risultante $S^\infty = \prod_i S_i^\infty$ è l'\textbf{insieme IESDS}.
\end{defbox}

\begin{proposition}
Ogni NE del gioco originale sopravvive all'IESDS. Se l'IESDS porta a un unico profilo, esso è l'unico NE del gioco.
\end{proposition}

\textit{Dimostrazione (sketch).} Supponiamo per assurdo che un NE $s^*$ venga eliminato al passo $k$: allora esiste $t_i$ che domina strettamente $s^*_i$ in $G^k$, ma in $G^k$ esistono ancora le strategie degli avversari in $s^*$, il che contraddice la proprietà di best response di $s^*_i$. \qed

\section{Strategie debolmente dominate}

L'IESDS con dominanza \emph{debole} (IEWDS) è più delicata: l'ordine di eliminazione può influenzare l'esito finale. È importante verificare i risultati caso per caso.

\section*{In sintesi per l'orale}
\addcontentsline{toc}{section}{In sintesi per l'orale (L4)}
\begin{itemize}
  \item Dominanza \emph{stretta}: $u_i(s_i,s_{-i})>u_i(t_i,s_{-i})$ per ogni $s_{-i}$; \emph{debole}: $\ge$ con almeno una disuguaglianza stretta. Nella definizione completa il dominante può essere una strategia \emph{mista}.
  \item Se ogni giocatore ha una strategia dominante, quel profilo è l'unico NE (PD).
  \item \textbf{IESDS}: elimina iterativamente le strettamente dominate. Preserva \emph{tutti} i NE, e il risultato non dipende dall'ordine di eliminazione. Se resta un solo profilo, quello è l'unico NE.
  \item Sketch: se un NE venisse eliminato al passo $k$, esisterebbe $t_i$ che lo domina strettamente contro \emph{tutte} le strategie ancora presenti, incluse quelle del NE: contraddice la best response.
  \item \textbf{IEWDS}: con la dominanza debole l'ordine conta e si possono distruggere equilibri. Se lo si usa, va detto esplicitamente in che ordine.
\end{itemize}

% ============================================================
\chapter{Eliminazione Iterata e Strategie di Sicurezza (L5)}
% ============================================================

\section{Razionalizzabilità e IENBR}

\begin{defbox}{Never Best Response (NBR)}
La strategia $s_i$ è una \textbf{never best response} (NBR) se non esiste alcuna credenza $\sigma_{-i}\in\Delta(S_{-i})$ sulle strategie degli avversari tale che $s_i\in\BR_i(\sigma_{-i})$.
\end{defbox}

\begin{defbox}{IENBR}
Il processo di \textbf{Iterated Elimination of Never Best Responses} elimina iterativamente le NBR finché non ne rimangono.
\end{defbox}

\begin{proposition}
Le strategie che sopravvivono all'IENBR coincidono con le strategie \textbf{razionalizzabili} (Bernheim 1984, Pearce 1984): quelle per cui esiste un sistema di credenze consistente con la razionalità comune.
\end{proposition}

IESDS implica IENBR (ogni strategia strettamente dominata è NBR), ma non viceversa.

\section{Strategie di sicurezza (maximin)}

\begin{defbox}{Valore di Sicurezza}
Il \textbf{valore di sicurezza} (security level) del giocatore $i$ è
\[
  v_i^* = \max_{s_i\in S_i}\; \min_{s_{-i}\in S_{-i}} u_i(s_i,\, s_{-i}).
\]
La strategia che lo raggiunge è la \textbf{strategia di sicurezza} (maximin strategy).
\end{defbox}

In un gioco a due giocatori, la strategia maximin garantisce a $i$ almeno $v_i^*$ indipendentemente dall'avversario. Simmetricamente, il \emph{minimax value} di $j$ contro $i$ è $\min_{s_j}\max_{s_i} u_i(s_i,s_j) \ge v_i^*$.

\section{Gioco di Hotelling}

\begin{example}[Hotelling (1929)]
Due politici scelgono la propria posizione $x_1,x_2\in[0,100]$ su uno spettro politico. Gli elettori votano per il candidato più vicino (distribuzione uniforme). Ogni candidato massimizza la propria quota di voti. L'unico NE è $x_1^*=x_2^*=50$ (posizione mediana): se un candidato si discosta, l'altro può catturare più voti avvicinandosi al centro.
\end{example}

Il gioco di Hotelling illustra la forza dell'IESDS: tutte le posizioni $<25$ e $>75$ vengono eliminate al primo passo, poi si procede iterativamente fino all'unico sopravvissuto $\{50\}$.

\section*{In sintesi per l'orale}
\addcontentsline{toc}{section}{In sintesi per l'orale (L5)}
\begin{itemize}
  \item \textbf{NBR}: $s_i$ è never best response se \emph{non esiste alcuna} credenza mista $\sigma_{-i}\in\Delta(S_{-i})$ sulle strategie altrui a cui $s_i$ risponda ottimamente. Il quantificatore su \emph{tutte} le credenze è il cuore della definizione.
  \item IENBR $=$ eliminazione iterata delle NBR; ciò che sopravvive $=$ strategie \textbf{razionalizzabili} (Bernheim, Pearce 1984). Ogni strettamente dominata è NBR, quindi IESDS $\subseteq$ IENBR; a due giocatori le due nozioni coincidono.
  \item \textbf{Strategie di sicurezza}: $w_i(x_i)=\min_{x_{-i}}u_i(x_i,x_{-i})$, livello di sicurezza $v_i=\max_{x_i}w_i(x_i)$, e la massimizzante è la strategia di sicurezza (maximin). Garantisce $v_i$ comunque giochino gli altri.
  \item Sicurezza e equilibrio sono cose diverse: in RPS ogni strategia è di sicurezza; in Hawk-Dove la sicurezza è $(D,D)$, che non è un NE; nel Cournot continuo la sicurezza è produrre $0$.
  \item \textbf{Hotelling}: due candidati su $[0,100]$, elettori uniformi, si vota il più vicino. IESDS elimina progressivamente gli estremi e resta l'unico NE $x_1^*=x_2^*=50$ (teorema dell'elettore mediano).
\end{itemize}

% ============================================================
\chapter{Strategie Miste e Teoremi di Esistenza (L6)}
\label{ch:mixed}
% ============================================================

\section{Strategie miste ed estensione mista}

\begin{defbox}{Strategia mista}
Una \textbf{strategia mista} $\sigma_i\in\Delta(S_i)$ del giocatore $i$ è una
distribuzione di probabilità sulle sue strategie pure:
\[
  \Delta(S_i)=\Bigl\{\sigma_i:S_i\to[0,1]\ :\ \textstyle\sum_{x_i\in S_i}\sigma_i(x_i)=1\Bigr\}.
\]
Equivale a un vettore di $m_i=|S_i|$ componenti non negative con somma $1$ (il
\emph{simplesso} $\Delta_{m_i}$). Il \textbf{supporto} è
$\supp(\sigma_i)=\{x_i : \sigma_i(x_i)>0\}$; una strategia pura è la mista
degenere che mette probabilità $1$ su un solo punto.
\end{defbox}

L'utilità attesa nel profilo misto $\sigma=(\sigma_1,\dots,\sigma_n)$ è
\[
  h_i(\sigma)=\sum_{x\in S} u_i(x_1,\dots,x_n)\,\sigma_1(x_1)\cdots\sigma_n(x_n)
  =\sum_{x\in S} u_i(x)\,\sigma_i(x_i)\,\sigma_{-i}(x_{-i}),
  \qquad
  \sigma_{-i}(x_{-i})=\prod_{j\ne i}\sigma_j(x_j).
\]
La proprietà da cui discende tutto il capitolo è che $h_i$ è \textbf{multilineare}:
lineare in $\sigma_i$ una volta fissate le altre.

\begin{defbox}{Estensione mista}
Dato un gioco finito $G=(N,(S_i),(u_i))$, il gioco
$G_{\mathrm{me}}=(N,(\Delta(S_i)),(h_i))$ è l'\textbf{estensione di $G$ a strategie
miste}. Per definizione, gli equilibri di Nash in strategie miste di $G$ sono gli
equilibri di Nash di $G_{\mathrm{me}}$.
\end{defbox}

\begin{proposition}
Se $x^*\in S$ è un NE di $G$, il profilo di miste degeneri corrispondente è un NE
di $G_{\mathrm{me}}$. Il viceversa è falso: RPS ha un NE misto e nessun NE puro.
\end{proposition}

\section{Come si calcola un NE misto: il principio di indifferenza}

\begin{thmbox}{Caratterizzazione per indifferenza}
Il profilo $\sigma^*$ è un NE in strategie miste se e solo se, per ogni giocatore
$i$ e ogni $x_i\in S_i$:
\[
  x_i\in\supp(\sigma^*_i)\ \Longrightarrow\
  h_i(x_i,\sigma^*_{-i})=\max_{t_i\in S_i} h_i(t_i,\sigma^*_{-i}).
\]
Cioè: \emph{tutte} le pure nel supporto danno lo stesso payoff atteso, pari al
massimo; quelle fuori dal supporto danno un payoff minore o uguale.
\end{thmbox}

\textit{Perché vale.} $h_i(\sigma_i,\sigma_{-i})=\sum_{x_i}\sigma_i(x_i)\,
h_i(x_i,\sigma_{-i})$ è una media pesata dei payoff delle pure: una media è
massima se e solo se tutto il peso sta sui massimi. \qed

\begin{defbox}{Ricetta operativa per un NE misto $2\times 2$}
\begin{enumerate}[leftmargin=*]
  \item Chiama $p$ la probabilità con cui \emph{il giocatore I} gioca la sua prima
        strategia e $q$ quella con cui \emph{il giocatore II} gioca la sua.
  \item Scrivi i due payoff attesi del giocatore I contro $q$ e imponi che siano
        \emph{uguali}: ottieni $q$. \textbf{Attenzione}: l'indifferenza del
        giocatore I determina la mista dell'\emph{avversario}, non la sua.
  \item Simmetricamente, l'indifferenza del giocatore II determina $p$.
  \item Verifica che $p,q\in(0,1)$: se esce fuori range, il NE completamente misto
        non esiste e vanno esaminati i supporti più piccoli (cioè i NE puri).
\end{enumerate}
\end{defbox}

\begin{example}[Battle of Sexes, NE misto]
Con lui che gioca $F$ con probabilità $p$ e lei $F$ con probabilità $q$:
\[
  h_{\text{lui}}(p,q)=2pq+(1-p)(1-q),\qquad
  h_{\text{lei}}(p,q)=pq+2(1-p)(1-q).
\]
Indifferenza di lui: $2q=1-q\Rightarrow q^*=1/3$.
Indifferenza di lei: $p=2(1-p)\Rightarrow p^*=2/3$.
Il NE misto è $\bigl((2/3,1/3),(1/3,2/3)\bigr)$ con payoff attesi
$h_{\text{lui}}=h_{\text{lei}}=2/3$, \emph{peggiori} di entrambi gli equilibri puri.
\end{example}

\begin{example}[Stesso conto con le corrispondenze di best response]
La lettura equivalente, quella disegnata a lezione, passa dalle best response:
\[
  R_{\text{lei}}(p)=\begin{cases}1 & p>2/3\\ [0,1] & p=2/3\\ 0 & p<2/3\end{cases}
  \qquad
  R_{\text{lui}}(q)=\begin{cases}1 & q>1/3\\ [0,1] & q=1/3\\ 0 & q<1/3\end{cases}
\]
I NE sono le intersezioni dei due grafici nel quadrato $[0,1]^2$:
$(p^*,q^*)=(1,1)$, $(0,0)$ e $(2/3,1/3)$. Le due curve ``a gradino'' si incrociano
tre volte: due agli angoli (i puri) e una nel punto di indifferenza (il misto).
Cambiando i payoff da $2$ e $1$ a $3$ e $2$, gli equilibri puri non si muovono ma
quello misto sì: è il modo rapido per mostrare che il misto dipende dai
\emph{valori}, i puri solo dall'ordine.
\end{example}

\begin{example}[Hawk-Dove]
Con $(H,H)=(-2,-2)$, $(H,D)=(2,0)$, $(D,H)=(0,2)$, $(D,D)=(1,1)$: l'avversario
gioca $H$ con probabilità $q$ tale che $-2q+2(1-q)=0\cdot q+1\cdot(1-q)$, cioè
$2-4q=1-q$ e $q^*=1/3$. Payoff atteso $2/3$ a testa.
\end{example}

\section{Giochi strettamente competitivi e a somma zero}

\begin{defbox}{Strettamente competitivo e somma zero}
Un gioco a due giocatori è \textbf{strettamente competitivo} se per ogni coppia di
profili $u_1$ cresce esattamente quando $u_2$ cala. È a \textbf{somma zero} se
$u_1(x)+u_2(x)=0$ per ogni $x$. Ogni gioco strettamente competitivo si riscrive
come gioco a somma zero sostituendo $u_2$ con $-u_1$: si tiene quindi una sola
funzione $u=u_1$, che I massimizza e II minimizza.
\end{defbox}

RPS è strettamente competitivo, Hawk-Dove no.

\begin{defbox}{Livelli di sicurezza e valore}
\[
  \underline{v}=\max_{x_1\in S_1}\min_{x_2\in S_2} u(x_1,x_2)
  \qquad
  \overline{v}=\min_{x_2\in S_2}\max_{x_1\in S_1} u(x_1,x_2).
\]
Vale sempre $\underline{v}\le\overline{v}$. Il gioco \textbf{ha valore} se
$\underline{v}=\overline{v}$; in tal caso le strategie di sicurezza si chiamano
\textbf{strategie minimax}.
\end{defbox}

\begin{example}[Un gioco a somma zero $3\times 4$ con valore in strategie pure]
\begin{center}
\begin{tabular}{c|cccc|c}
I/II & 1 & 2 & 3 & 4 & $\min$ \\\hline
1 & 4 & 7 & 2 & 3 & 2 \\
2 & 1 & 5 & 2 & 4 & 1 \\
3 & 3 & 4 & \textbf{3} & 6 & \textbf{3} \\\hline
$\max$ & 4 & 7 & \textbf{3} & 6 &
\end{tabular}
\end{center}
$\underline{v}=3$ (riga 3) e $\overline{v}=3$ (colonna 3): il gioco ha valore $3$ e
$(3,3)$ è un NE. Costo computazionale: $m_1+m_2+2$ scansioni di array, cioè
tempo lineare nella matrice. In RPS invece $\underline{v}=-1<1=\overline{v}$: in
strategie pure il valore non esiste.
\end{example}

\begin{thmbox}{Punto di sella e Nash $\equiv$ minimax}
$(x_1^*,x_2^*)$ è un NE se e solo se è un \textbf{punto di sella} di $u$:
\[
  u(x_1,x_2^*)\ \le\ u(x_1^*,x_2^*)\ \le\ u(x_1^*,x_2)
  \qquad\forall x_1\in S_1,\ \forall x_2\in S_2.
\]
Di conseguenza: (i) se esiste un NE, il gioco ha valore $u(x_1^*,x_2^*)$ e le due
componenti sono strategie minimax; (ii) se il gioco ha valore, ogni coppia di
strategie minimax è un NE. \emph{Esistenza di NE $\equiv$ esistenza del valore.}
\end{thmbox}

\begin{proposition}[Intercambiabilità]
Se $(x_1^*,x_2^*)$ e $(x_1^\circ,x_2^\circ)$ sono due NE di un gioco a somma zero,
allora $u(x_1^*,x_2^*)=u(x_1^\circ,x_2^\circ)$ e anche $(x_1^*,x_2^\circ)$ e
$(x_1^\circ,x_2^*)$ sono NE. Gli equilibri si ``mescolano'': una proprietà che
nei giochi generali non vale (in BoS mescolare i due equilibri dà $(F,D)$, che non
è un equilibrio).
\end{proposition}

\begin{remark}[Zero-sum simmetrico]
Se il gioco è a somma zero e simmetrico ($A=-A^{\mathsf T}$, come RPS), allora il
valore è $V=0$: nessuno dei due può garantirsi più dell'altro, e i due giocatori
usano la stessa strategia ottima.
\end{remark}

\section{Teorema Minimax di von Neumann}

\begin{thmbox}{Teorema Minimax (von Neumann, 1928)}
In ogni gioco finito a due giocatori a somma zero, l'estensione mista ha valore:
\[
  \max_{p\in\Delta_{m_1}}\ \min_{q\in\Delta_{m_2}} p^{\mathsf T}\!Aq
  \;=\;
  \min_{q\in\Delta_{m_2}}\ \max_{p\in\Delta_{m_1}} p^{\mathsf T}\!Aq
  \;=:\;V .
\]
Le strategie miste ottimali dei due giocatori formano gli NE del gioco, e tutti
gli NE danno lo stesso payoff $V$.
\end{thmbox}

Il punto concettuale: in strategie \emph{pure} il valore può non esistere (RPS),
in strategie \emph{miste} esiste sempre. Randomizzare serve esattamente a chiudere
il divario $\underline{v}<\overline{v}$.

\section{Minimax e programmazione lineare}

Questa è la dimostrazione ``algoritmica'' del teorema minimax, ed è il legame che
il corso usa per calcolare gli equilibri. Sia $A\in\mathbb{R}^{m_1\times m_2}$ e
$h(p,q)=p^{\mathsf T}\!Aq$. Contro una strategia mista $p$ fissata, il minimizzatore
può sempre rispondere con una pura, quindi
\[
  w_1(p)=\min_{q\in\Delta_{m_2}} h(p,q)=\min_{\ell=1,\dots,m_2}\,(A^{\mathsf T}p)_\ell,
  \qquad
  w_2(q)=\max_{p\in\Delta_{m_1}} h(p,q)=\max_{k=1,\dots,m_1}\,(Aq)_k .
\]
Massimizzare un minimo di funzioni lineari è un problema di programmazione lineare:
si introduce una variabile scalare che fa da ``soglia''.

\begin{center}
\begin{tabular}{p{0.44\textwidth} p{0.44\textwidth}}
\textbf{Giocatore I} ($\max w_1$) & \textbf{Giocatore II} ($\min w_2$)\\[2pt]
$\max\ u$ & $\min\ v$ \\
\quad $u\le (A^{\mathsf T}p)_\ell,\quad \ell=1,\dots,m_2$ &
\quad $v\ge (Aq)_k,\quad k=1,\dots,m_1$ \\
\quad $p_1+\cdots+p_{m_1}=1$ & \quad $q_1+\cdots+q_{m_2}=1$ \\
\quad $p\ge 0$ & \quad $q\ge 0$
\end{tabular}
\end{center}

\begin{thmbox}{Dualità forte $\equiv$ teorema minimax}
I due programmi lineari qui sopra sono l'uno il \textbf{duale} dell'altro. Entrambi
sono ammissibili e limitati, quindi per il teorema di dualità forte hanno lo stesso
valore ottimo: $\max w_1=\min w_2=V$. Questo \emph{è} il teorema minimax, e dà anche
l'algoritmo: si calcolano gli equilibri di un gioco a somma zero risolvendo un LP
(quindi in tempo polinomiale).
\end{thmbox}

\begin{example}[RPS via LP]
Con $A$ la matrice di RPS,
$h(p,q)=q_1(p_2-p_3)+q_2(p_3-p_1)+q_3(p_1-p_2)$, da cui
$w_1(p)=\min\{p_2-p_3,\ p_3-p_1,\ p_1-p_2\}$. I tre termini hanno somma nulla,
quindi il minimo è $\le 0$ per ogni $p$, con uguaglianza solo quando sono tutti e
tre nulli, cioè $p_1=p_2=p_3=1/3$. Dunque $V=0$ e la strategia ottima è uniforme;
lo stesso vale per II. Il programma lineare corrispondente è
$\max u$ con $u\le p_2-p_3$, $u\le p_3-p_1$, $u\le p_1-p_2$, $p\in\Delta_3$.
\end{example}

\section{Esistenza degli equilibri: Nikaido--Isoda}

\begin{thmbox}{Teorema di Nikaido--Isoda (1955)}
Sia $G=(N,(S_i),(u_i))$ tale che per ogni $i$:
\begin{enumerate}
  \item $S_i\subseteq\mathbb{R}^{m_i}$ è \textbf{compatto} e \textbf{convesso};
  \item $u_i$ è \textbf{continua} in $s$;
  \item $u_i(\cdot,s_{-i})$ è \textbf{quasi-concava} in $s_i$ per ogni $s_{-i}$.
\end{enumerate}
Allora $G$ ammette almeno un equilibrio di Nash.
\end{thmbox}

\begin{corollary}[Nash, 1950]
Ogni gioco finito ha almeno un NE in strategie miste.
\end{corollary}
\textit{Dimostrazione.} Si applica Nikaido--Isoda a $G_{\mathrm{me}}$: i simplessi
$\Delta(S_i)$ sono compatti e convessi, e le utilità attese $h_i$ sono
multilineari, quindi continue e (quasi-)concave in $\sigma_i$. \qed

\begin{remark}[Come si dimostra Nikaido--Isoda, in una riga]
Si considera la corrispondenza di best response congiunta
$\BR(\sigma)=\prod_i \BR_i(\sigma_{-i})$ definita sul convesso compatto
$\prod_i\Delta(S_i)$. Le ipotesi garantiscono che $\BR$ abbia valori non vuoti e
convessi e grafo chiuso: il teorema di punto fisso di \textbf{Kakutani} dà
$\sigma^*\in\BR(\sigma^*)$, che per definizione è un NE. La funzione di
Nikaido--Isoda $\Psi(\sigma,\tau)=\sum_i\bigl[h_i(\tau_i,\sigma_{-i})-h_i(\sigma)\bigr]$
rende la cosa operativa: $\sigma$ è NE se e solo se $\max_\tau\Psi(\sigma,\tau)=0$.
\end{remark}

\begin{remark}[Le tre vie all'esistenza, da tenere distinte all'orale]
\begin{itemize}
  \item \emph{Generale}: Nikaido--Isoda (Kakutani) $\Rightarrow$ esistenza nei
        giochi finiti misti e nei giochi continui quasi-concavi.
  \item \emph{Somma zero}: dualità forte dell'LP $\Rightarrow$ teorema minimax
        (costruttiva e polinomiale).
  \item \emph{Potential games} (L7--L8): massimo del potenziale $\Rightarrow$ NE in
        strategie \emph{pure}, senza punto fisso.
\end{itemize}
\end{remark}

\section*{In sintesi per l'orale}
\addcontentsline{toc}{section}{In sintesi per l'orale (L6)}
\begin{itemize}
  \item Mista $=$ distribuzione su $S_i$; $h_i$ è multilineare; estensione mista
        $G_{\mathrm{me}}$. NE misti di $G$ $=$ NE di $G_{\mathrm{me}}$.
  \item \textbf{Indifferenza}: nel supporto tutte le pure danno lo stesso payoff
        atteso. L'indifferenza di un giocatore determina la mista
        \emph{dell'altro}. BoS: $q^*=1/3$, $p^*=2/3$, payoff $2/3$ (peggio dei puri).
  \item Somma zero: $\underline v\le\overline v$; il gioco ha valore se sono uguali;
        NE $\equiv$ punto di sella $\equiv$ coppia di strategie minimax; gli
        equilibri sono intercambiabili. Simmetrico $\Rightarrow V=0$.
  \item \textbf{Minimax (von Neumann)}: in miste $\max_p\min_q p^{\mathsf T}Aq=
        \min_q\max_p p^{\mathsf T}Aq$. In pure può fallire (RPS: $-1<1$).
  \item \textbf{LP}: $\max u$ con $u\le (A^{\mathsf T}p)_\ell$, $p\in\Delta$, e il
        suo duale $\min v$ con $v\ge (Aq)_k$, $q\in\Delta$: dualità forte $=$
        teorema minimax $=$ algoritmo polinomiale.
  \item \textbf{Esistenza}: Nikaido--Isoda (compatto, convesso, continua,
        quasi-concava) via Kakutani sulla best response congiunta; corollario di
        Nash per i giochi finiti.
  \item Esempio pronto: la matrice $3\times 4$ con $\underline v=\overline v=3$ per
        il caso ``valore in pure''; RPS per il caso ``serve randomizzare''.
\end{itemize}

% ============================================================
\chapter{Best Response Dynamics e Potential Games (L7)}
% ============================================================

\section{Best response dynamics}

\begin{defbox}{Best Response}
La \textbf{best response} del giocatore $i$ al profilo $x_{-i}$ è
\[
  \BR_i(x_{-i}) = \argmax_{s_i\in S_i} u_i(s_i,\, x_{-i}).
\]
\end{defbox}

\begin{defbox}{Best Response Dynamics (BRD)}
Partendo da un profilo iniziale $x^0$, ad ogni passo $k$ uno (o più) giocatori aggiornano la propria strategia con una best response: $x^{k+1}_i\in\BR_i(x^k_{-i})$. Il processo si dice \textbf{convergente} se raggiunge un punto fisso (che è necessariamente un NE).
\end{defbox}

\subsection{Algoritmo di Jacobi (sincrono)}
Tutti i giocatori aggiornano \emph{simultaneamente} al passo $k$:
\begin{enumerate}
  \item $x^0\in S$, $k=0$.
  \item Per ogni $i$: $x_i^{k+1}\in\BR_i(x^k_{-i})$ (se $x_i^k\in\BR_i(x^k_{-i})$, si lascia $x_i^{k+1}=x_i^k$).
  \item Se $x^{k+1}=x^k$: STOP (NE trovato).
  \item $k\leftarrow k+1$, vai al passo 2.
\end{enumerate}

\subsection{Algoritmo di Gauss-Seidel (asincrono)}
Un solo giocatore per volta aggiorna la propria strategia. L'ordine può essere deterministico o casuale.

\begin{example}[Duopolio di Cournot -- traiettoria BRD]
Con $T=10$, $c=3$, quantità $x_i\in\{0,\dots,6\}$, partendo da $(x_1,x_2)=(2,5)$: la BR di $x_1$ contro $x_2=5$ è $1$ (troppa offerta dall'avversario); poi la BR di $x_2$ contro $x_1=1$ è $2$; poi la BR di $x_1$ contro $x_2=2$ è $3$; e così via finché si raggiunge l'NE $(2,2)$.
\end{example}

\section{Potential games}

\begin{defbox}{Potential Game (Monderer--Shapley, 1996)}
Un gioco $G$ ha un \textbf{potenziale esatto} $P:S\to\mathbb{R}$ se per ogni $i$, ogni $s_{-i}$ e ogni $s_i,t_i\in S_i$:
\[
  P(s_i,s_{-i}) - P(t_i,s_{-i}) = u_i(s_i,s_{-i}) - u_i(t_i,s_{-i}).
\]
Ha un \textbf{potenziale ordinale} $P$ se l'uguaglianza è sostituita da una condizione sul segno:
\[
  P(s_i,s_{-i}) > P(t_i,s_{-i}) \iff u_i(s_i,s_{-i}) > u_i(t_i,s_{-i}).
\]
\end{defbox}

Intuitivamente: in un potential game, ogni miglioramento unilaterale di un giocatore si traduce in un aumento del potenziale globale.

\begin{example}[Cournot -- potenziale ordinale]
Il duopolio di Cournot con qualsiasi funzione di domanda inversa $p$ ha il potenziale ordinale
\[
  P(x_1,x_2) = x_1\,x_2\,(p(x_1+x_2)-c).
\]
\end{example}

\begin{example}[BoS non ha potenziale esatto]
Per il Battle of Sexes si verifica che nessuna funzione $P$ soddisfa la condizione del potenziale esatto (il ciclo $FF\to FD\to DD\to DF\to FF$ porta a un'inconsistenza numerica).
\end{example}

\subsection{Caratterizzazione mediante cicli}

\begin{proposition}
Un gioco finito ha un potenziale ordinale se e solo se non ammette \emph{weak improvement cycles}: nessun ciclo $(x^0,x^1,\dots,x^\ell=x^0)$ dove ciascuno step è un miglioramento unilaterale (con almeno uno stretto).
\end{proposition}

\section*{In sintesi per l'orale}
\addcontentsline{toc}{section}{In sintesi per l'orale (L7)}
\begin{itemize}
  \item \textbf{BRD}: si parte da un profilo e si aggiorna con best response. Jacobi $=$ sincrono, tutti insieme; Gauss--Seidel $=$ asincrono, uno alla volta. Un punto fisso della dinamica è per definizione un NE.
  \item \textbf{Potenziale esatto}: $u_i(s_i,s_{-i})-u_i(t_i,s_{-i})=P(s_i,s_{-i})-P(t_i,s_{-i})$, la variazione di utilità \emph{coincide} con quella del potenziale. \textbf{Ordinale}: solo lo stesso \emph{segno}.
  \item Conseguenza: ogni miglioramento unilaterale aumenta $P$, quindi in un gioco finito non si può ciclare e ogni massimo di $P$ è un NE puro.
  \item Esempi: Cournot ha il potenziale ordinale $P(x_1,x_2)=x_1x_2(p(x_1+x_2)-c)$; BoS \emph{non} ha potenziale esatto (il ciclo $FF\to FD\to DD\to DF\to FF$ non torna a zero).
  \item \textbf{Caratterizzazione di Monderer--Shapley}: basta controllare i cicli di lunghezza 4 con due giocatori; se la somma alternata delle variazioni di utilità è zero su tutti, il potenziale esatto esiste. Per l'ordinale: nessun ciclo di miglioramento.
\end{itemize}

% ============================================================
\chapter{Convergenza e Potential Games Avanzati (L8)}
% ============================================================

\section{Comportamento dell'algoritmo sincrono}

\begin{example}[Dilemma del Prigioniero -- convergenza]
Da qualsiasi stato iniziale, l'algoritmo sincrono converge a $(D,D)$ in un passo (poiché $D$ è dominante per entrambi).
\end{example}

\begin{example}[Battle of Sexes -- ciclo]
Con $x^0=(F,D)$: entrambi trovano conveniente deviare, $x^1=(D,F)$; poi $x^2=(F,D)=x^0$: ciclo infinito.
\end{example}

\begin{example}[BoS con strategie miste -- ciclo]
Con $x^0=(1/2,1/2)$ per entrambi, l'algoritmo sincrono salta tra $(1,0)$ e $(0,1)$ senza convergere. Notare che l'NE misto $(2/3,1/3)$ non è un punto fisso dell'algoritmo sincrono (la best response a $(1/3,2/3)$ è la strategia pura $F$, non la mista).
\end{example}

Per evitare il looping si introduce la regola: se $x_i^k\in\BR_i(x_{-i}^k)$, mantieni $x_i^{k+1}=x_i^k$ (\emph{avoid useless switches}).

\section{Risultati di esistenza per potential games}

\begin{thmbox}{Esistenza per giochi continui}
Sia $G$ un gioco a potenziale esatto. Se $u_i$ è continua e $S_i\subseteq\mathbb{R}^{m_i}$ è compatto per ogni $i\in N$, allora $G$ ha almeno un NE.
\end{thmbox}
\textit{Dimostrazione.} Il potenziale $P$ è continuo su $S$ compatto, quindi raggiunge il massimo $s^*$. Per la definizione di potenziale esatto, $s^*$ è un NE. \qed

\begin{thmbox}{Esistenza per giochi finiti}
Ogni gioco finito con potenziale ordinale ha almeno un NE (in strategie pure).
\end{thmbox}

\begin{remark}
Non ogni NE massimizza il potenziale. Nel Battle of Sexes asimmetrico (payoff $(3,2)$ vs $(1,2)$), il profilo $(2,2)$ è NE ma non massimizza il potenziale esatto.
\end{remark}

\section{Convergenza dell'algoritmo asincrono}

\begin{thmbox}{Convergenza algoritmo asincrono}
Se $G$ è un gioco finito con potenziale ordinale, allora l'algoritmo distribuito asincrono (Gauss--Seidel con best responses) converge a un NE in un numero finito di passi.
\end{thmbox}
\textit{Dimostrazione.} Il potenziale $P$ è intero e cresce strettamente ad ogni iterazione che cambia la strategia (poiché ogni mossa è un miglioramento unilaterale). Essendoci un numero finito di profili, il processo deve terminare a un NE. \qed

\section*{In sintesi per l'orale}
\addcontentsline{toc}{section}{In sintesi per l'orale (L8)}
\begin{itemize}
  \item L'algoritmo \emph{sincrono} può ciclare: in BoS da $(F,D)$ si passa a $(D,F)$ e si torna indietro all'infinito. Rimedio: se sei già in best response non cambiare (``avoid useless switches'').
  \item Anche con le miste il sincrono può non convergere: in BoS l'NE misto $(2/3,1/3)$ non è un punto fisso della dinamica sincrona, perché la best response a una mista è una pura.
  \item \textbf{Esistenza nei giochi continui}: se $G$ ha potenziale esatto, $u_i$ continue e $S_i$ compatti, allora $P$ ha massimo e quel massimo è un NE. Dimostrazione in due righe, senza punto fisso: è il modo più economico di dimostrare esistenza.
  \item \textbf{Esistenza nei giochi finiti}: ogni gioco finito con potenziale ordinale ha un NE in strategie \emph{pure}.
  \item \textbf{Convergenza dell'asincrono}: in un gioco finito con potenziale ordinale, Gauss--Seidel con best response converge a un NE in un numero finito di passi (il potenziale cresce strettamente e i profili sono finiti).
  \item Attenzione: non ogni NE massimizza il potenziale; il massimo del potenziale è un NE, non il viceversa.
\end{itemize}

% ============================================================
\chapter{Apprendimento nei Giochi Finiti: Fictitious Play e Regret Matching (L9)}
% ============================================================

\section{Il problema dell'apprendimento}

La domanda della lezione è una sola: \textbf{gli equilibri si possono imparare
giocando?} Non si assume che i giocatori calcolino un equilibrio: si assume solo
che ripetano il gioco e reagiscano alla storia.

\begin{defbox}{Framework generale}
Si gioca ripetutamente lo \emph{stesso} gioco finito in forma normale. I giocatori
conoscono tutte le azioni passate di tutti, e aggiornano il proprio comportamento
in base alla storia osservata. La domanda è se le \emph{frequenze empiriche}
convergano a un equilibrio (di Nash o correlato).
\end{defbox}

\section{Fictitious play}

\begin{defbox}{Processo di fictitious play (Brown, 1949, 1951)}
\begin{enumerate}[leftmargin=*]
  \item Si parte da $\sigma^1=(\sigma^1_1,\dots,\sigma^1_n)$, $k=1$.
  \item Si calcola il \textbf{comportamento medio} degli avversari
        $\hat\sigma^k_{-i}=\frac{1}{k}\sum_{\ell=1}^{k}\sigma^\ell_{-i}$.
  \item Si gioca una best response alla media:
        $\sigma^{k+1}_i\in R_i(\hat\sigma^k_{-i})$.
  \item $k\leftarrow k+1$ e si torna al passo 2.
\end{enumerate}
\end{defbox}

\begin{remark}
Due osservazioni che il prof sottolinea: (i) \emph{non} serve conoscere le funzioni
di utilità degli avversari, bastano le loro azioni; (ii) \emph{non} esiste un
criterio di arresto, se non ``$\sigma^k$ è un equilibrio di Nash''.
\end{remark}

\begin{example}[Fictitious play in rock-paper-scissors]
Ordine delle strategie: (carta, forbici, sasso), $h(p,q)=p^{\mathsf T}Aq$ con
$A$ antisimmetrica. Partendo da carta contro forbici:
\begin{center}
\begin{tabular}{c|cc|cc}
$k$ & \multicolumn{2}{c|}{strategie giocate $\sigma^k$} & \multicolumn{2}{c}{comportamenti medi $\hat\sigma^k$}\\\hline
1 & $(1,0,0)$ & $(0,1,0)$ & $(1,0,0)$ & $(0,1,0)$\\
2 & $(0,0,1)$ & $(0,1,0)$ & $(1/2,0,1/2)$ & $(0,1,0)$\\
3 & $(0,0,1)$ & $(1,0,0)$ & $(1/3,0,2/3)$ & $(1/3,2/3,0)$\\
4 & $(0,1,0)$ & $(1,0,0)$ & $(1/4,1/4,1/2)$ & $(1/2,1/2,0)$\\
\end{tabular}
\end{center}
Le medie girano attorno a $((1/3,1/3,1/3),(1/3,1/3,1/3))$ e, essendo RPS a somma
zero, vi convergono (Robinson).
\end{example}

\begin{example}[``Gentle'' RPS di Shapley (1964): il fictitious play non converge]
\begin{center}
\begin{tabular}{c|ccc}
I/II & Left & Middle & Right\\\hline
Top & $(0,0)$ & $(1,0)$ & $(0,1)$\\
Middle & $(0,1)$ & $(0,0)$ & $(1,0)$\\
Down & $(1,0)$ & $(0,1)$ & $(0,0)$
\end{tabular}
\end{center}
Chi perde non perde davvero (payoff $0$ invece di $-1$): il gioco \emph{non} è a
somma zero. La traccia
\[
  (1,0,0)(0,1,0)\ \to\ (1,0,0)(0,0,1)\ \to\ (1,0,0)(0,0,1)\ \to\ (0,1,0)(0,0,1)\ \to\cdots
\]
dà medie $(1,0,0),(0,1/2,1/2)$; $(1,0,0),(0,1/3,2/3)$; $(3/4,1/4,0),(0,1/4,3/4)$: le
frequenze \emph{non} convergono all'unico NE $((1/3,1/3,1/3),(1/3,1/3,1/3))$, ma
finiscono in un \textbf{ciclo limite asintoticamente stabile}. È il controesempio
da citare quando si dice ``non converge in generale''.
\end{example}

\begin{thmbox}{Quando il fictitious play converge}
Le frequenze empiriche convergono a un NE per:
\begin{itemize}
  \item giochi a due giocatori a somma zero \hfill (Robinson 1951);
  \item giochi $2\times 2$ con la \emph{diagonal property}
        $a_{11}+a_{22}\ne a_{12}+a_{21}$ e $b_{11}+b_{22}\ne b_{12}+b_{21}$
        \hfill (Miyasawa 1961);
  \item giochi $2\times n$ non degeneri (best response unica alle pure)
        \hfill (Berger 2005);
  \item giochi risolvibili con IESDS \hfill (Milgrom--Roberts 1991);
  \item giochi a potenziale ordinale \hfill (Monderer--Shapley 1996).
\end{itemize}
Tasso nel caso somma zero: $0\le w_2(\sigma^k_2)-w_1(\sigma^k_1)=O(1/\sqrt[m]{k})$
con $m=m_1+m_2-2$, e $w_1(\sigma^k_1)\le V\le w_2(\sigma^k_2)$: i due livelli di
sicurezza stringono il valore del gioco.
\end{thmbox}

\begin{remark}
La convergenza è all'\emph{insieme} degli NE, e riguarda le \emph{frequenze
empiriche}, non le azioni giocate: in RPS nessun giocatore gioca mai la mista
uniforme, eppure la media delle sue mosse ci tende.
\end{remark}

\section{Regret matching}

Qui si cambia framework: si giocano strategie \emph{pure} (niente ambiente
fittizio) e si guarda al rimpianto.

\begin{defbox}{Processo di regret matching (Hart--Mas Colell, 2000)}
Sia $T_i^k(x_i^k)=\{t\le k:\ x_i^t=x_i^k\}$ l'insieme dei round passati in cui $i$
ha giocato la \emph{stessa} strategia che sta giocando ora. Il \textbf{regret} di
aver scelto $x_i^k$ invece di $s$ è
\[
  R_i^k(s,x_i^k)=\Bigl[\ \sum_{t\in T_i^k(x_i^k)}
      \bigl(u_i(s,x^t_{-i})-u_i(x_i^k,x^t_{-i})\bigr)\Bigr]^{+},
  \qquad [a]^+=\max\{a,0\}.
\]
Al round successivo si estrae una pura da
\[
  p_i^k(s)=\frac{R_i^k(s,x_i^k)}{\mu_i^k}\ \ (s\ne x_i^k),
  \qquad
  p_i^k(x_i^k)=1-\sum_{s\ne x_i^k}p_i^k(s),
\]
con $\mu_i^k\ge k\,(m_i-1)\max|u_i(x_i,x_{-i})-u_i(x_i',x_{-i})|$ abbastanza grande
da rendere $p_i^k$ una distribuzione.
\end{defbox}

\begin{remark}[Tre dettagli che è facile sbagliare]
\begin{enumerate}[leftmargin=*]
  \item La parte positiva $[\cdot]^+$ si prende \emph{dopo} la somma, non termine a
        termine: i round in cui l'alternativa sarebbe andata peggio compensano.
  \item Si somma solo sui round in cui si è giocata l'azione \emph{corrente}
        ($t\in T_i^k$): è un regret \emph{condizionato}, ed è il motivo per cui si
        arriva agli equilibri correlati e non solo ai loro parenti grossolani.
  \item La probabilità di restare sull'azione corrente è il \emph{residuo}
        $1-\sum_s p_i^k(s)$: più regret accumulato, più è probabile cambiare.
\end{enumerate}
\end{remark}

\begin{example}[Il calcolo del regret fatto a lezione]
Gioco $3\times 4$ (payoff $(u_{\mathrm I},u_{\mathrm{II}})$):
\begin{center}
\begin{tabular}{c|cccc}
I/II & 1 & 2 & 3 & 4\\\hline
1 & $(5,3)$ & $(7,4)$ & $(5,2)$ & $(3,4)$\\
2 & $(5,5)$ & $(5,7)$ & $(1,1)$ & $(2,5)$\\
3 & $(3,4)$ & $(4,2)$ & $(5,5)$ & $(6,3)$
\end{tabular}
\end{center}
Storia dei primi quattro round e regret del giocatore I, che ora sta giocando la
strategia $3$:
\begin{center}
\begin{tabular}{c|cc|cc|cc}
$k$ & profilo & $u_{\mathrm I}$ & se avessi giocato $1$ & se avessi giocato $2$ & regret vs $1$ & regret vs $2$\\\hline
1 & $(3,3)$ & 5 & 5 & 1 & $0$ & $-4$\\
2 & $(2,1)$ & 5 & 5 & 5 & \multicolumn{2}{c}{\emph{escluso}: I giocò $2$}\\
3 & $(3,1)$ & 3 & 5 & 5 & $2$ & $2$\\
4 & $(3,2)$ & 4 & 7 & 5 & $3$ & $1$\\\hline
\multicolumn{5}{r|}{totale su $T=\{1,3,4\}$} & $\mathbf{5}$ & $-1\to \mathbf{0}$
\end{tabular}
\end{center}
Quindi $p(1)=5/\mu^k$, $p(2)=0$, $p(3)=1-5/\mu^k$. Si noti il round $2$ escluso
(I aveva giocato un'altra azione) e il $-1$ azzerato dalla parte positiva
\emph{finale}: se si prendesse $[\cdot]^+$ termine a termine verrebbe $3$, sbagliato.
\end{example}

\begin{defbox}{Distribuzione empirica}
Dopo $k$ round, $z_k(x)=|\{t\le k:\ x^t=x\}|/k$ è la frequenza con cui è stato
giocato il profilo congiunto $x\in S$. Attenzione: è una distribuzione su $S$,
\emph{non} un prodotto di distribuzioni individuali. È proprio questa correlazione
che permette di superare i NE.
\end{defbox}

\begin{thmbox}{Convergenza del regret matching}
La successione delle distribuzioni empiriche $z_k$ converge quasi certamente,
per $k\to+\infty$, \textbf{all'insieme degli equilibri correlati}. Il regret medio
tende a zero con tasso $O(1/\sqrt{k})$.
\end{thmbox}

\section{Equilibrio correlato}

\begin{defbox}{Equilibrio correlato (Aumann, 1974)}
Una distribuzione $\mu$ su $S$ è un \textbf{equilibrio correlato} se per ogni
giocatore $i$, ogni $s_i$ con probabilità marginale positiva e ogni deviazione
$s_i'\in S_i$:
\[
  \sum_{s_{-i}\in S_{-i}}\mu(s_i,s_{-i})\,
  \bigl[u_i(s_i,s_{-i})-u_i(s_i',s_{-i})\bigr]\ \ge\ 0 .
\]
Interpretazione: un mediatore estrae $s\sim\mu$ e comunica \emph{in privato} a
ciascuno solo la propria componente $s_i$; il vincolo dice che, condizionatamente
al segnale ricevuto, obbedire è ottimo.
\end{defbox}

\begin{remark}[Perché piace: è un poliedro]
I vincoli sono \emph{lineari} in $\mu$ (più $\mu\ge 0$ e $\sum_s\mu(s)=1$): l'insieme
degli equilibri correlati è un politopo, quindi si calcola con la programmazione
lineare e si può ottimizzare su di esso (per esempio massimizzando il benessere
sociale). Per contro, calcolare un NE è un problema di complementarità, non un LP.
Inoltre l'esistenza segue dalla dualità LP, senza teoremi di punto fisso.
\end{remark}

\begin{proposition}
Ogni NE (misto) è un equilibrio correlato: basta prendere $\mu=\prod_i\sigma_i$.
L'insieme degli equilibri correlati è convesso e contiene l'inviluppo convesso dei
NE; in generale lo contiene \emph{strettamente}.
\end{proposition}

\begin{example}[Il semaforo: un CE migliore di ogni NE]
Gioco del pollo (Hawk-Dove riscalato), $C=$ cedere, $D=$ tirare dritto:
\begin{center}
\begin{tabular}{c|cc}
I/II & $C$ & $D$\\\hline
$C$ & $(6,6)$ & $(2,7)$\\
$D$ & $(7,2)$ & $(0,0)$
\end{tabular}
\end{center}
NE puri $(C,D)$ e $(D,C)$ con payoff $(2,7)$ e $(7,2)$; NE misto: l'avversario
cede con probabilità $q$ tale che $6q+2(1-q)=7q$, cioè $q=2/3$, con payoff
$14/3\approx 4{,}67$ a testa.

Sia ora $\mu$ uniforme su $\{(C,C),(C,D),(D,C)\}$, cioè $1/3$ ciascuno: è il
semaforo che dice a uno di passare e all'altro di fermarsi, e ogni tanto lascia
passare entrambi con prudenza. Verifica dei vincoli per il giocatore I:
\begin{itemize}
  \item segnale $D$ (probabilità $1/3$): l'altro sta giocando $C$ con certezza,
        obbedire dà $7$, deviare a $C$ dà $6$. \checkmark
  \item segnale $C$ (probabilità $2/3$): l'altro gioca $C$ o $D$ con probabilità
        $1/2$ ciascuna, obbedire dà $(6+2)/2=4$, deviare a $D$ dà $(7+0)/2=3{,}5$.
        \checkmark
\end{itemize}
Payoff atteso: $(6+2+7)/3=5$ a testa, quindi $10$ di benessere complessivo, contro
$9{,}33$ del NE misto e $9$ di qualunque combinazione pubblica dei due NE puri.
La correlazione con \emph{informazione privata} fa strettamente meglio.
\end{example}

\begin{remark}[Correlato e correlato grossolano]
Se i giocatori minimizzano il regret \emph{esterno} (confronto con una singola
azione fissa, senza condizionare sull'azione giocata) le frequenze convergono agli
equilibri \emph{correlati grossolani}, un insieme più grande. Il regret
condizionato del regret matching, che somma solo sui round con la stessa azione,
è ciò che porta agli equilibri correlati veri e propri.
\end{remark}

\section{Giochi sequenziali e minimax}\label{sec:minimax-intro}

Un gioco sequenziale a \textbf{somma zero} tra un massimizzatore (MAX) e un
minimizzatore (MIN) si risolve con la procedura \textbf{minimax} (backward induction):
\begin{itemize}
  \item Ai nodi di MAX: si sceglie il figlio con valore massimo.
  \item Ai nodi di MIN: si sceglie il figlio con valore minimo.
  \item Il valore si propaga verso la radice.
\end{itemize}
Per alberi di grandi dimensioni (es.\ scacchi), la backward induction completa è
computazionalmente impraticabile e si ricorre all'\textbf{alfa-beta pruning}
(descritto nel Capitolo~\ref{ch:alpha-beta}) e alle \textbf{funzioni di valutazione}
che stimano il valore di posizioni non terminali.

\section*{In sintesi per l'orale}
\addcontentsline{toc}{section}{In sintesi per l'orale (L9)}
\begin{itemize}
  \item Domanda della lezione: gli equilibri si imparano? Due risposte, due
        algoritmi, due insiemi limite diversi.
  \item \textbf{Fictitious play}: best response al comportamento \emph{medio} degli
        avversari. Non serve conoscere le utilità altrui, non c'è criterio di
        arresto. Converge in: somma zero, $2\times2$ con diagonal property,
        $2\times n$ non degeneri, IESDS-risolvibili, potenziale ordinale.
  \item Controesempio obbligatorio: gentle RPS di Shapley (1964), ciclo limite
        asintoticamente stabile, non converge.
  \item \textbf{Regret matching}: si giocano pure; regret condizionato
        \[
          R^k(s,x^k)=\Bigl[\sum_{t\,:\,x^t=x^k}\bigl(u(s,x^t_{-i})-u(x^k,x^t_{-i})\bigr)\Bigr]^{+},
        \]
        probabilità proporzionale ai regret positivi, il resto sull'azione
        corrente. La parte positiva si prende \emph{dopo} la somma.
  \item Le distribuzioni empiriche $z_k$ convergono quasi certamente all'insieme
        degli \textbf{equilibri correlati}; regret medio $O(1/\sqrt{k})$.
  \item \textbf{Equilibrio correlato}: $\sum_{s_{-i}}\mu(s_i,s_{-i})[u_i(s_i,s_{-i})
        -u_i(s_i',s_{-i})]\ge 0$. Vincoli lineari $\Rightarrow$ politopo,
        calcolabile con LP, esistenza senza punto fisso; contiene l'inviluppo
        convesso dei NE, in generale in senso stretto.
  \item Esempio pronto: il semaforo nel gioco del pollo, $1/3$ su
        $(C,C),(C,D),(D,C)$, payoff $5$ a testa contro $4{,}67$ del NE misto.
  \item Collegamento con L6: nei giochi a somma zero il valore si calcola con un LP
        (dualità forte); qui è l'insieme degli equilibri correlati a essere
        descritto da un LP. In entrambi i casi la linearità è ciò che rende il
        concetto algoritmico.
\end{itemize}

% ============================================================
\chapter{Giochi di Stackelberg e Backward Induction (L10)}
\label{ch:stackelberg}
% ============================================================

\section{Giochi di Stackelberg}

In un \textbf{gioco di Stackelberg} (von Stackelberg, 1934) un giocatore (\emph{leader}) si muove per primo, e gli altri (\emph{follower}) osservano la sua mossa prima di rispondere. Il leader può quindi \emph{anticipare} le risposte razionali dei follower.

\begin{defbox}{Subgame Perfect Equilibrium (SPE)}
In un gioco esteso, un profilo di strategie $\sigma^*$ è un \textbf{subgame perfect equilibrium} (equilibrio perfetto nei sottogiochi) se costituisce un NE in ogni sottogioco del gioco originale.
\end{defbox}

L'SPE è il concetto di soluzione naturale per i giochi sequenziali: rimuove le minacce non credibili (minacce che, se attuate, danneggerebbero anche chi le fa).

\subsection{Backward induction e SPE}

Nei giochi a informazione perfetta, l'SPE si calcola tramite backward induction:
\begin{enumerate}
  \item Ai nodi terminali, il payoff è noto.
  \item Risalendo, ogni giocatore sceglie la mossa che massimizza il proprio payoff (dato ciò che accade in seguito).
\end{enumerate}
Con best responses non uniche, si distinguono:
\begin{itemize}
  \item \textbf{Atteggiamento ottimistico}: il leader assume che i tie vengano risolti a suo favore.
  \item \textbf{Atteggiamento pessimistico}: il leader assume il peggior caso.
\end{itemize}

\section{Duopolio di Stackelberg con bene indivisibile}

Si riprende il duopolio con $x_i\in\mathbb{Z}_+$, $T=10$, $c=3$:
\[
  u_i(x_1,x_2) = x_i\max\{T-(x_1+x_2),0\} - c\,x_i.
\]
Il follower (giocatore 2) osserva $x_1$ e risponde con
\[
  x_2^*(x_1) = \begin{cases}
    \displaystyle\left\lfloor\frac{T-x_1-c}{2}\right\rfloor & \text{se }T-x_1>c \\
    0 & \text{altrimenti}
  \end{cases}
\]
Il leader anticipa questa risposta e massimizza $u_1(x_1, x_2^*(x_1))$. Con atteggiamento ottimistico, la soluzione di Stackelberg è $x_1^*=4$, $x_2^*=1$, con payoff $(8,2)$ (migliore per il leader rispetto al NE di Cournot $(2,2)$ con payoff $(6,6)$).

\section{Gioco a due stadi con investitori}

Due investitori devono decidere se rimanere (\emph{stay in}) o ritirarsi (\emph{withdraw}) da un progetto bancario in due stadi. Ogni investitore ha investito 4M€. Il secondo stadio ha come unico NE \emph{(stay in, stay in)} con payoff $(10,10)$. Sostituendo nell'albero del primo stadio e applicando la backward induction, si ottengono due NE al primo stadio (a seconda del comportamento in caso di informazione imperfetta).

\section*{In sintesi per l'orale}
\addcontentsline{toc}{section}{In sintesi per l'orale (L10)}
\begin{itemize}
  \item \textbf{Stackelberg}: un leader muove per primo, il follower osserva e risponde. Il leader risolve $\max_{x_1} u_1(x_1, r(x_1))$ dove $r$ è la best response del follower: è ottimizzazione a due livelli, non un gioco simultaneo.
  \item Se la best response del follower non è unica servono due letture: \emph{ottimistica} (il follower sceglie, tra le sue ottime, quella che conviene al leader) e \emph{pessimistica} (quella che gli conviene meno). Vanno nominate entrambe.
  \item La soluzione di Stackelberg è un SPE del gioco sequenziale, e in generale dà al leader almeno quanto il NE simultaneo: è il \emph{vantaggio della prima mossa}.
  \item Esempi: duopolio di Stackelberg con bene indivisibile (il leader produce di più, il follower di meno rispetto a Cournot) e il gioco a due stadi con investitori.
  \item Collegamento: backward induction $\to$ SPE (L11), e la stessa struttura a due livelli torna nei giochi Stackelberg--Nash con più follower (L13).
\end{itemize}

% ============================================================
\chapter{Giochi Sequenziali Finiti: SPE e Backward Induction (L11)}
\label{ch:seq-finite}
% ============================================================

\section{L'albero di enumerazione}

Un \textbf{gioco sequenziale finito con informazione perfetta} si modella formalmente come un \textbf{albero di enumerazione} (\emph{enumeration tree}).

\begin{defbox}{Albero di Enumerazione}
Un \textbf{albero di enumerazione} è un albero radicato orientato verso le foglie (\emph{directed rooted out-tree}): tutti gli archi puntano dalla radice verso le foglie. La struttura è:
\begin{itemize}
  \item \textbf{Nodo} (\emph{node}): rappresenta uno \emph{stato} del gioco.
  \item \textbf{Foglia} (\emph{leaf}): nodo terminale, rappresenta un \emph{outcome} (risultato finale). Ogni foglia è etichettata con il vettore dei payoff $(u_1,\dots,u_n)$.
  \item \textbf{Turn / Ply}: l'insieme di nodi allo stesso livello di profondità (stessa distanza dalla radice). Ogni turn/ply corrisponde a un ``turno'' del gioco.
  \item \textbf{Arco} (\emph{arc}): rappresenta un'azione disponibile per il giocatore proprietario del nodo coda dell'arco.
  \item Ogni nodo non-foglia appartiene a \emph{esattamente un} giocatore.
  \item Le etichette sulle foglie sono le utilità/payoff dei giocatori.
\end{itemize}
\end{defbox}

\begin{remark}
La \emph{definizione formale matematica} di albero di enumerazione come tupla $(V, E, r, P, A, \ell)$ è piuttosto pesante. Le slides del corso osservano esplicitamente che ``the formal mathematical definition is not very handy''; è più utile ragionare direttamente sulla struttura ad albero.
\end{remark}

\begin{remark}[\textbf{Informazione perfetta}]
``Informazione perfetta'' significa che ogni giocatore conosce \emph{tutte} le mosse precedenti: non c'è alcuna incertezza sulla storia del gioco. In termini dell'albero, ogni giocatore sa esattamente in quale nodo si trova.
\end{remark}

\begin{remark}[\textbf{Crescita esponenziale}]
L'albero di enumerazione cresce \emph{esponenzialmente} con il numero di mosse. Per giochi complessi (scacchi: $\approx 10^{43}$ posizioni, go: ancora di più), è computazionalmente impraticabile costruire l'albero per intero. Questa osservazione motiva tecniche di potatura come l'alfa-beta pruning (Capitolo~\ref{ch:alpha-beta}) e l'uso di funzioni di valutazione (Deep Blue).
\end{remark}

\section{Strategie nei giochi sequenziali}

\begin{defbox}{Strategia pura (gioco sequenziale)}
Una \textbf{strategia pura} per il giocatore $i$ è una funzione
\[
  \sigma_i : D_i \;\longrightarrow\; \bigcup_{v \in D_i} A(v),
  \qquad \sigma_i(v)\in A(v)\;\forall v\in D_i,
\]
dove $D_i$ è l'insieme dei nodi di decisione di $i$ e $A(v)$ è l'insieme delle azioni disponibili al nodo $v$.

Un profilo di strategie pure $\sigma = (\sigma_1,\dots,\sigma_n)$ determina un \emph{unico cammino} dalla radice a una foglia, e quindi un unico outcome del gioco.
\end{defbox}

\begin{remark}
Il numero di strategie pure per il giocatore $i$ è
\[
  |S_i| = \prod_{v \in D_i} |A(v)|.
\]
Questo numero è in generale esponenziale nel numero di nodi di decisione. Per esempio, se il giocatore~I ha 5 nodi ognuno con 2 azioni disponibili, $|S_I| = 2^5 = 32$; se il giocatore~II ha 10 nodi, $|S_{II}| = 2^{10} = 1024$.
\end{remark}

\section{Sottogiochi e Subgame Perfect Equilibrium}

\begin{defbox}{Sottogioco}
Un \textbf{sottogioco} (\emph{subgame}) $G_v$ di un gioco sequenziale $G$ è il sottoalbero di enumerazione radicato in un nodo non-foglia $v$: esso è a sua volta un gioco sequenziale finito completo (con la propria radice, le proprie foglie e gli stessi giocatori).
\end{defbox}

\begin{defbox}{Subgame Perfect Equilibrium -- SPE (Selten, 1965)}
Un profilo di strategie $\sigma^*$ è un \textbf{Subgame Perfect Equilibrium} (SPE) se $\sigma^*$ costituisce un Nash Equilibrium \emph{in ogni sottogioco} di $G$ (incluso $G$ stesso).
\end{defbox}

\begin{remark}
L'SPE raffina il concetto di NE eliminando le \emph{minacce non credibili}: un NE può prescrivere azioni che, se effettivamente eseguite, danneggerebbero il giocatore stesso. L'SPE richiede razionalità in ogni punto dell'albero, non solo lungo il percorso di equilibrio.
\end{remark}

\begin{thmbox}{Esistenza dell'SPE (Selten, 1965)}
\label{thm:spe-existence}
Ogni gioco sequenziale finito con informazione perfetta ammette almeno un Subgame Perfect Equilibrium in strategie pure.
\end{thmbox}

\begin{proof}
Poiché il gioco è finito, l'albero ha un numero finito di nodi. La procedura di backward induction (definizione seguente) costruisce esplicitamente un SPE e termina in tempo finito. L'esistenza è quindi garantita per costruzione.
\end{proof}

\section{Procedura di Backward Induction}

\begin{defbox}{Backward Induction}
La \textbf{procedura di backward induction} calcola un SPE in un gioco sequenziale finito con informazione perfetta, procedendo dai nodi foglia verso la radice:
\begin{enumerate}
  \item \textbf{Inizializzazione}: per ogni foglia $\ell$, poni $\mathit{val}(\ell)$ uguale al vettore di payoff di $\ell$.
  \item \textbf{Passo induttivo}: per ogni nodo non-foglia $v$ proprietà del giocatore $i$, i cui tutti i figli hanno già $\mathit{val}$ assegnato:
    \begin{itemize}
      \item scegli il figlio $c^*$ che massimizza $(u_i)$-esima componente di $\mathit{val}(c)$;
      \item poni $\mathit{val}(v) := \mathit{val}(c^*)$ e assegna a $v$ l'azione corrispondente all'arco $(v,c^*)$.
    \end{itemize}
  \item \textbf{Terminazione}: ripeti il passo~2 finché non si raggiunge la radice. Il valore $\mathit{val}(\text{root})$ è il payoff di equilibrio e le azioni assegnate a tutti i nodi formano un SPE.
\end{enumerate}
\end{defbox}

\begin{remark}
La backward induction effettua un'esplorazione completa dell'albero (al peggio $O(|V|)$ operazioni). Per alberi di grandi dimensioni, l'alfa-beta pruning ne riduce drasticamente il costo computazionale.
\end{remark}

\section{Esempio: Pure Coordination Game}

\begin{example}[Pure Coordination Game]
Due giocatori devono coordinare la propria scelta tra due azioni: \emph{leave} (L) e \emph{cooperate} (C). In forma simultanea:
\begin{center}
\begin{tabular}{c|cc}
I/II & L & C \\\hline
L & $(1,1)$ & $(1,0)$ \\
C & $(0,1)$ & $(3,3)$
\end{tabular}
\end{center}
Vi sono due NE puri: $(L,L)$ e $(C,C)$.

In forma sequenziale (il giocatore~I muove per primo), la backward induction produce:
\begin{itemize}
  \item Se I sceglie L: II risponde con L (payoff $1>0$ per lui) → outcome $(1,1)$.
  \item Se I sceglie C: II risponde con C (payoff $3>1$ per lui) → outcome $(3,3)$.
\end{itemize}
Il giocatore~I ottiene $1$ scegliendo L, $3$ scegliendo C → I sceglie \textbf{C}.

\textbf{Unico SPE}: $(C, C)$ con payoff $(3,3)$. Il NE $(L,L)$ non è SPE perché se I sceglie C, II risponde con C (non con L): la ``minaccia'' di II di giocare L è non credibile.
\end{example}

\section{Esempio: Sequential Battle of Sexes}

\begin{example}[Sequential Battle of Sexes]
Si riprende il Battle of Sexes con lei (giocatrice~1) che si muove per prima, scegliendo Football (F) o Ballet (B); lui (giocatore~2) osserva la scelta e risponde.

\textbf{Payoff}: (F,F)=$(2,1)$; (F,B)=$(0,0)$; (B,F)=$(0,0)$; (B,B)=$(1,2)$.

Il giocatore~2 ha due nodi di decisione (uno dopo F, uno dopo B), ciascuno con 2 azioni: pertanto ha $2^2=4$ strategie pure. Denotiamo $\sigma_2(F)/\sigma_2(B)$ le sue scelte condizionate.

\textbf{Backward induction}:
\begin{itemize}
  \item Se lei gioca F: lui sceglie F (payoff $1>0$ per lui) → outcome $(2,1)$.
  \item Se lei gioca B: lui sceglie B (payoff $2>0$ per lui) → outcome $(1,2)$.
\end{itemize}
Lei ottiene $2$ scegliendo F, $1$ scegliendo B → lei sceglie \textbf{F}.

\textbf{Unico SPE}: lei gioca F; lui gioca ``segui lei'' $\sigma_2=(F\mid F,\; B\mid B)$. Outcome di equilibrio: $(F,F)$ con payoff $(2,1)$.

\textbf{Confronto con il gioco simultaneo}: nel BoS simultaneo ci sono 2 NE puri e 1 NE misto. Nella versione sequenziale, chi muove per primo ha un vantaggio: può \emph{impegnarsi} credibilmente e il secondo giocatore si adatta. Il gioco sequenziale produce un unico SPE che favorisce chi muove prima.
\end{example}

\section{Veto driven choice}

\begin{example}[Veto Driven Choice]
Due giocatori devono accordarsi su un'alternativa tra $\{a,b,c\}$. Le preferenze sono: $a\succ b\succ c$ per il giocatore 1; $c\succ b\succ a$ per il giocatore 2. Ad ogni turno ciascun giocatore \emph{veta} un'alternativa; l'alternativa finale è quella non vetata.

Struttura: l'albero ha $3\times 3 = 9$ foglie; si costruisce la matrice $8\times 8$ delle strategie (ognuno ha 3 possibili prime veto e poi 2 seconde veto). L'equilibrio per backward induction porta il giocatore 1 a scegliere $c$ come primo veto (eliminare la sua alternativa peggiore) e il giocatore 2 a rispondere in modo da ottenere $b$: payoff $(0,2)$ -- un esito di compromesso.
\end{example}

\section{Centipede game}

\begin{example}[Centipede Game (Rosenthal, 1981)]
Due giocatori si alternano a decidere se \emph{continuare} (C) o \emph{fermarsi} (S). Se si ferma al turno $t$, il giocatore corrente ottiene il proprio portafoglio; se continua, trasferisce $1$€ all'avversario che guadagna anche $1$€ extra. I payoff crescono esponenzialmente:
\[
  S: (1,1)\to (0,3)\to (2,2)\to (1,4)\to (3,3)\to\cdots
\]
La \textbf{backward induction} porta all'SPE: fermarsi immediatamente, con payoff $(1,1)$.

Il \textbf{paradosso di inefficienza}: se entrambi cooperassero fino all'ultimo turno, otterrebbero payoff arbitrariamente grandi. L'SPE è Pareto-dominato dalla cooperazione.
\end{example}

\section{Chain store paradox}

\begin{example}[Chain Store Paradox (Selten, 1978)]
Un chain store ha filiali in $n$ città. In ogni città un potenziale concorrente deve scegliere se \emph{entrare} nel mercato o \emph{restare fuori}; il chain store risponde cooperando o agendo aggressivamente. Le città vengono visitate in ordine, con informazione perfetta. I payoff per ogni stadio sono:
\begin{itemize}
  \item (stai fuori): $(1+5,\cdot)$ per il chain store.
  \item (entra, coopera): $(2,2)$.
  \item (entra, aggressivo): $(0,0)$.
\end{itemize}
La backward induction: all'ultimo stadio il chain store preferisce cooperare (payoff $2>0$); quindi il concorrente entra; questo si propaga a ritroso -- il chain store coopera \emph{sempre}. \textbf{Paradosso}: il comportamento aggressivo iniziale potrebbe dissuadere i futuri concorrenti, ma non è razionale nella logica backward.
\end{example}

\section{Informazione imperfetta}

\begin{defbox}{Information Set}
In un gioco in forma estesa, un \textbf{information set} è un insieme di nodi che:
\begin{itemize}
  \item appartengono allo stesso giocatore;
  \item hanno lo stesso nodo genitore (o hanno avuto la stessa storia ``visibile'');
  \item hanno le stesse azioni disponibili.
\end{itemize}
Il giocatore non può distinguere i nodi appartenenti allo stesso information set.
\end{defbox}

Un gioco in forma normale equivalente al gioco esteso si costruisce enumerando le strategie pure come funzioni dagli information sets alle azioni. Nell'\emph{Hawk-Hare-Hunt} (stag hunt) rappresentato come gioco sequenziale con information set sul nodo del giocatore 2, la matrice equivale esattamente al gioco simultaneo.

\section{Forward induction}

\begin{description}
  \item[Backward induction:] le mosse future saranno razionali (previsione del futuro).
  \item[Forward induction:] le mosse passate sono state razionali (inferenza sul passato).
\end{description}

La forward induction consente di raffinare gli NE eliminando quelli inconsistenti con la razionalità osservata. Esempio: se il giocatore 1 in un gioco esteso ha scelto azione $s$ (anzichè $h$) raggiungendo l'information set $A$ del giocatore 2, quest'ultimo può inferire che il giocatore 1 ha probabilmente scelto $s$ razionalmente, aggiornando la propria credenza. Questo restringe le best response ammissibili.

\section*{In sintesi per l'orale}
\addcontentsline{toc}{section}{In sintesi per l'orale (L11)}
\begin{itemize}
  \item In un gioco sequenziale una \emph{strategia} è un piano completo: un'azione per ogni insieme di informazione, anche nei nodi che non si raggiungeranno mai. Da qui il numero di strategie pure come prodotto.
  \item \textbf{Sottogioco}: sottoalbero radicato in un nodo che è un singoletto del suo insieme di informazione e che contiene \emph{interi} tutti gli insiemi di informazione che tocca. Senza questa clausola la definizione è sbagliata.
  \item \textbf{SPE}: profilo che induce un NE in \emph{ogni} sottogioco. Ogni SPE è un NE, non viceversa: l'SPE elimina le \emph{minacce non credibili} (chain store paradox).
  \item \textbf{Backward induction}: su alberi finiti a informazione perfetta produce un SPE e, se i payoff sono tutti distinti, è unico (Kuhn, versione costruttiva di Zermelo).
  \item Esempi: pure coordination, BoS sequenziale (il primo a muovere impone il suo equilibrio preferito), veto driven choice, \textbf{centipede} (l'SPE dice di uscire al primo nodo, i soggetti reali no: limite del concetto), forward induction come lettura opposta.
\end{itemize}

% ============================================================
\chapter{Alpha-Beta Pruning, Giochi Multi-Stage e Oligopoli (L12)}
\label{ch:alpha-beta}
% ============================================================

\section{Gioco sequenziale minimax: un esempio completo}

Prima di presentare l'alpha-beta pruning in dettaglio, analizziamo un esempio di backward induction su un gioco a somma zero tra un \textbf{massimizzatore} (giocatore~1, MAX) e un \textbf{minimizzatore} (giocatore~2, MIN).

\begin{example}[Sequential Minimax Game]
Consideriamo un albero di gioco in cui MAX muove per primo (radice). L'albero ha profondità tale che:
\begin{itemize}
  \item Il giocatore~I (MAX) ha $2^5 = 32$ strategie pure (5 nodi di decisione, ciascuno con 2 azioni).
  \item Il giocatore~II (MIN) ha $2^{10} = 1024$ strategie pure (10 nodi di decisione, ciascuno con 2 azioni).
\end{itemize}
La backward induction procede foglia per foglia, livello per livello. Consideriamo un sottoalbero con le seguenti foglie (da sinistra a destra): $3,\,5,\,2,\,9,\,0,\,6,\,4,\,2$.

Applicando la backward induction:
\begin{enumerate}
  \item \textbf{Livello MIN (profondità 3)}: i nodi MIN scelgono il minimo tra i propri figli:
    \[\min(3,5)=3,\quad \min(2,9)=2,\quad \min(0,6)=0,\quad \min(4,2)=2.\]
  \item \textbf{Livello MAX (profondità 2)}: i nodi MAX scelgono il massimo:
    \[\max(3,2)=3,\quad \max(0,2)=2.\]
  \item \textbf{Livello MIN (profondità 1)}: il nodo MIN sceglie il minimo:
    \[\min(3,2)=2.\]
  \item \textbf{Livello MAX (radice)}: MAX sceglie il massimo; supponiamo ci sia un unico figlio: il valore è $\boxed{3}$.
\end{enumerate}

\textbf{Osservazione}: la backward induction effettua un'esplorazione \emph{completa} dell'albero (nessuna potatura). Per alberi profondi, il costo computazionale è $O(b^d)$ dove $b$ è il fattore di ramificazione e $d$ la profondità.
\end{example}

\section{Alfa-beta pruning: algoritmo e esempio dettagliato}

L'alfa-beta pruning è un'ottimizzazione della backward induction che elimina sottorami dell'albero senza comprometterne la correttezza.

\begin{defbox}{Alfa-beta Pruning}
L'algoritmo mantiene due parametri globali lungo il percorso radice-nodo corrente:
\begin{itemize}
  \item $\alpha$ (\emph{alpha}): il valore minimo garantito al massimizzatore lungo il percorso corrente (``il meglio che MAX può assicurarsi'');
  \item $\beta$ (\emph{beta}): il valore massimo garantito al minimizzatore lungo il percorso corrente (``il meglio che MIN può assicurarsi'').
\end{itemize}
Si usa la coppia $(\alpha,\beta)$ del percorso padre-figlio per rilevare sottorami non proficui:
\begin{itemize}
  \item \textbf{Potatura $\alpha$} (al nodo MIN): se il valore corrente $v \le \alpha$, il padre MAX non sceglierà mai questo nodo (poiché può già garantirsi $\alpha>\mathit{v}$): si potano tutti i fratelli non ancora esplorati del nodo corrente.
  \item \textbf{Potatura $\beta$} (al nodo MAX): se il valore corrente $v \ge \beta$, il padre MIN non sceglierà mai questo nodo: si potano i fratelli.
\end{itemize}
\end{defbox}

\begin{thmbox}{Complessità dell'alfa-beta pruning}
Nel caso migliore (ordinamento ottimale dei nodi), l'alfa-beta pruning riduce la complessità da $O(b^d)$ a $O(b^{d/2})$: si raddoppia la profondità esplorabile a parità di risorse computazionali.
\end{thmbox}

\begin{example}[Alfa-beta pruning: esempio con $\alpha=3, \beta=3$]
Consideriamo un albero di gioco a due livelli con MAX alla radice e MIN al primo livello. Esploriamo il primo sottoalbero sinistro del nodo MIN:

\textbf{Fase 1}: la prima foglia vale $6$. Poniamo $\beta_{\text{MIN}} = 6$ (il nodo MIN può garantirsi al più $6$).

\textbf{Fase 2}: la seconda foglia vale $3$. Il nodo MIN aggiorna $\beta_{\text{MIN}} = \min(6,3) = 3$. Il nodo MAX (radice) vede $v=3$ e aggiorna $\alpha_{\text{MAX}} = 3$.

\textbf{Fase 3}: la terza foglia vale $2$. Il nodo MIN aggiorna $\beta_{\text{MIN}} = \min(3,2) = 2 < \alpha_{\text{MAX}} = 3$.

$\Rightarrow$ \textbf{Potatura $\alpha$}: il nodo MAX alla radice non sceglierà mai questo nodo MIN (che garantisce al più $2 < 3$). Si pota il fratello successivo della foglia.

\textbf{Risultato}: la foglia di valore $5$ (fratello potato) \emph{non viene esplorata}. Il valore propagato alla radice è $\beta_{\text{MIN}} = 3$.

\begin{center}
\begin{tabular}{clll}
\toprule
Passo & Foglia esplorata & Aggiornamento & Azione \\
\midrule
1 & $6$ & $\beta_{\text{MIN}} = 6$ & Continua \\
2 & $3$ & $\beta_{\text{MIN}} = 3$, $\alpha_{\text{MAX}} = 3$ & Continua \\
3 & $2$ & $\beta_{\text{MIN}} = 2 \le \alpha_{\text{MAX}} = 3$ & \textbf{Pota} il ramo con $5$ \\
\bottomrule
\end{tabular}
\end{center}

\textbf{Interpretazione}: il nodo MIN, avendo già trovato un outcome di valore $2$, non può migliorare rispetto alla garanzia $\alpha=3$ di MAX. Esplorare il fratello (valore $5$) sarebbe inutile: MAX non sceglierebbe comunque questo sottoalbero.
\end{example}

\section{Deep guessing e funzioni di valutazione}

In giochi con alberi estremamente profondi — scacchi ($\approx 10^{40}$ stati), go ($\approx 10^{170}$) — esplorare l'albero fino alle foglie terminali è computazionalmente impossibile. Si usa la tecnica del \textbf{deep guessing}:

\begin{enumerate}
  \item Si fissa un \textbf{limite di profondità} $d$.
  \item Si esegue minimax (o alfa-beta pruning) fino alla profondità $d$.
  \item I nodi al livello $d$ (non terminali) vengono valutati da una \textbf{funzione di valutazione euristica} $f: S \to \mathbb{R}$ che approssima il valore ``vero'' del nodo.
\end{enumerate}

\begin{defbox}{Funzione di valutazione euristica}
Una funzione di valutazione $f: S \to \mathbb{R}$ assegna uno score a ogni stato $s\in S$, anche non terminale, approssimando il valore che otterrebbe il giocatore MAX se la partita continuasse in modo ottimale. Le proprietà desiderabili sono:
\begin{itemize}
  \item \textbf{Correttezza sui terminali}: $f(s)=u(s)$ per ogni foglia $s$.
  \item \textbf{Efficienza}: calcolabile in tempo $O(1)$ o $O(\text{poly})$ rispetto alla dimensione dello stato.
  \item \textbf{Informativit\`a}: cattura caratteristiche rilevanti del dominio (vantaggio materiale, struttura posizionale, ecc.).
\end{itemize}
\end{defbox}

\begin{example}[Deep Blue -- scacchi]
Deep Blue (IBM, 1997) usava una funzione di valutazione lineare
\[
  f(s) = \sum_{k=1}^{8000} w_k \cdot \phi_k(s)
\]
dove $\phi_k(s)$ sono $8000$ feature dello stato scacchistico (vantaggio materiale, struttura dei pedoni, sicurezza del re, mobilità dei pezzi, \ldots) e i pesi $w_k$ erano affinati da partite di grandi maestri. Combinato con alfa-beta pruning a profondità 12--17, Deep Blue sconfisse il campione mondiale Garry Kasparov nel 1997 con punteggio $3.5$--$2.5$.
\end{example}

\begin{remark}
La qualità della strategia dipende \emph{criticamente} dalla funzione di valutazione: un limite di profondità elevato con una funzione mediocre può dare risultati peggiori di un limite basso con una funzione eccellente. Le moderne reti neurali profonde (AlphaZero, 2017) rimpiazzano le feature hand-crafted con valutazioni apprese per self-play, raggiungendo livelli sovrumani a scacchi, shogi e go.
\end{remark}

\section{Backward induction nei giochi multi-stage}

I \textbf{giochi multi-stage} decompongono il gioco in $T$ stadi. Ad ogni stadio si gioca un sotto-gioco, il cui payoff (determinato dall'NE del sotto-gioco) si aggiunge al payoff globale.

\begin{example}[Gioco a due stadi con information set]
Riprendo il gioco dei due investitori (Capitolo~\ref{ch:stackelberg}). Nel primo stadio ciascun investitore sceglie simultaneamente (\emph{withdraw} o \emph{stay in}); per l'informazione imperfetta i due nodi del giocatore 2 sono raggruppati in un information set. Il secondo stadio è identico al primo.

Applicando la backward induction:
\begin{itemize}
  \item Secondo stadio: l'unico NE è \emph{(stay in, stay in)} con payoff $(10,10)$.
  \item Primo stadio: il payoff $(10,10)$ del secondo stadio si aggiunge ai payoff del primo. L'analisi del gioco ridotto produce due NE nel primo stadio: \emph{(stay in, stay in)} (payoff $20,20$) e \emph{(withdraw, withdraw)} (payoff $3+10=13$, $3+10=13$).
\end{itemize}
Con informazione imperfetta nel primo stadio (information set), rimane l'NE \emph{(stay in, stay in)} come equilibrio di cooperazione.
\end{example}

\section{Oligopolio con un leader e follower multipli}

Tre imprese producono lo stesso bene omogeneo, competendo sulla quantità. Le strategie sono $S_1=S_2=S_3=[0,+\infty)$ e le utilità:
\[
  u_i(x_1,x_2,x_3) = x_i\max\{T-(x_1+x_2+x_3),0\} - c\,x_i,\quad T>c.
\]
L'impresa 1 è il \textbf{leader}: sceglie $x_1$ per prima. Le imprese 2 e 3 sono i \textbf{follower}: osservano $x_1$ e reagiscono simultaneamente in un \textbf{duopolio di Cournot} con prezzo massimo $T-x_1$.

\subsection{Equilibrio dei follower}

Fissato $x_1$, le imprese 2 e 3 giocano un Cournot con prezzo di intercetta $T-x_1$. Le best responses simmetriche danno
\[
  x_2^*(x_1) = x_3^*(x_1) = \begin{cases}
    \dfrac{T-x_1-c}{3} & \text{se }x_1 < T-c\\
    0 & \text{altrimenti}
  \end{cases}
\]
e il profilo $(x_2^*(x_1), x_3^*(x_1))$ è l'unico NE del sottogioco dei follower.

\subsection{Scelta ottimale del leader}

Il leader massimizza $u_1(x_1, x_2^*(x_1), x_3^*(x_1))$ sostituendo le best responses:
\[
  u_1 = x_1\Bigl(T - x_1 - \tfrac{2(T-x_1-c)}{3}\Bigr) - c\,x_1
       = x_1\cdot\frac{T-x_1-c}{3} - c\,x_1.
\]
Ottimizzando in $x_1$ (continuo): $x_1^* = \frac{T-c}{2}$, e quindi $x_2^*=x_3^*=\frac{T-c}{6}$. L'equilibrio globale è
\[
  \Bigl(\frac{T-c}{2},\;\frac{T-c}{6},\;\frac{T-c}{6}\Bigr).
\]

\subsection{Tabella comparativa degli oligopoli}

\begin{center}
\begin{tabular}{lcccc}
\toprule
Struttura & Produzione totale & Prezzo & Payoff per impresa & Payoff totale \\
\midrule
Monopolio & $\tfrac{T-c}{2}$ & $\tfrac{T+c}{2}$ & $\tfrac{(T-c)^2}{4}$ & $\tfrac{(T-c)^2}{4}$ \\
Duopolio Cournot & $\tfrac{2(T-c)}{3}$ & $\tfrac{T+2c}{3}$ & $\tfrac{(T-c)^2}{9}$ & $\tfrac{2(T-c)^2}{9}$ \\
Duopolio Stackelberg & $\tfrac{3(T-c)}{4}$ & $\tfrac{T+3c}{4}$ & L: $\tfrac{(T-c)^2}{8}$, F: $\tfrac{(T-c)^2}{16}$ & $\tfrac{3(T-c)^2}{16}$ \\
1 leader, 2 follower & $\tfrac{5(T-c)}{6}$ & $\tfrac{T+5c}{6}$ & L: $\tfrac{(T-c)^2}{12}$, F: $\tfrac{(T-c)^2}{36}$ & $\tfrac{5(T-c)^2}{36}$ \\
\bottomrule
\end{tabular}
\end{center}

Osservazione: all'aumentare della concorrenza, la produzione totale cresce, il prezzo scende (bene per i consumatori) ma il payoff per impresa e totale diminuisce; il leader ha sempre vantaggio sul follower.

\section*{In sintesi per l'orale}
\addcontentsline{toc}{section}{In sintesi per l'orale (L12)}
\begin{itemize}
  \item \textbf{Alpha-beta}: $\alpha$ $=$ il valore minimo già garantito a MAX lungo il cammino, $\beta$ $=$ il massimo già concesso a MIN. Si pota il sottoalbero appena $\alpha\ge\beta$, perché quel ramo non potrà influenzare la scelta alla radice.
  \item Il risultato è \emph{esatto}: alpha-beta restituisce lo stesso valore minimax, solo visitando meno nodi.
  \item Complessità: da $O(b^d)$ a $O(b^{d/2})$ con ordinamento ottimo delle mosse, cioè profondità doppia a parità di tempo; con ordinamento casuale circa $O(b^{3d/4})$.
  \item Quando l'albero resta troppo grande: \emph{deep guessing} con funzioni di valutazione sulle posizioni non terminali (Deep Blue).
  \item Oligopolio con un leader e più follower: prima si calcola l'equilibrio (di Nash) tra i follower dato $x_1$, poi il leader ottimizza sopra. Tabella comparativa Cournot / Stackelberg / leader con $n$ follower: al crescere dei follower la quantità totale cresce e il prezzo cala.
\end{itemize}

% ============================================================
\chapter{Stackelberg-Nash Games (L13)}
\label{ch:stackelberg-nash}
% ============================================================

\section{Definizione e motivazione}

I \textbf{Stackelberg-Nash games} (single-leader multi-follower games) generalizzano il duopolio di Stackelberg a $n$ follower che reagiscono simultaneamente all'scelta del leader, giocando a loro volta un gioco di Nash.

\begin{defbox}{Stackelberg-Nash Game}
Un gioco Stackelberg-Nash è una tripla $(N,(S_i)_{i\in N},(u_i)_{i\in N})$ con $N=\{1\}\cup F$, $F=\{2,\dots,n+1\}$, dove:
\begin{itemize}
  \item Il \textbf{leader} (giocatore 1) sceglie $x_1\in S_1$ per primo.
  \item I \textbf{follower} $i\in F$ osservano $x_1$ e reagiscono \emph{simultaneamente}, giocando il gioco strategico
    \[G(x_1) = \bigl(F,\;(S_i)_{i\in F},\;(u_i(x_1,\cdot))_{i\in F}\bigr).\]
  \item Il leader \emph{anticipa} l'esito del gioco dei follower.
\end{itemize}
\end{defbox}

\textbf{Working assumption:} $\NE(x_1) = \{\text{Nash equilibria di }G(x_1)\}\subseteq S_{-1}$ è un singleton per ogni $x_1\in S_1$.

\begin{defbox}{Stackelberg-Nash Equilibrium (Sherali--Soyster--Murphy, 1983)}
Se la working assumption vale, $x_1^*\in S_1$ è una \textbf{Stackelberg-Nash solution} se
\[
  x_1^* \in \argmax\{u_1(x_1,\,\NE(x_1)) : x_1\in S_1\}.
\]
Il profilo $(x_1^*,\NE(x_1^*))$ è un \textbf{Stackelberg-Nash equilibrium}.
\end{defbox}

\section{Best response dynamics in giochi Stackelberg-Nash}

In un gioco finito, il leader utilizza la BRD:
\begin{enumerate}
  \item Il leader sceglie una strategia $x_1^{(k)}$.
  \item I follower reagiscono con $\NE(x_1^{(k)})$.
  \item Il leader risponde con $x_1^{(k+1)}\in\BR_1(\NE(x_1^{(k)}))$.
\end{enumerate}
Questo processo può convergere a un NE oppure ciclare (vedere esempio numerico con 4 strategie per il leader e 4 per coppia di follower).

\section{Mancanza di unicità dell'equilibrio dei follower}

Se $\NE(x_1)$ non è singleton, $u_1(x_1,\NE(x_1))$ non è ben definita. Si introducono due varianti:

\begin{defbox}{Optimistic Stackelberg-Nash Problem (OSN)}
\[
  \max\{u_1(x_1,x_{-1}) : x_1\in S_1,\; x_{-1}\in\NE(x_1)\}. \tag{OSN}
\]
I punti di massimo $(x_1^*,x_{-1}^*)$ sono gli \textbf{equilibri ottimistici}: il leader assume che i follower scelgano l'NE più favorevole per lui.
\end{defbox}

\begin{defbox}{Pessimistic Stackelberg-Nash Problem (PSN)}
\[
  \max\{\min\{u_1(x_1,x_{-1}) : x_{-1}\in\NE(x_1)\} : x_1\in S_1\}. \tag{PSN}
\]
Gli \textbf{equilibri pessimistici} sono definiti da: $x_1^*$ è il punto di massimo di (PSN) e $x_{-1}^*\in\NE(x_1^*)$.
\end{defbox}

\section{Teoremi di esistenza}

\begin{thmbox}{Esistenza in giochi finiti}
Sia $(N,(S_i),(u_i))$ un gioco Stackelberg-Nash finito con $N=\{1\}\cup F$. Se il gioco $G(x_1)$ ha almeno un NE per ogni $x_1\in S_1$, allora esistono almeno un equilibrio Stackelberg-Nash ottimistico e uno pessimistico.
\end{thmbox}
\textit{Dimostrazione.} (OSN) e (PSN) sono problemi di ottimizzazione su insiemi finiti, quindi i massimi esistono (il dominio è finito). \qed

\begin{remark}
Se $\NE(x_1)=\emptyset$ per qualche $x_1$, allora (OSN) e (PSN) possono non essere ben definiti. Esempio: nella matrice del gioco con 3 strategie per follower 2 e 3, la strategia $\mathbf{2}$ del leader non ha NE.
\end{remark}

\begin{thmbox}{Esistenza per giochi continui}
Sia $(N,(S_i),(u_i))$ un gioco Stackelberg-Nash tale che:
\begin{itemize}
  \item[(i)] $S_1\subseteq\mathbb{R}^{m_1}$ è compatto;
  \item[(ii)] $u_1$ è continua;
  \item[(iii)] Per ogni $i\in F$: $S_i\subseteq\mathbb{R}^{m_i}$ è convesso e compatto;
  \item[(iv)] $u_i$ è continua per $i\in F$;
  \item[(v)] L'insieme delle best responses $R_i(x_{-i})$ è convesso per ogni $x_{-i}$.
\end{itemize}
Allora il gioco ha almeno un equilibrio Stackelberg-Nash ottimistico.
\end{thmbox}
\textit{Dimostrazione (sketch).} Il teorema di Nikaido--Isoda garantisce $\NE(x_1)\ne\emptyset$ per ogni $x_1$ (condizioni iii, iv, v). Il teorema di Weierstrass applicato a (OSN) -- che è continuo su $S_1$ compatto -- garantisce l'esistenza del massimo. \qed

\begin{corollary}
Ogni gioco Stackelberg-Nash finito ha almeno un equilibrio ottimistico in strategie miste.
\end{corollary}

\section{Pool mining nelle blockchain}

\begin{example}[Pool Mining]
Un leader (\textit{platform}) offre una ricompensa $r$ per il mining di un blocco di valore $g$. I follower sono $n$ miner che forniscono potenza di hashing $p_i$. L'utilità di ogni miner è
\[
  u_i(r,p_i,p_{-i}) = \frac{p_i}{p_1+\cdots+p_n}\,r - c_i\,p_i
\]
dove $c_i$ è il costo unitario di potenza per il miner $i$. I miner giocano un oligopolio di Cournot con funzione di domanda inversa $p(t)=r/t$ (dove $t=\sum_j p_j$). L'utilità del leader è
\[
  u_1(r,p) = \frac{p_1+\cdots+p_n}{T+p_1+\cdots+p_n}\,g - r
\]
con $T$ = potenza totale della rete. Nel caso omogeneo $c_i\equiv c$:
\begin{itemize}
  \item $G(r)$ è un gioco a potenziale ordinale.
  \item L'unico NE dei miner è $p_i=r(n-1)/cn^2$ per ogni $i$.
  \item La ricompensa ottimale per il leader è:
    \[
      r^* = \begin{cases}
        0 & \text{se }g\le cT\,n/(n-1)\\
        c\Bigl(\sqrt{gT(n-1)/cn} - T\Bigr)\frac{n}{n-1} & \text{altrimenti.}
      \end{cases}
    \]
\end{itemize}
\end{example}

\section{Giochi multi-leader common-follower}

\begin{defbox}{Multi-leader Common-follower Game}
Un gioco con $n$ leader e 1 follower: $N=L\cup\{n+1\}$, $L=\{1,\dots,n\}$. I leader scelgono $x_L=(x_1,\dots,x_n)\in S_L=S_1\times\cdots\times S_n$ \emph{simultaneamente}; il follower osserva $x_L$ e risponde con $x_{n+1}\in\argmax_{x_{n+1}}u_{n+1}(x_L,x_{n+1})$.

Se $R_{n+1}(x_L)$ è singleton, i leader giocano il gioco
\[
  (L,\;(S_i)_{i\in L},\;(u_i(x_i,x_{-i}^L,R_{n+1}(x_L)))_{i\in L}).
\]
\end{defbox}

\begin{defbox}{Generalized Stackelberg-Nash Equilibrium (Sherali, 1984)}
$x_L^*\in S_L$ è una \textbf{soluzione Stackelberg-Nash generalizzata} se ogni leader $i\in L$ soddisfa
\[
  x_i^*\in\argmax\{u_i(x_i,x_{-i}^{L*},R_{n+1}(x_i,x_{-i}^{L*})) : x_i\in S_i\}.
\]
Il profilo $(x_L^*,R_{n+1}(x_L^*))$ è un \textbf{Stackelberg-Nash equilibrium generalizzato}.
\end{defbox}

\subsection{Non-unicità e giochi indotti multipli}

Il problema fondamentale nei giochi multi-leader common-follower è che la best response del follower $R_{n+1}(x_L)$ può non essere un singleton. Quando ciò accade per almeno una coppia di strategie $(x_1,x_2)\in S_1\times S_2$, si presenta ambiguità: i leader non possono determinare univocamente il gioco indotto tra loro.

\begin{example}[Quattro giochi indotti]
Supponiamo due leader (I e II) con 2 strategie ciascuno e un follower con 4 strategie. Per ogni profilo $(x_I,x_{II})\in\{1,2\}\times\{1,2,3\}$, i leader calcolano la best response del follower $R_f(x_I,x_{II})$. Se per la coppia $(x_I,x_{II})=(1,1)$ la risposta del follower è ambigua (ad esempio strategie 1 o 2 sono entrambe ottimali), il valore della cella per i leader è incerto.

I leader anticipano le risposte del follower su tutte le celle non ambigue, costruendo una matrice sommario; le celle ambigue rimangono un punto interrogativo ``?''. A seconda della scelta del follower per le celle ambigue, si inducono \textbf{4 giochi diversi} tra i leader (uno per ogni combinazione di risposte nelle celle ambigue), ognuno con potenzialmente diversi equilibri di Nash o nessun equilibrio.

\begin{itemize}
  \item Alcuni dei 4 giochi indotti possono avere NE; altri no.
  \item Nessuno dei giochi indotti potrebbe ammettere un equilibrio Stackelberg-Nash generalizzato.
  \item Le formulazioni OSN e PSN (ottimistica e pessimistica) gestiscono questa ambiguità scegliendo, rispettivamente, la risposta del follower più favorevole e più sfavorevole per il leader.
\end{itemize}
\end{example}

\begin{remark}
La mancanza di unicità della risposta del follower è un problema \textbf{strutturale} nei giochi multi-leader common-follower. Non dipende dalla specificità dell'esempio, ma emerge naturalmente ogniqualvolta il follower ha più best responses equivalent. Le formulazioni OSN e PSN sono gli strumenti standard per gestire questa incertezza.
\end{remark}

% ============================================================
%  APPENDICE A – Bargaining
% ============================================================

% ============================================================
%  Capitoli L14 -- L16 : Bargaining
\section*{In sintesi per l'orale}
\addcontentsline{toc}{section}{In sintesi per l'orale (L13)}
\begin{itemize}
  \item \textbf{Stackelberg--Nash}: un leader, più follower che tra loro giocano un gioco di Nash parametrico in $x_1$. Struttura a due livelli: $\max_{x_1}u_1(x_1,y(x_1))$ con $y(x_1)$ equilibrio dei follower.
  \item Se l'equilibrio dei follower non è unico, il problema del leader non è ben posto: si sceglie la lettura \emph{ottimistica} (strong Stackelberg) o \emph{pessimistica} (weak Stackelberg); quest'ultima è quella che garantisce, la prima quella che spesso esiste.
  \item Esistenza: servono compattezza, continuità e una selezione misurabile della mappa degli equilibri; nella lettura pessimistica l'esistenza è più delicata perché la funzione valore può essere solo semicontinua.
  \item BRD in questi giochi: si alternano l'aggiornamento del leader e la ri-convergenza dei follower; può non convergere se i follower hanno equilibri multipli.
  \item Applicazione del corso: \textbf{pool mining} nelle blockchain (il pool è il leader, i miner i follower) e i giochi multi-leader common-follower, dove la non unicità genera più giochi indotti.
\end{itemize}

% ============================================================

\chapter{Bargaining assiomatico: la soluzione di Nash (L14)}
\label{ch:bargaining-assiomatico}
% ============================================================

\section{Il problema di bargaining}

Il \textbf{bargaining} (contrattazione) descrive una situazione in cui $n$ giocatori devono
\emph{mettersi d'accordo} su un esito, cioè su un vettore di pay-off, uno per giocatore.
A differenza dei giochi strategici, qui non si modellano le mosse: si modella soltanto
\emph{che cosa} è raggiungibile e \emph{che cosa} succede se l'accordo salta.

\begin{defbox}{Bargaining problem (Nash, 1950)}
Un \textbf{problema di bargaining} è una coppia $(U,u^*)$ dove
\begin{itemize}
  \item $U\subseteq\mathbb{R}^n$ è l'insieme degli \textbf{outcome fattibili} (vettori di pay-off
        ottenibili con un accordo);
  \item $u^*\in\mathbb{R}^n$ è il \textbf{disagreement point}: l'esito che si realizza se
        \emph{nessun} accordo viene raggiunto.
\end{itemize}
Gli \textbf{assiomi di Nash} sul dato sono:
\begin{itemize}
  \item[(i)] $U$ è convesso e compatto;
  \item[(ii)] esiste $u\in U$ con $u>u^*$, cioè $u_i>u^*_i$ per ogni $i=1,\dots,n$.
\end{itemize}
\end{defbox}

L'assioma (i) è naturale quando gli esiti si ottengono randomizzando fra accordi
(la convessità è la possibilità di \emph{lotterie} fra accordi); l'assioma (ii) dice che
c'è un guadagno \emph{stretto} da cooperare per tutti: altrimenti il problema è degenere.

\begin{defbox}{Bargaining function}
Sia $\mathcal{B}=\{(U,u^*) : \text{valgono (i) e (ii)}\}$ la famiglia dei bargaining games.
Una \textbf{bargaining function} è una qualsiasi mappa
\[
  \psi:\mathcal{B}\longrightarrow\mathbb{R}^n .
\]
È un \emph{solution concept}: un meccanismo (o arbitrato) che a ogni problema associa
un esito.
\end{defbox}

\section{Due sorgenti di problemi di bargaining}

\begin{example}[Cooperazione in giochi non cooperativi]
Se i giocatori di un gioco finito possono \emph{accordarsi} su quali strategie miste giocare,
l'insieme degli esiti raggiungibili è l'inviluppo convesso dei vettori di pay-off.
Per il gioco
\[
\begin{array}{c|ccc}
 \text{I}/\text{II} & \ell_2 & c_2 & r_2\\\hline
 \ell_1 & (2,2) & (3,0) & (4,1)\\
 r_1    & (1,0) & (1,1) & (3,1)
\end{array}
\]
si ottiene
\[
  U=\mathrm{conv}\{(2,2),(3,0),(4,1),(1,0),(1,1),(3,1)\},
\]
che è convesso e compatto per costruzione: l'assioma (i) è automaticamente soddisfatto.
Il disagreement point si sceglie tipicamente come il pay-off che ciascuno si garantisce
da solo (maxmin) o come un equilibrio di Nash del gioco.
\end{example}

\begin{example}[Bankruptcy]
Un'azienda fallisce; l'attivo residuo vale $a>0$ e ci sono $n$ creditori, con $d_i$ il
debito verso il creditore $i$ e $d_1+\dots+d_n>a$ (l'attivo non basta). L'insieme degli
esiti di un accordo/arbitrato è
\[
  U=\Bigl\{u\in\mathbb{R}^n_+ : \sum_{k=1}^n u_k\le a,\ u_i\le d_i\ \ i=1,\dots,n\Bigr\},
\]
un politopo convesso e compatto, con $u^*=(0,\dots,0)$. La domanda ``come dividere
l'attivo fra i creditori?'' è esattamente la scelta di una bargaining function.
\end{example}

\section{Proprietà desiderabili per una bargaining function}

\begin{defbox}{Assiomi su $\psi$}
Sia $\psi$ una bargaining function.
\begin{description}
  \item[feasibility] $\psi(U,u^*)\in U$;
  \item[rationality] $\psi(U,u^*)\ge u^*$ (nessuno accetta meno del disaccordo);
  \item[Pareto optimality] $u\in U,\ u\ge\psi(U,u^*)\implies u=\psi(U,u^*)$;
  \item[independence of irrelevant alternatives (IIA)] se $(U,u^*),(\widehat U,u^*)\in\mathcal{B}$,
        $\widehat U\subseteq U$ e $\psi(U,u^*)\in\widehat U$, allora
        $\psi(\widehat U,u^*)=\psi(U,u^*)$;
  \item[covariance sotto trasformazioni affini] per $\alpha_i>0$, $\beta_i\in\mathbb{R}$ e
        $\varphi(u)=(\alpha_iu_i+\beta_i)_i$ si ha
        $\psi(\varphi(U),\varphi(u^*))=\varphi(\psi(U,u^*))$;
  \item[symmetry] se $U$ è $ij$-simmetrico e $u^*_i=u^*_j$, allora
        $\psi_i(U,u^*)=\psi_j(U,u^*)$.
\end{description}
Qui $U$ è \textbf{$ij$-simmetrico} se
$(u_1,\dots,u_i,\dots,u_j,\dots,u_n)\in U\iff(u_1,\dots,u_j,\dots,u_i,\dots,u_n)\in U$.
\end{defbox}

\begin{remark}
La covarianza affine dice che la soluzione non deve dipendere dalla scala e dall'origine con
cui ciascun giocatore misura la propria utilità: applicare $\varphi$ ai dati del problema
equivale ad applicare $\varphi$ alla soluzione.
\end{remark}

\begin{remark}[Rilevanza della Pareto optimality]
Se $u^*\in U$, la funzione costante $\psi\equiv u^*$ soddisfa \emph{tutte} le proprietà
elencate tranne la Pareto optimality. È il ``solution concept del disaccordo'': banale,
ma mostra che senza efficienza gli assiomi non hanno mordente.
\end{remark}

\section{Il teorema di Nash}

\begin{thmbox}{Nash solution (Nash, 1950)}
Esiste un'unica bargaining function che soddisfa tutte e sei le proprietà precedenti, ed è
\[
  \psi^{na}(U,u^*)=\argmax\Bigl\{\prod_{k=1}^n (x_k-u^*_k)\ :\ x\in U,\ x\ge u^*\Bigr\}.
\]
\end{thmbox}

\noindent Osservazioni sull'enunciato:
\begin{itemize}
  \item il punto di massimo \emph{esiste ed è unico}: la funzione $g(x)=\prod_k(x_k-u^*_k)$
        è continua sul compatto $\{x\in U : x\ge u^*\}$ (esistenza per Weierstrass), e
        \[
          \argmax\Bigl\{\prod_{k=1}^n(x_k-u^*_k)\Bigr\}
          =\argmax\Bigl\{\sum_{k=1}^n\log(x_k-u^*_k)\Bigr\}
        \]
        sul sottoinsieme dove tutte le componenti sono $>u^*_k$ (assioma (ii) garantisce che
        non è vuoto): la funzione a destra è \emph{strettamente concava} su un insieme
        convesso, quindi il massimo è unico;
  \item $\psi\equiv u^*$ è l'\emph{unica} altra funzione che soddisfa le sei proprietà
        rinunciando alla Pareto optimality, insieme a tutte le combinazioni convesse delle due.
\end{itemize}

\textit{Dimostrazione (sketch).}
\emph{(a) $\psi^{na}$ soddisfa gli assiomi.} Feasibility e rationality sono per costruzione;
la Pareto optimality segue perché $g$ è strettamente crescente in ogni coordinata su
$\{x>u^*\}$; la IIA è immediata (se il massimizzatore su $U$ sta in $\widehat U\subseteq U$,
è anche il massimizzatore su $\widehat U$); la covarianza affine segue da
$\prod_k(\alpha_kx_k+\beta_k-\alpha_ku^*_k-\beta_k)=\bigl(\prod_k\alpha_k\bigr)\prod_k(x_k-u^*_k)$,
cioè $g$ cambia per una costante positiva e l'$\argmax$ non si muove; la simmetria segue
dalla simmetria di $g$ e dall'unicità del massimo.

\emph{(b) Unicità.} L'unicità della bargaining function che soddisfa i sei assiomi è
la parte non banale dell'enunciato di Nash, ed è data per acquisita nel corso. \qed

\begin{remark}[Significato geometrico]
Per $n=2$ la quantità $g(x)=(x_1-u^*_1)(x_2-u^*_2)$ è l'\textbf{area del rettangolo} di
vertici opposti $u^*$ e $x$: la soluzione di Nash è il punto della frontiera di $U$ che
\textbf{massimizza quest'area}.
\end{remark}

\section{Proprietà aggiuntive della soluzione di Nash}

Vale sempre
\[
  \psi^{na}(U,u^*)=\psi^{na}\bigl((U-\mathbb{R}^n_+)\cap(u^*+\mathbb{R}^n_+),\,u^*\bigr),
\]
cioè si può sempre rimpiazzare $U$ con la sua chiusura verso il basso a partire da $u^*$
senza cambiare la soluzione.

\begin{defbox}{Comprehensiveness}
$U$ è \textbf{$u^*$-comprehensive} se $(U-\mathbb{R}^n_+)\cap(u^*+\mathbb{R}^n_+)=U$, ovvero
\[
  x\in U,\ u^*\le y\le x\implies y\in U .
\]
(Si può sempre ``buttare via'' utilità, restando sopra il disagreement point.)
\end{defbox}

\begin{defbox}{Midpoint domination}
Posto $d^i=\argmax\{x_i : x\in U,\ x_{-i}=u^*_{-i}\}$ (il massimo che $i$ può ottenere
lasciando gli altri al disaccordo), $\psi$ soddisfa la \textbf{midpoint domination} se
\[
  \psi(U,u^*)\ \ge\ \sum_{i=1}^n \frac{d^i}{n}.
\]
Cioè: ogni giocatore ottiene almeno quanto gli darebbe la lotteria uniforme sulle
$n$ ``dittature''. La midpoint domination implica la rationality.
\end{defbox}

\begin{thmbox}{Nash e midpoint domination (Moulin, 1983)}
La soluzione di Nash è l'unica bargaining function che soddisfa feasibility,
midpoint domination e independence of irrelevant alternatives.
\end{thmbox}

\begin{remark}
È una caratterizzazione \emph{alternativa}: la simmetria e la covarianza affine sono
sostituite dalla sola midpoint domination, che è un requisito di equità più concreto
(un ``benchmark'' esplicito) ma più forte della rationality.
\end{remark}

\section{Critiche agli assiomi}

\begin{example}[Critica all'IIA (Luce--Raiffa, 1957)]
Sia $n=2$, $u^*=(0,0)$ e
\[
  U=\Bigl\{x\in\mathbb{R}^2_+:\ \tfrac{x_1}{10}+\tfrac{x_2}{4}\le 1\Bigr\},
  \qquad
  \widehat U=U\cap\{x_1\le 5\}\subseteq U .
\]
Su $U$: massimizzando $x_1x_2$ con $x_2=4(1-x_1/10)$ si trova $x_1=5$, $x_2=2$, quindi
$\psi^{na}(U,u^*)=(5,2)$. Poiché $(5,2)\in\widehat U$, per IIA anche
$\psi^{na}(\widehat U,u^*)=(5,2)$.

\emph{Il punto della critica.} In $U$ le aspettative massime sono $m(U)=(10,4)$: la
soluzione $(5,2)$ dà a ciascuno esattamente metà della propria aspettativa massima, e
sembra equa. In $\widehat U$ le aspettative massime sono $m(\widehat U)=(5,4)$: la stessa
soluzione $(5,2)$ dà al giocatore~1 \emph{tutto} ciò che può sperare e al giocatore~2
solo la metà. È davvero ``irrilevante'' ciò che è stato tolto?
\end{example}

\begin{example}[Critica alla mancanza di monotonicità (Kalai--Smorodinsky, 1975)]
Con gli stessi $U$ e $\widehat U$: $\widehat U\subseteq U$, eppure la soluzione di Nash non
peggiora per il giocatore~1. In generale si costruiscono $\widehat U\subseteq U$ con
soluzioni di Nash $b$ e $a$ tali che $a_1>b_1$ ma $b_2>a_2$: il giocatore~2 \emph{guadagna}
dal fatto che l'insieme si sia \emph{ristretto}, benché per ogni $\hat u\in\widehat U$ con
$\hat u>u^*$ esista $u\in U$ con $u>\hat u$ (``$U$ è uniformemente migliore''). Questo è
percepito come paradossale.
\end{example}

\section{Weak monotonicity e la soluzione di Kalai--Smorodinsky}

\begin{defbox}{Utopia point e weak monotonicity}
Il \textbf{punto utopia} di $(U,u^*)$ è $m(U)\in\mathbb{R}^n$ con
\[
  m(U)_i=\max\{x_i\in\mathbb{R} : (x_i,x_{-i})\in U \text{ per qualche } x_{-i}\in\mathbb{R}^{n-1}\}.
\]
$\psi$ soddisfa la \textbf{weak monotonicity} se
\[
  (\widehat U,u^*),(U,u^*)\in\mathcal{B},\ \widehat U\subseteq U,\ m(\widehat U)=m(U)
  \implies \psi(\widehat U,u^*)\le\psi(U,u^*).
\]
\end{defbox}

\begin{thmbox}{Soluzione di Kalai--Smorodinsky (1975), $n=2$}
$\psi^{ks}(U,u^*)$ è l'intersezione del segmento fra $u^*$ e il punto utopia $m(U)$ con la
frontiera (di Pareto) di $U$. È l'unica bargaining function che soddisfa feasibility,
rationality, Pareto optimality, symmetry e weak monotonicity.
\end{thmbox}

L'idea normativa è esplicita: \emph{i guadagni rispetto a $u^*$ sono proporzionali ai
massimi guadagni possibili},
\[
  \frac{\psi^{ks}_1-u^*_1}{m(U)_1-u^*_1}=\frac{\psi^{ks}_2-u^*_2}{m(U)_2-u^*_2}.
\]

\begin{exbox}{Nash contro Kalai--Smorodinsky (esempio delle slide)}
Con $U=\{x\in\mathbb{R}^2_+ : x_1/10+x_2/4\le1\}$, $u^*=(0,0)$: $m(U)=(10,4)$, la retta
$x_2=\tfrac{4}{10}x_1$ interseca la frontiera in $x_1/10+x_1/10=1$, cioè $x_1=5$, $x_2=2$:
\[
  \psi^{ks}(U,u^*)=(5,2)=\psi^{na}(U,u^*).
\]
Con $\widehat U=U\cap\{x_1\le5\}$: $m(\widehat U)=(5,4)$, la retta è $x_2=\tfrac45x_1$ e
sulla frontiera $\tfrac{x_1}{10}+\tfrac{1}{4}\cdot\tfrac45x_1=\tfrac{3}{10}x_1=1$, da cui
\[
  \psi^{ks}(\widehat U,u^*)=\Bigl(\tfrac{10}{3},\tfrac{8}{3}\Bigr)\ne(5,2)=\psi^{na}(\widehat U,u^*).
\]
KS ridistribuisce a favore del giocatore~2, che nel problema ristretto ha una aspettativa
relativamente più alta: è esattamente la risposta alla critica di Luce--Raiffa.
\end{exbox}

\begin{thmbox}{Risultato di impossibilità (Roth, 1979)}
Se $n\ge3$ non esiste alcuna bargaining function che soddisfi feasibility, rationality,
Pareto optimality, symmetry e weak monotonicity.
\end{thmbox}

\textit{Dimostrazione (sketch, controesempio con $n=3$).} Siano $u^*=(0,0,0)$,
\[
  S=\mathrm{conv}\{(0,0,0),(1,0,1),(0,1,1)\},\qquad
  T=\{u\in\mathbb{R}^3_+ : u_i\le1,\ u_1+u_2+u_3\le2\}.
\]
Allora $S\subseteq T$ e $m(S)=m(T)=(1,1,1)$. Il problema $(T,u^*)$ è simmetrico, quindi per
symmetry e Pareto optimality $\psi(T,u^*)=(\alpha,\alpha,\alpha)$ con $3\alpha=2$, cioè
$\psi(T,u^*)=(\tfrac23,\tfrac23,\tfrac23)$. Ogni $s\in S$ si scrive
$s=(\lambda_2,\lambda_3,\lambda_2+\lambda_3)$ con $\lambda_i\ge0$,
$\lambda_1+\lambda_2+\lambda_3=1$; la Pareto optimality forza $\lambda_1=0$, dunque
$\lambda_2+\lambda_3=1$ e $[\psi(S,u^*)]_3=1$. Ma la weak monotonicity richiederebbe
$\psi(S,u^*)\le\psi(T,u^*)$, cioè $1\le\tfrac23$: assurdo. \qed

\section{Strong monotonicity e la soluzione egalitaria}

\begin{defbox}{Strong monotonicity e weak Pareto optimality}
\begin{description}
  \item[strong monotonicity] $(\widehat U,u^*),(U,u^*)\in\mathcal{B},\ \widehat U\subseteq U
        \implies\psi(\widehat U,u^*)\le\psi(U,u^*)$ (senza richiedere $m(\widehat U)=m(U)$);
  \item[weak Pareto optimality] $u>\psi(U,u^*)\implies u\notin U$ (nessun miglioramento
        \emph{stretto per tutti}).
\end{description}
\end{defbox}

\begin{defbox}{Soluzione egalitaria}
\[
  \psi^{eg}(U,u^*)=u^*+\tau\,(1,\dots,1),\qquad
  \tau=\min\bigl\{\max\{x_i-u^*_i : x\in U\} : i=1,\dots,n\bigr\}.
\]
Tutti guadagnano \emph{lo stesso} rispetto al disagreement point.
\end{defbox}

\begin{thmbox}{Strong monotonicity implica egalitarismo (Kalai, 1977)}
La soluzione egalitaria è l'unica bargaining function che soddisfa feasibility, rationality,
weak Pareto optimality, symmetry e strong monotonicity, purché $U$ sia $u^*$-comprehensive.
\end{thmbox}

\begin{remark}
Si noti l'indebolimento: la Pareto optimality piena è sostituita da quella debole, e in
compenso la monotonicità è rafforzata da debole a forte.

Sostituendo $(1,\dots,1)$ con un qualunque $p>0$ si ottiene la \textbf{soluzione
proporzionale}: in sostanza \emph{strong monotonicity $\simeq$ proporzionalità}.
\end{remark}

\section*{In sintesi per l'orale}
\addcontentsline{toc}{section}{In sintesi per l'orale (L14)}
\begin{itemize}
  \item Un bargaining problem è $(U,u^*)$ con $U$ convesso e compatto e almeno un $u\in U$
        con $u>u^*$; una bargaining function $\psi:\mathcal{B}\to\mathbb{R}^n$ è un solution
        concept \emph{assiomatico} (non si modellano le mosse, solo gli esiti).
  \item Sei assiomi: feasibility, rationality, Pareto optimality, IIA, covarianza affine,
        simmetria. Teorema di Nash: esiste \emph{una sola} $\psi$ che li soddisfa, il
        massimizzatore del prodotto $\prod_k(x_k-u^*_k)$ (equivalentemente di
        $\sum_k\log(x_k-u^*_k)$: stretta concavità $\Rightarrow$ unicità).
  \item Geometricamente ($n=2$) la soluzione di Nash massimizza l'area del rettangolo di
        vertici $u^*$ e $x$.
  \item Senza Pareto optimality gli assiomi sono soddisfatti anche da $\psi\equiv u^*$ (e da
        tutte le combinazioni convesse con $\psi^{na}$): la Pareto optimality è ciò che
        rende il teorema non banale.
  \item Caratterizzazione alternativa (Moulin 1983): feasibility + midpoint domination + IIA.
  \item Critiche: Luce--Raiffa attacca l'IIA (esempio $U$ vs $\widehat U=U\cap\{x_1\le5\}$,
        soluzione $(5,2)$ in entrambi); Kalai--Smorodinsky attacca la mancanza di monotonicità.
  \item KS ($n=2$): intersezione fra il segmento $u^*\,m(U)$ e la frontiera; caratterizzata da
        feasibility, rationality, Pareto optimality, symmetry, weak monotonicity. Per $n\ge3$
        quel pacchetto di assiomi è \emph{inconsistente} (Roth 1979; controesempio con
        $S=\mathrm{conv}\{0,(1,0,1),(0,1,1)\}\subseteq T$).
  \item Rafforzando a strong monotonicity e indebolendo a weak Pareto optimality si ottiene la
        soluzione egalitaria (Kalai 1977), che però perde la covarianza affine.
  \item \textbf{Collegamenti}: la logica ``assiomi $\Rightarrow$ soluzione unica'' è la stessa
        della caratterizzazione del valore di $\Sh$ nei giochi TU; il bankruptcy problem è
        anche un gioco TU (e lì si confronta con nucleolo e Shapley); il disagreement point
        gioca il ruolo che nel core hanno i valori $v(\{i\})$ delle coalizioni singoletto.
\end{itemize}


% ============================================================
\chapter{Nash generalizzato e Nash program (L15)}
\label{ch:nash-program}
% ============================================================

\section{Richiamo e indebolimento degli assiomi}

Si riprende il problema cooperativo di bargaining $(U,u^*)\in\mathcal{B}$ con gli assiomi (i)
e (ii) del capitolo precedente. Qui si usa la lista di proprietà leggermente \emph{indebolita}:
feasibility, rationality, \textbf{weak} Pareto optimality
($u>\psi(U,u^*)\implies u\notin U$), independence of irrelevant alternatives,
\textbf{affine positive covariance} e symmetry.

Il punto interessante è: che cosa si ottiene se si \emph{rinuncia alla simmetria}?

\section{La soluzione di Nash generalizzata}

\begin{thmbox}{Generalized Nash solution}
Se una bargaining function $\psi$ soddisfa feasibility, rationality, weak Pareto optimality,
affine positive covariance e independence of irrelevant alternatives, allora esistono
\textbf{bargaining powers} $\alpha_k\ge0$ tali che
\[
  \psi(U,u^*)=\argmax\Bigl\{\prod_{k=1}^n(x_k-u^*_k)^{\alpha_k}\ :\ x\in U,\ x\ge u^*\Bigr\}.
\]
\end{thmbox}

\begin{itemize}
  \item Se tutti gli $\alpha_k$ sono uguali si ritrova esattamente la soluzione di Nash: la
        \emph{simmetria} è precisamente ciò che impone pesi uguali.
  \item Come prima si può passare ai logaritmi:
        \[
          \argmax\Bigl\{\prod_{k=1}^n(x_k-u^*_k)^{\alpha_k}\Bigr\}
          =\argmax\Bigl\{\sum_{k=1}^n\alpha_k\log(x_k-u^*_k)\Bigr\}
        \]
        su $\{x\in U : x\ge u^*\}$; la funzione a destra è concava, strettamente concava nelle
        coordinate con $\alpha_k>0$.
  \item Gli $\alpha_k$ misurano il \emph{potere contrattuale} del giocatore $k$: non derivano
        dagli assiomi, sono un dato ulteriore del modello (chi ha $\alpha_k=0$ viene di fatto
        ignorato dall'ottimizzazione, e ottiene solo ciò che gli garantiscono rationality e
        weak Pareto optimality).
\end{itemize}

\begin{exbox}{Divide the euro con bargaining powers}
Sia $n=2$, $U=\{x\in\mathbb{R}^2_+ : x_1+x_2\le1\}$, $u^*=(0,0)$. Si massimizza
$x_1^{\alpha_1}x_2^{\alpha_2}$ con $x_2=1-x_1$, cioè
$h(x_1)=\alpha_1\log x_1+\alpha_2\log(1-x_1)$. Da $h'(x_1)=\frac{\alpha_1}{x_1}-\frac{\alpha_2}{1-x_1}=0$
si ottiene $\alpha_1(1-x_1)=\alpha_2x_1$, quindi
\[
  \psi(U,u^*)=\Bigl(\frac{\alpha_1}{\alpha_1+\alpha_2},\ \frac{\alpha_2}{\alpha_1+\alpha_2}\Bigr).
\]
L'euro si divide \emph{in proporzione ai bargaining powers}; con $\alpha_1=\alpha_2$ si
ritrova la soluzione di Nash $(\tfrac12,\tfrac12)$.
\end{exbox}

\section{Il Nash program}

La soluzione di Nash è costruita per via assiomatica: è un \emph{arbitrato}, non il
risultato di una interazione strategica. Il \textbf{Nash program} è il progetto di
\emph{giustificare} le soluzioni cooperative come equilibri di opportuni giochi
non cooperativi.

\begin{defbox}{Nash program}
Dare, per ogni solution concept cooperativo, un gioco non cooperativo esplicito
(``bargaining procedure'') i cui equilibri -- tipicamente equilibri perfetti nei sottogiochi
(SPE) -- coincidano con la soluzione cooperativa.
\end{defbox}

I modelli non cooperativi di bargaining trattati nel corso sono:
\begin{itemize}
  \item il \textbf{Nash demand game}: i giocatori dichiarano simultaneamente e
        indipendentemente il proprio pay-off, cercando una combinazione fattibile;
  \item i \textbf{giochi di bargaining sequenziale} (due giocatori che devono dividere 1
        euro):
        \begin{itemize}
          \item \textbf{ultimatum}: un'unica trattativa, offerta ``take-it-or-leave-it'';
          \item \textbf{two stage}: offerte alternate, un solo round di offerta per giocatore;
          \item \textbf{infinite horizon}: offerte alternate fino all'accordo (Rubinstein);
        \end{itemize}
  \item il modello di \textbf{bargaining legislativo di Baron--Ferejohn}: $n$ giocatori,
        proponente estratto a sorte a ogni round, approvazione a maggioranza qualificata.
\end{itemize}
Gli ultimi due gruppi sono sviluppati nel Capitolo~\ref{ch:bargaining-sequenziale}.

\section{Il Nash demand game}

\begin{defbox}{Nash demand game}
Dato un bargaining game $(U,u^*)$, il \textbf{Nash demand game} è il gioco strategico in cui
ogni giocatore \emph{dichiara simultaneamente} il pay-off che pretende; se le richieste non
sono congiuntamente fattibili, tutti ricevono il pay-off di disaccordo. Formalmente:
\[
  S_i=[u^*_i,m_i],\qquad
  m_i=\max\{x_i\in\mathbb{R} : (x_i,x_{-i})\in U\ \text{per qualche}\ x_{-i}\in\mathbb{R}^{n-1}\},
\]
cioè $S=S_1\times\dots\times S_n=[u^*,m]$ con $m=(m_1,\dots,m_n)=m(U)$ il punto utopia, e
\[
  u_i(x_i,x_{-i})=
  \begin{cases}
    x_i & \text{se } (x_i,x_{-i})\in U,\\
    u^*_i & \text{se } (x_i,x_{-i})\notin U.
  \end{cases}
\]
\end{defbox}

\begin{remark}[Struttura degli equilibri]
Direttamente dalla definizione dei pay-off: se $x\in U$ è Pareto ottimale e $x\ge u^*$,
nessun giocatore ha convenienza a deviare unilateralmente. Chiedere meno di $x_i$ dà
semplicemente meno; chiedere più di $x_i$ rende il profilo non fattibile e quindi porta al
pay-off di disaccordo $u^*_i\le x_i$. Il demand game ha dunque un \emph{continuo} di
equilibri di Nash lungo la frontiera di Pareto: la sola nozione di equilibrio di Nash non
individua un unico esito.
\end{remark}

\begin{exbox}{Demand game per ``divide the euro''}
$U=\{x\in\mathbb{R}^2_+:x_1+x_2\le1\}$, $u^*=(0,0)$, $m(U)=(1,1)$, quindi
$S_1=S_2=[0,1]$. Se il giocatore~II chiede $\tfrac12$, la best response di~I è chiedere
$\tfrac12$: chiedendo di più la coppia diventa infeasible e I prende $0$; chiedendo di meno
prende meno. Ogni coppia $(t,1-t)$ con $t\in[0,1]$ è un equilibrio.
\end{exbox}

\section{Bargaining sequenziale: l'ultimatum game}

\begin{defbox}{Ultimatum game}
\emph{Divide the euro}: due giocatori devono accordarsi su come dividere $1$ euro, in
\emph{un'unica} fase di negoziazione. Il giocatore~I propone una allocazione, il giocatore~II
accetta o rifiuta; in caso di mancato accordo entrambi ottengono $0$. Le strategie sono
\[
  S_1=[0,1]\quad (x=\text{frazione per il giocatore I}),\qquad
  S_2=\{\sigma:[0,1]\to\{\text{yes},\text{no}\}\},
\]
e i pay-off
\[
  u_1(x,\sigma)=\begin{cases}x&\sigma(x)=\text{yes}\\0&\sigma(x)=\text{no}\end{cases}
  \qquad
  u_2(x,\sigma)=\begin{cases}1-x&\sigma(x)=\text{yes}\\0&\sigma(x)=\text{no}.\end{cases}
\]
Identificando $\text{yes}\mapsto1$, $\text{no}\mapsto0$ si scrive compattamente
$u_1(x,\sigma)=x\,\sigma(x)$ e $u_2(x,\sigma)=(1-x)\,\sigma(x)$.
\end{defbox}

\begin{remark}[Strategie del secondo giocatore]
Si noti che $S_2$ è l'insieme di \emph{tutte} le funzioni da $[0,1]$ in $\{$yes,no$\}$: II
sceglie una regola di accettazione \emph{prima} di vedere l'offerta, come sempre nella forma
strategica di un gioco estensivo. È un insieme di strategie enorme, e questo spiega la
molteplicità di equilibri di Nash.
\end{remark}

\begin{proposition}[Esito perfetto nei sottogiochi]
Nell'unico equilibrio perfetto nei sottogiochi dell'ultimatum game il giocatore~II accetta
ogni offerta che gli lasci un pay-off $\ge0$ e il giocatore~I si tiene tutto: il pay-off di
equilibrio è $(1,0)$.
\end{proposition}

\textit{Dimostrazione (sketch).} Induzione a ritroso. Nel sottogioco che segue l'offerta $x$
il giocatore~II confronta $1-x$ con $0$: la razionalità sequenziale impone
$\sigma^*(x)=\text{yes}$ per ogni $x$ con $1-x\ge0$. Anticipando $\sigma^*$, il giocatore~I
risolve $\max_{x\in[0,1]}x$ e propone $x=1$. \qed

\begin{remark}
Si osservi che molte altre coppie $(x,\sigma)$ sono equilibri di Nash del gioco in forma
strategica: la regola $\sigma_t(x)=\text{yes}\iff x\le t$ sostiene l'esito $(t,1-t)$. Sono
però basate su minacce di rifiuto che non sono sequenzialmente razionali, e vengono eliminate
dalla perfezione nei sottogiochi.
\end{remark}

\section*{In sintesi per l'orale}
\addcontentsline{toc}{section}{In sintesi per l'orale (L15)}
\begin{itemize}
  \item Indebolendo la Pareto optimality a weak Pareto optimality e \emph{togliendo la
        simmetria}, gli assiomi (feasibility, rationality, weak Pareto, covarianza affine
        positiva, IIA) caratterizzano la famiglia delle \textbf{soluzioni di Nash
        generalizzate} $\argmax\prod_k(x_k-u^*_k)^{\alpha_k}$, con bargaining powers
        $\alpha_k\ge0$.
  \item La simmetria è esattamente ciò che impone $\alpha_1=\dots=\alpha_n$ e riporta alla
        soluzione di Nash classica.
  \item Divide the euro: $\psi=(\alpha_1/(\alpha_1+\alpha_2),\alpha_2/(\alpha_1+\alpha_2))$;
        l'euro si divide in proporzione al potere contrattuale.
  \item \textbf{Nash program}: giustificare le soluzioni cooperative come equilibri di giochi
        non cooperativi espliciti. È il ponte fra bargaining assiomatico e SPE.
  \item Nash demand game: $S_i=[u^*_i,m_i]$ con $m$ punto utopia; si ottiene la richiesta se
        il profilo è fattibile, altrimenti $u^*_i$. Ha un \emph{continuo} di equilibri di
        Nash lungo la frontiera di Pareto: l'equilibrio di Nash da solo non seleziona un esito.
  \item Ultimatum game: $S_1=[0,1]$, $S_2=\{\sigma:[0,1]\to\{\text{yes,no}\}\}$,
        $u_1=x\sigma(x)$, $u_2=(1-x)\sigma(x)$; molti equilibri di Nash (sostenuti da minacce
        non sequenzialmente razionali), \emph{unico} esito SPE $(1,0)$.
  \item \textbf{Collegamenti}: il passaggio ``molti NE $\to$ un solo SPE'' è lo stesso
        meccanismo dei giochi di Stackelberg (Cap.~\ref{ch:stackelberg-nash}); il punto
        utopia $m(U)$ è lo stesso oggetto usato per la soluzione di Kalai--Smorodinsky
        (Cap.~\ref{ch:bargaining-assiomatico}); i bargaining powers $\alpha_k$ hanno il ruolo
        che nei giochi TU hanno i pesi nel valore di $\Sh$ pesato.
\end{itemize}


% ============================================================
\chapter{Bargaining sequenziale: Rubinstein e Baron--Ferejohn (L16)}
\label{ch:bargaining-sequenziale}
% ============================================================

Si continua il Nash program con i modelli \emph{dinamici}: due giocatori che si scambiano
offerte, e poi il modello legislativo a $n$ giocatori. Lo strumento è sempre l'induzione a
ritroso / la razionalità sequenziale (SPE).

\section{Bargaining sequenziale in due stadi}

\begin{defbox}{Two stage game}
\emph{Divide the euro} con un solo round di offerta per giocatore:
\begin{enumerate}
  \item il giocatore~I propone $x\in[0,1]$ (la sua frazione);
  \item il giocatore~II accetta, oppure rifiuta e formula una \emph{controproposta finale}
        $y\in[0,1]$;
  \item il giocatore~I accetta o rifiuta la controproposta; se rifiuta entrambi prendono $0$.
\end{enumerate}
I giocatori hanno \textbf{fattori di sconto} $\delta_i\in\,]0,1[$: un pay-off $z$ ottenuto al
secondo stadio vale $\delta_i z$. Gli spazi di strategie sono
\[
  S_1=[0,1]\times\{\sigma_1:[0,1]\to\{\text{yes},\text{no}\}\},\qquad
  S_2=\{\sigma_2:[0,1]\to\{\text{yes}\}\cup[0,1]\},
\]
cioè II sceglie, per ogni offerta ricevuta, se accettare o quale controproposta fare.
\end{defbox}

\begin{thmbox}{SPE del gioco a due stadi}
Il gioco a due stadi ha un unico esito SPE: il giocatore~I propone
$x^*=1-\delta_2$ e il giocatore~II accetta immediatamente. I pay-off di equilibrio sono
\[
  (u_1,u_2)=(1-\delta_2,\ \delta_2).
\]
\end{thmbox}

\textit{Dimostrazione (sketch).} Induzione a ritroso.
\emph{Ultimo stadio}: di fronte alla controproposta finale $y$ (frazione per~II), il
giocatore~I confronta $\delta_1(1-y)$ con $0$ e quindi accetta ogni $y\le1$; con la solita
convenzione sull'indifferenza, II sceglie $y=1$ e ottiene $\delta_2\cdot1=\delta_2$
(scontato di un periodo), mentre I ottiene $0$.
\emph{Primo stadio}: la \emph{continuation value} di~II se rifiuta è dunque $\delta_2$;
quindi II accetta $x$ se e solo se $1-x\ge\delta_2$, cioè $x\le1-\delta_2$. Il giocatore~I
massimizza $x$ sotto questo vincolo e offre $x^*=1-\delta_2$. \qed

\begin{remark}
Per $\delta_2\to0$ si ritrova l'esito dell'ultimatum game, $(1,0)$: il secondo stadio non ha
più valore per~II. Si noti che $\delta_1$ non compare nella soluzione, perché l'ultimo a
proporre e'~II.
\end{remark}

\section{Bargaining sequenziale a orizzonte infinito}

È la situazione più complessa: entrambi i giocatori possono sempre formulare una
controproposta, finché non si raggiunge finalmente un accordo.

\begin{defbox}{Assunzioni preliminari}
\begin{itemize}
  \item il tempo è discreto, con passi temporali unitari $t\in\mathbb{N}$;
  \item i fattori di sconto non cambiano nel tempo;
  \item se il gioco prosegue all'infinito, entrambi i giocatori ottengono $0$.
\end{itemize}
Le utilità dei due giocatori sono date da
\[
  u_i(x,t)=\delta_i^{\,t}\,x ,
\]
e devono valere le seguenti proprietà:
\begin{itemize}
  \item $u_i(x,t)>u_i(z,t)\iff x>z$;
  \item $u_i(x,t)>u_i(x,s)\iff t<s$;
  \item $u_i(x,t)>u_i(z,t+1)\iff u_i(x,0)>u_i(z,1)$.
\end{itemize}
\end{defbox}

\begin{remark}
Nelle slide le tre proprietà sono enunciate a parole come: \emph{pie is desirable} (la torta
è desiderabile), \emph{time matters} (il tempo conta), \emph{stationarity over time}
(stazionarietà nel tempo), cui si aggiunge la \emph{increasing loss to delay}.
\end{remark}

\subsection{Strategie stazionarie e condizione di indifferenza}

Data la complessità del gioco si considerano soltanto \textbf{strategie stazionarie}: i due
giocatori giocano sempre la stessa strategia nel tempo. Si consideri un generico istante $t$
corrispondente a un sottogioco radicato nel giocatore~1.

In questi giochi l'induzione a ritroso esplicita \emph{non} può essere applicata, perché
l'albero del gioco ha altezza infinita. Tuttavia, per calcolare l'equilibrio di Nash perfetto
nei sottogiochi è necessario che il proponente corrente renda il rispondente
\textbf{indifferente} fra accettare e rifiutare l'offerta: il pay-off ottenuto accettando
l'offerta corrente deve essere uguale al pay-off atteso che il rispondente otterrebbe
rifiutando e formulando una controproposta.

\begin{thmbox}{Condizioni di indifferenza e offerte di equilibrio}
Per il giocatore~2 (che riceve l'offerta $x$ del giocatore~1):
\[
  \delta_2^{\,t}(1-x)=\delta_2^{\,t+1}y
  \iff
  1-x=\delta_2\,y .
\]
Simmetricamente, per il giocatore~1 (che riceve l'offerta $y$ del giocatore~2):
\[
  \delta_1^{\,t}(1-y)=\delta_1^{\,t+1}x
  \iff
  1-y=\delta_1\,x .
\]
Risolvendo il sistema di queste due equazioni si ottengono le strategie di offerta
\[
  x^*=\frac{1-\delta_2}{1-\delta_1\delta_2},
  \qquad
  y^*=\frac{1-\delta_1}{1-\delta_1\delta_2}.
\]
\end{thmbox}

\textit{Dimostrazione (sketch).} Sostituendo $y=\delta_1x+\ldots$, cioè ricavando
$y=(1-\delta_1x)$ dalla seconda equazione e inserendolo nella prima, si ha
$1-x=\delta_2(1-\delta_1x)$, da cui $x(1-\delta_1\delta_2)=1-\delta_2$ e quindi
$x^*=(1-\delta_2)/(1-\delta_1\delta_2)$; in modo del tutto analogo si ricava $y^*$. \qed

\subsection{L'equilibrio perfetto nei sottogiochi}

\begin{defbox}{Funzioni di accettazione}
\[
  \sigma_1^*(y)=
  \begin{cases}
    \text{yes} & \text{se } y\le 1-y^*\\[2pt]
    x^* & \text{se } y> 1-y^*
  \end{cases}
  \qquad
  \sigma_2^*(x)=
  \begin{cases}
    \text{yes} & \text{se } x\le 1-x^*\\[2pt]
    y^* & \text{se } x> 1-x^*
  \end{cases}
\]
Cioè: se l'offerta ricevuta è accettabile si risponde \emph{yes}, altrimenti si rilancia con
la propria offerta stazionaria.
\end{defbox}

\begin{thmbox}{Teorema di Rubinstein (1982)}
L'unico equilibrio di Nash perfetto nei sottogiochi è il profilo di strategie
\[
  \bigl((x^*,\sigma_1^*),\,(y^*,\sigma_2^*)\bigr).
\]
Come nei casi precedenti, se entrambi i giocatori si comportano razionalmente il gioco
termina dopo la prima offerta, con pay-off
\[
  \frac{1-\delta_2}{1-\delta_1\delta_2}
  \qquad\text{e}\qquad
  1-\frac{1-\delta_2}{1-\delta_1\delta_2}=\frac{\delta_2(1-\delta_1)}{1-\delta_1\delta_2}.
\]
Se invece l'offerta viene accettata al tempo $t$, i pay-off sono
$\bigl((x^*)^t,\,1-(x^*)^t\bigr)$.
\end{thmbox}

\begin{thmbox}{Collegamento con il Nash program}
Al tendere di $t\to0$ i pay-off convergono a
\[
  \Bigl(\frac{\rho_1}{\rho_1+\rho_2},\ \frac{\rho_2}{\rho_1+\rho_2}\Bigr),
\]
che corrisponde alla \textbf{generalized Nash bargaining solution} con bargaining powers
$\rho_i$. In particolare
\[
  \rho_i=\log\frac{1}{\delta_i},
\]
e i $\rho_i$ sono detti \textbf{instantaneous rates of interest}.
\end{thmbox}

\begin{remark}
Questo è il punto in cui il \emph{Nash program} si chiude: una procedura di contrattazione
puramente non cooperativa (offerte alternate, SPE) restituisce esattamente la soluzione di
Nash generalizzata introdotta per via assiomatica in L15, con bargaining powers determinati
dall'impazienza dei giocatori.
\end{remark}

\section{Il modello di Baron--Ferejohn}

Quest'ultimo gioco di bargaining sequenziale differisce dai precedenti perché coinvolge
$n$ giocatori.

\begin{defbox}{Baron--Ferejohn model}
A ogni round un \textbf{proponente} viene selezionato casualmente fra gli $n$ giocatori e deve
ottenere l'approvazione di una maggioranza fissata di $k$ giocatori, formulando offerte
opportune. Tutti i giocatori condividono lo stesso fattore di sconto $\delta$. Come nei giochi
a orizzonte infinito, il gioco procede finché un'offerta non viene accettata; se un'offerta
viene rifiutata, un nuovo proponente viene scelto casualmente e l'istante temporale aumenta.
Come nei casi precedenti, il gioco si analizza tramite \textbf{strategie stazionarie}.
\end{defbox}

Nello specifico, il proponente offre una quota $s\le1$ a ciascuno dei $k$ rispondenti.
L'offerta ottimale $s$ è determinata rendendo i rispondenti indifferenti fra accettare la
proposta corrente e rifiutarla. Il proponente impone quindi l'equazione

\begin{thmbox}{Condizione di indifferenza e reservation share}
\[
  s=\delta\Bigl[\frac{1}{n}(1-ks)+\frac{k}{n}\,s\Bigr]=\frac{\delta}{n}.
\]
Il valore $s^*=\dfrac{\delta}{n}$ è detto \textbf{reservation share}.
\end{thmbox}

Questa espressione rappresenta la condizione di indifferenza per un generico rispondente. Il
pay-off immediato che deriva dall'accettare l'offerta, $s$, deve eguagliare il pay-off atteso
scontato che si ottiene aspettando fino al periodo successivo. Nel round seguente il giocatore
può essere selezionato come proponente con probabilità $1/n$, nel qual caso riceverebbe un
pay-off pari a $1-ks$, oppure essere uno dei $k$ rispondenti con probabilità $k/n$, ricevendo
quindi $s$. La media pesata di questi pay-off possibili dà il valore di continuazione atteso,
che viene moltiplicato per $\delta$ a causa del tempo trascorso.

\begin{remark}
Si noti che nell'espressione fra parentesi i termini in $ks$ si cancellano:
$\frac1n(1-ks)+\frac kn s=\frac1n-\frac{ks}{n}+\frac{ks}{n}=\frac1n$, da cui immediatamente
$s=\delta/n$. Il proponente si tiene quindi $1-ks^*$.
\end{remark}
\section*{In sintesi per l'orale}
\addcontentsline{toc}{section}{In sintesi per l'orale (L16)}
\begin{itemize}
  \item Two stage game: strategie
        $S_1=[0,1]\times\{\sigma_1:[0,1]\to\{\text{yes,no}\}\}$,
        $S_2=\{\sigma_2:[0,1]\to\{\text{yes}\}\cup[0,1]\}$, fattori di sconto
        $\delta_i\in\,]0,1[$. Per induzione a ritroso l'esito SPE è $(1-\delta_2,\delta_2)$.
  \item Orizzonte infinito -- assunzioni: tempo discreto con passi unitari $t\in\mathbb{N}$;
        fattori di sconto costanti nel tempo; se il gioco non finisce mai entrambi prendono $0$.
        Utilità: $u_i(x,t)=\delta_i^{\,t}x$.
  \item Le tre proprietà richieste: $u_i(x,t)>u_i(z,t)\iff x>z$;
        $u_i(x,t)>u_i(x,s)\iff t<s$; $u_i(x,t)>u_i(z,t+1)\iff u_i(x,0)>u_i(z,1)$.
  \item Si considerano solo \textbf{strategie stazionarie}: l'induzione a ritroso esplicita non
        si applica (albero di altezza infinita), quindi si impone che il proponente renda il
        rispondente \emph{indifferente} fra accettare e rifiutare.
  \item Le due equazioni chiave: $\delta_2^{\,t}(1-x)=\delta_2^{\,t+1}y\iff 1-x=\delta_2y$ e
        $\delta_1^{\,t}(1-y)=\delta_1^{\,t+1}x\iff 1-y=\delta_1x$, da cui
        $x^*=\frac{1-\delta_2}{1-\delta_1\delta_2}$ e $y^*=\frac{1-\delta_1}{1-\delta_1\delta_2}$.
  \item Funzioni di accettazione $\sigma_1^*$, $\sigma_2^*$ (yes se l'offerta ricevuta è
        accettabile, altrimenti si rilancia con la propria offerta stazionaria); l'unico SPE è
        $((x^*,\sigma_1^*),(y^*,\sigma_2^*))$, il gioco finisce dopo la prima offerta con
        pay-off $\frac{1-\delta_2}{1-\delta_1\delta_2}$ e $\frac{\delta_2(1-\delta_1)}{1-\delta_1\delta_2}$
        (accettando al tempo $t$: $((x^*)^t,1-(x^*)^t)$).
  \item \textbf{Nash program}: al limite i pay-off convergono a
        $\bigl(\frac{\rho_1}{\rho_1+\rho_2},\frac{\rho_2}{\rho_1+\rho_2}\bigr)$, cioè alla
        \emph{generalized Nash bargaining solution} con bargaining powers
        $\rho_i=\log\frac{1}{\delta_i}$ (instantaneous rates of interest). È il collegamento
        diretto con L15.
  \item \textbf{Baron--Ferejohn}: $n$ giocatori, proponente estratto a caso ogni round,
        approvazione di una maggioranza fissata di $k$ giocatori, sconto comune $\delta$,
        orizzonte infinito, strategie stazionarie. Il proponente offre $s$ a ciascuno dei $k$
        rispondenti.
  \item Condizione di indifferenza: $s=\delta\bigl[\frac1n(1-ks)+\frac kn s\bigr]=\frac{\delta}{n}$
        (con probabilità $1/n$ si è proponente e si prende $1-ks$, con probabilità $k/n$ si
        è rispondenti e si prende $s$; il tutto scontato di $\delta$). Il valore
        $s^*=\delta/n$ è la \textbf{reservation share}.
  \item \textbf{Collegamenti}: ultimatum $\to$ two stage $\to$ orizzonte infinito è la catena
        del Nash program iniziato in L15; il two stage è il caso limite in cui il gioco si
        arresta dopo un round; Baron--Ferejohn estende il bargaining sequenziale da 2 a $n$
        giocatori con maggioranza qualificata.
\end{itemize}

% ============================================================
%  Capitoli L17 -- L19
%  (da includere nella dispensa: nessun preambolo, nessun document)
% ============================================================

\chapter{Exchange Markets e One-sided Matching (L17)}
\label{ch:exchange-ttc}
% ============================================================

\section{Economie di scambio e mercati di matching}

Nelle \textbf{exchange economies} (economie di scambio) ogni agente possiede una \emph{dotazione iniziale} (initial endowment) di beni; i beni vengono scambiati fra gli agenti \textbf{senza pagamenti monetari}. Il criterio di valutazione non è quindi un'utilità cardinale trasferibile, ma la sola \emph{preferenza ordinale} di ciascun agente sui panieri che può ottenere.

Il caso limite, e il più studiato dal punto di vista algoritmico, è quello in cui ogni bene è \textbf{unico e indivisibile} e ogni dotazione è \textbf{unitaria} (un bene per agente): l'economia di scambio degenera in un problema di \textbf{matching}.

\begin{defbox}{Matching markets}
Si devono accoppiare i membri di (al più) due gruppi. L'accoppiamento può essere \emph{one-to-one}, \emph{many-to-one} oppure \emph{many-to-many}. Si distinguono:
\begin{itemize}
  \item \textbf{one-sided matching}: solo uno dei due lati (gruppo) esprime preferenze sui possibili matching; l'altro lato è costituito da oggetti passivi (case, posti, organi, \dots);
  \item \textbf{two-sided matching}: entrambi i lati esprimono preferenze (Cap.~\ref{ch:two-sided}).
\end{itemize}
Un \textbf{mechanism} è una procedura (un algoritmo) che, dati i dati del problema e le preferenze dichiarate, restituisce un matching.
\end{defbox}

\section{House allocation (Shapley--Scarf 1974)}

\begin{defbox}{Problema di house allocation}
Dati: un insieme di agenti $N$ e un insieme di case $H$ con $|N|=|H|=n$.

Ogni agente $i\in N$ dichiara \textbf{preferenze strette} su $H$ tramite una relazione binaria $\succ_i$ che soddisfa
\begin{itemize}
  \item \emph{completezza}: per ogni $h_1,h_2\in H$ vale $h_1\succ_i h_2$ oppure $h_2\succ_i h_1$;
  \item \emph{irriflessività}: $h\not\succ_i h$;
  \item \emph{transitività}: $h_1\succ_i h_2$ e $h_2\succ_i h_3$ implicano $h_1\succ_i h_3$.
\end{itemize}
(L'\emph{asimmetria} segue: $h_1\succ_i h_2\iff h_2\not\succ_i h_1$.)

Un \textbf{matching} è una qualunque funzione iniettiva (quindi biiettiva, essendo $|N|=|H|$) $f:N\to H$, cui corrisponde l'insieme di coppie
\[ P=\bigl\{(i,f(i)) : i\in N\bigr\}. \]
\end{defbox}

Il dato aggiuntivo che rende il problema un'\emph{economia di scambio} è un \textbf{matching iniziale} $f_0$: $f_0(i)$ è la casa \emph{posseduta} da $i$ prima degli scambi. L'idea di fondo è allora:

\begin{center}\itshape
sfruttare le permutazioni di allocazione mutuamente vantaggiose.
\end{center}

\section{Il digrafo degli scambi desiderati}

\begin{defbox}{Digrafo degli scambi unilateralmente top-desiderati}
Dato un matching $f$ e un sottoinsieme $S\subseteq N$, si definisce $h_S(i)\in S$ come l'agente che, \emph{nella restrizione $f|_S$}, detiene l'oggetto preferito da $i$:
\[ j=h_S(i)\iff f(j)\succ_i f(k)\quad\forall k\in S,\ k\ne j. \]
È possibile che $h_S(i)=i$ (l'agente possiede già il meglio disponibile in $S$). Si pone
\[ A_f(S)=\bigl\{(i,h_S(i)) : i\in S\bigr\}, \]
e $(S,A_f(S))$ è il \textbf{digrafo} degli scambi unilateralmente top-desiderati dentro $S$.
\end{defbox}

\begin{remark}
Poiché le preferenze sono strette, $h_S$ è ben definita e ogni nodo ha \emph{esattamente un} arco uscente. In un digrafo finito con out-degree $1$ esiste sempre almeno un ciclo orientato. Un ciclo orientato corrisponde a una \textbf{permutazione mutuamente vantaggiosa} interna a $S$; un cappio (self-loop) indica che l'agente ha già la massima soddisfazione ottenibile in $f|_S$.
\end{remark}

\section{Top Trading Cycle (TTC)}

\begin{defbox}{Algoritmo Top Trading Cycle (Gale)}
\begin{enumerate}
  \item $S_0=N$, $f_0$ matching iniziale, $k=0$;
  \item se $S_k=\emptyset$ allora STOP;
  \item costruisci il digrafo $G_k=(S_k,A_{f_0}(S_k))$;
  \item trova un ciclo orientato $C_k$ in $G_k$;
  \item poni $f(i)=f_0(h_{S_k}(i))$ per ogni $i\in C_k$;
  \item $S_{k+1}=S_k\setminus C_k$, $k=k+1$ e torna al passo 2.
\end{enumerate}
\end{defbox}

\begin{remark}
Osservazioni sull'algoritmo:
\begin{itemize}
  \item l'algoritmo termina in al più $n$ iterazioni e restituisce un matching $f$ (a ogni iterazione almeno un agente viene rimosso);
  \item al passo 4 si possono individuare \emph{tutti} i cicli e rimuoverli in un solo aggiornamento dei passi 5--6: i cicli di un digrafo con out-degree $1$ sono infatti a due a due disgiunti;
  \item il matching finale $f$ è \textbf{fortemente sensibile} al matching iniziale $f_0$: cambiando le dotazioni cambia l'esito, pur restando le preferenze le stesse.
\end{itemize}
\end{remark}

\begin{exbox}{Esecuzione di TTC su 4 agenti}
$N=\{1,2,3,4\}$, $H=\{1,2,3,4\}$, dotazioni $f_0(i)=i$ (l'agente $i$ possiede la casa $i$). Preferenze (dalla più alla meno preferita):
\begin{center}
\begin{tabular}{c|cccc}
\toprule
agente & \multicolumn{4}{c}{ordine di preferenza sulle case}\\
\midrule
$1$ & $3$ & $4$ & $1$ & $2$\\
$2$ & $1$ & $3$ & $4$ & $2$\\
$3$ & $1$ & $2$ & $4$ & $3$\\
$4$ & $2$ & $1$ & $3$ & $4$\\
\bottomrule
\end{tabular}
\end{center}

\emph{Iterazione $k=0$.} $S_0=\{1,2,3,4\}$. Ciascuno punta al \emph{proprietario} della casa che preferisce in assoluto:
$1\to 3$ (vuole la casa $3$, posseduta da $3$), $2\to 1$, $3\to 1$, $4\to 2$. L'unico ciclo orientato è $1\to 3\to 1$, cioè $C_0=\{1,3\}$ (gli archi $2\to1$ e $4\to2$ non chiudono cicli). Si assegna
\[ f(1)=f_0(h_{S_0}(1))=f_0(3)=3,\qquad f(3)=f_0(1)=1. \]
$S_1=\{2,4\}$: entrambi ottengono la loro prima scelta assoluta.

\emph{Iterazione $k=1$.} Restano gli agenti $\{2,4\}$ e quindi le case $\{2,4\}$. Fra queste, $2$ preferisce la casa $4$ (le case $1$ e $3$ non ci sono più), posseduta da $4$: $2\to 4$. L'agente $4$ preferisce la casa $2$, posseduta da $2$: $4\to 2$. Il ciclo è $C_1=\{2,4\}$ e si scambiano:
\[ f(2)=f_0(4)=4,\qquad f(4)=f_0(2)=2. \]
$S_2=\emptyset$, STOP.

Matching finale: $f=(1\mapsto 3,\;2\mapsto 4,\;3\mapsto 1,\;4\mapsto 2)$. Si verifica l'\emph{individual rationality}: $1$ e $3$ ottengono la prima scelta; $2$ passa dalla casa $2$ (ultima in lista) alla $4$ (terza); $4$ passa dalla casa $4$ (ultima) alla $2$ (prima). Nessuno peggiora rispetto a $f_0$.
\end{exbox}

\begin{exbox}{Un ciclo di lunghezza $3$}
$N=H=\{1,2,3\}$, $f_0(i)=i$, con $2\succ_1 3\succ_1 1$, $3\succ_2 1\succ_2 2$, $1\succ_3 2\succ_3 3$. Il digrafo iniziale è $1\to 2\to 3\to 1$: un unico ciclo che coinvolge tutti. TTC termina in una sola iterazione con la permutazione ciclica $f=(1\mapsto 2,\;2\mapsto 3,\;3\mapsto 1)$: ogni agente ottiene la sua prima scelta.

Se invece tutti preferissero la casa $1$ ($1\succ_i h$ per ogni $i$ e $h\neq 1$), il digrafo avrebbe $1\to 1$ e $2\to 1$, $3\to 1$: il ciclo è il cappio su $1$, che tiene la propria casa; il gioco prosegue su $\{2,3\}$.
\end{exbox}

\section{Proprietà di TTC: stabilità e strategy-proofness}

\begin{thmbox}{Stabilità di TTC (Roth--Postlewaite 1977)}
Il matching $f$ prodotto dall'algoritmo TTC è
\begin{itemize}
  \item \textbf{individually rational}: per ogni $i\in N$ vale $f(i)\succ_i f_0(i)$ oppure $f(i)=f_0(i)$;
  \item \textbf{stabile}: non esistono $S\subseteq N$ e una permutazione $\pi:S\to S$ tali che ogni $i\in S$ soddisfi
  \[ \pi(i)\ne i\quad\text{e}\quad f(\pi(i))\succ_i f(i), \]
  cioè nessun gruppo di agenti può riorganizzare le proprie allocazioni migliorando tutti.
\end{itemize}
Inoltre \textbf{non esiste alcun altro meccanismo} individually rational e stabile.
\end{thmbox}

\textit{Dimostrazione (sketch).} \emph{Individual rationality}: un agente $i$ rimosso all'iterazione $k$ riceve $f_0(h_{S_k}(i))$, cioè l'oggetto che preferisce fra quelli posseduti dagli agenti ancora in $S_k$; poiché $i\in S_k$, la propria dotazione $f_0(i)$ è fra i candidati, dunque $f(i)\succ_i f_0(i)$ oppure $f(i)=f_0(i)$.

\emph{Stabilità}: si ragiona sull'ordine di rimozione dei cicli. Gli agenti del primo ciclo $C_0$ ricevono la loro \emph{prima scelta assoluta} su tutto $H$, quindi nessuna permutazione può migliorarli e una coalizione bloccante non può contenerli (deve migliorare \emph{tutti} i suoi membri). Rimossi $C_0$ e i relativi oggetti, gli agenti di $C_1$ ricevono la prima scelta fra gli oggetti rimasti, e così via: la coalizione bloccante deve essere vuota. \qed

\begin{remark}
La condizione di stabilità è una condizione di tipo \emph{blocco di coalizione}: la permutazione $\pi$ ridistribuisce a $S$ \emph{soltanto} gli oggetti già assegnati a $S$, quindi la coalizione deve poter migliorare tutti i suoi membri con le sole risorse interne. È la stessa logica con cui, nei giochi coalizionali (L19), una coalizione valuta se conviene staccarsi dalla grande coalizione.
\end{remark}

\begin{thmbox}{Truthfulness pays back (Roth 1982)}
Il meccanismo TTC è \textbf{strategy-proof}: per ogni agente $i\in N$, detto $f$ il matching TTC sotto le preferenze $(\succ_i,\succ_{-i})$ e $f'$ quello sotto $(\succ_i',\succ_{-i})$, vale
\[ f'(i)\ne h\qquad\text{per ogni } h\succ_i f(i). \]
Nessun agente può cioè ottenere un'allocazione migliore dichiarando preferenze diverse da quelle reali.
\end{thmbox}

(Le slide riportano il solo enunciato: nessun agente può ottenere un'allocazione migliore falsificando le proprie preferenze.)

\section{Varianti e altri contesti}

\begin{itemize}
  \item \textbf{nessun matching iniziale disponibile}: se ne genera uno casualmente per far girare TTC (\emph{random endowment}), oppure si usa la \textbf{serial dictatorship} quando gli agenti hanno priorità;
  \item \textbf{matching iniziale disponibile solo per alcuni agenti}: esistono varianti di TTC (agenti con e senza dotazione, ``houses'' libere);
  \item \textbf{preferenze non strette} (indifferenze): si rompono i pareggi arbitrariamente (tie-breaking) oppure si gestiscono internamente al meccanismo;
  \item \textbf{preferenze mancanti}: alternative inaccettabili, che l'agente rifiuta comunque;
  \item $|N|\ne|H|$: occorre gestire la mancanza di accoppiamenti possibili;
  \item estensioni \textbf{many-to-one} e \textbf{many-to-many}.
\end{itemize}

\section*{In sintesi per l'orale}
\addcontentsline{toc}{section}{In sintesi per l'orale (L17)}
\begin{itemize}
  \item \textbf{Exchange economy}: dotazioni iniziali, scambi \emph{senza denaro}, preferenze ordinali strette. Caso unitario e indivisibile $\Rightarrow$ house allocation (Shapley--Scarf).
  \item \textbf{One-sided vs two-sided}: nel one-sided solo gli agenti hanno preferenze; nel two-sided (L18) entrambi i lati.
  \item \textbf{Digrafo $A_f(S)$}: out-degree $1$ $\Rightarrow$ esiste sempre un ciclo; i cicli sono disgiunti $\Rightarrow$ si possono rimuovere tutti insieme.
  \item \textbf{TTC}: al più $n$ iterazioni; l'esito dipende in modo essenziale da $f_0$.
  \item \textbf{Roth--Postlewaite}: TTC è individually rational e stabile, ed è l'\emph{unico} meccanismo con queste due proprietà.
  \item \textbf{Roth 1982}: TTC è strategy-proof. Contrasto forte con il two-sided matching, dove la strategy-proofness è impossibile (L18) e con le social choice functions (Gibbard--Satterthwaite).
  \item \textbf{Collegamento}: la stabilità è una condizione di blocco coalizionale, la stessa logica che in L19 governa la formazione delle coalizioni nei giochi TU.
  \item \textbf{Sketch chiave da saper ripetere}: individual rationality (la propria dotazione è sempre fra le opzioni) e stabilità per induzione sui cicli (il primo ciclo ottiene la prima scelta assoluta).
\end{itemize}

% ============================================================
\chapter{Two-sided Matching e Social Choice (L18)}
\label{ch:two-sided}
% ============================================================

\section{Il problema del matrimonio (Gale--Shapley 1962)}

\begin{defbox}{Marriage problem}
Dati: due insiemi $M$ (uomini) e $W$ (donne) con $|M|=|W|=n$; per ogni $m\in M$ una preferenza stretta $\succ_m$ su $W$, per ogni $w\in W$ una preferenza stretta $\succ_w$ su $M$.

Un \textbf{matching} è una funzione iniettiva $f:W\to M$, cui corrisponde
\[ P=\bigl\{(w,f(w)) : w\in W\bigr\}=\bigl\{(f^{-1}(m),m) : m\in M\bigr\}, \]
dove $f^{-1}:M\to W$ è la funzione inversa.
\end{defbox}

\begin{definition}[Partner più e meno preferito]
Per $\hat M\subseteq M$ si pone
\begin{itemize}
  \item[(i)] $\bar m=\max_w \hat M$ se $\bar m\succ_w m$ per ogni $m\in\hat M$, $m \neq \bar m$;
  \item[(ii)] $\bar m=\min_w \hat M$ se $m\succ_w \bar m$ per ogni $m\in\hat M$, $m \neq \bar m$;
\end{itemize}
e analogamente $\max_m \hat W$, $\min_m \hat W$ per $\hat W\subseteq W$.
\end{definition}

\section{Stabilità}

\begin{defbox}{Matching stabile}
\begin{itemize}
  \item[(i)] Una coppia $(w,m)\in W\times M$ ha un'\textbf{obiezione} (blocking pair) al matching $f$ quando
  \[ m\succ_w f(w)\quad\text{e}\quad w\succ_m f^{-1}(m); \]
  \item[(ii)] $f$ è \textbf{stabile} se nessuna coppia donna-uomo ha obiezioni.
\end{itemize}
Equivalentemente, $f$ è stabile se e solo se ogni coppia $(w,m)$ soddisfa
\[ m\succ_w f(w)\ \Longrightarrow\ f^{-1}(m)\succ_m w, \]
o, in modo del tutto equivalente,
\[ w\succ_m f^{-1}(m)\ \Longrightarrow\ f(w)\succ_w m. \]
\end{defbox}

\begin{center}\itshape
obiezione $\equiv$ i due si preferiscono reciprocamente ai partner assegnati.
\end{center}

\begin{remark}[Stabilità come non-dominanza]
Come nel one-sided matching, la stabilità si può leggere come assenza di blocco coalizionale (appunti del corso). Dati due matching $f,\hat f$, si dice che \textbf{$f$ è dominato da $\hat f$} se esiste $S\subseteq M\cup W$, $S\ne\emptyset$, tale che
\begin{itemize}
  \item[(i)] $\hat f(S\cap W)=S\cap M=f(S\cap W)$ (la coalizione si ridistribuisce \emph{internamente});
  \item[(ii)] $\hat f(w)\succ_w f(w)$ per ogni $w\in S\cap W$;
  \item[(iii)] $\hat f^{-1}(m)\succ_m f^{-1}(m)$ per ogni $m\in S\cap M$.
\end{itemize}
Vale allora: \emph{$f$ è stabile se e solo se $f$ non è dominato}. In particolare basta considerare coalizioni $S$ di due elementi (una coppia), il che riconduce la non-dominanza all'assenza di blocking pairs.
\end{remark}

\begin{thmbox}{Esistenza di matching stabili (Gale--Shapley 1962)}
Per il marriage problem esiste sempre almeno un matching stabile.
\end{thmbox}

La dimostrazione è \emph{costruttiva}: l'algoritmo di deferred acceptance produce un matching stabile.

\section{Il meccanismo di deferred acceptance}

\begin{defbox}{Deferred acceptance: men's courtship algorithm (Gale--Shapley 1962)}
Idea: ogni uomo propone alla sua \emph{top-available} donna (la più preferita fra quelle che non lo hanno ancora rifiutato); ogni donna trattiene provvisoriamente il proponente che preferisce (fra i proponenti \emph{compreso} l'eventuale match corrente) e rifiuta gli altri; ogni uomo rifiutato passa alla donna successiva nella sua lista.

Formalmente, con $p_m=(w_1,\dots,w_n)$ lista di preferenza di $m$ (cioè $p_m(i)\succ_m p_m(j)\iff i<j$):
\begin{enumerate}
  \item $r_m=1$ per ogni $m\in M$, $k=1$;
  \item $L_k(w)=\{m : p_m(r_m)=w\}$ per ogni $w\in W$;
  \item se $|L_k(w)|=1$ per ogni $w\in W$ allora STOP, con $f(w)=L_k(w)$;
  \item $\bar m=\max_w L_k(w)$, $f(w)=\bar m$;
  \item $r_m=r_m+1$ per ogni $m\in L_k(w)$, $m\ne\bar m$;\quad (4--5 ripetuti per ogni $w\in W$)
  \item $k=k+1$ e torna al passo 2.
\end{enumerate}
Il \textbf{women's courtship algorithm} si ottiene scambiando i ruoli di uomini e donne.
\end{defbox}

\begin{remark}
\begin{itemize}
  \item L'algoritmo termina dopo al più $n(n-1)+1$ iterazioni: ogni iterazione non terminale comporta almeno un rifiuto, e ogni uomo può essere rifiutato al più $n-1$ volte.
  \item Nel corso delle iterazioni gli uomini \textbf{scendono} lungo la propria lista di preferenze (proposte via via a donne meno preferite), mentre le donne \textbf{salgono} (il match provvisorio può solo migliorare). Questo è il motivo strutturale della asimmetria fra i due lati.
\end{itemize}
\end{remark}

\begin{exbox}{Esecuzione completa di Gale--Shapley ($n=3$)}
$M=\{m_1,m_2,m_3\}$, $W=\{w_1,w_2,w_3\}$, con preferenze
\begin{center}
\begin{tabular}{c|ccc@{\qquad}c|ccc}
\toprule
uomo & \multicolumn{3}{c}{lista} & donna & \multicolumn{3}{c}{lista}\\
\midrule
$m_1$ & $w_1$ & $w_2$ & $w_3$ & $w_1$ & $m_3$ & $m_1$ & $m_2$\\
$m_2$ & $w_1$ & $w_3$ & $w_2$ & $w_2$ & $m_1$ & $m_3$ & $m_2$\\
$m_3$ & $w_2$ & $w_1$ & $w_3$ & $w_3$ & $m_2$ & $m_1$ & $m_3$\\
\bottomrule
\end{tabular}
\end{center}

\emph{Men's courtship.} $k=1$: $m_1\to w_1$, $m_2\to w_1$, $m_3\to w_2$. Dunque $L_1(w_1)=\{m_1,m_2\}$, $L_1(w_2)=\{m_3\}$, $L_1(w_3)=\emptyset$: non tutti singoletti. $\max_{w_1}\{m_1,m_2\}=m_1$ (poiché $m_1\succ_{w_1}m_2$): $w_1$ trattiene $m_1$ e rifiuta $m_2$, quindi $r_{m_2}=2$.

$k=2$: $m_2\to w_3$. Ora $L_2(w_1)=\{m_1\}$, $L_2(w_2)=\{m_3\}$, $L_2(w_3)=\{m_2\}$: tutti singoletti, STOP.

Matching \emph{men-optimal}: $f_M=\{(w_1,m_1),(w_2,m_3),(w_3,m_2)\}$.

\emph{Verifica di stabilità.} $m_1$ ha $w_1$, la sua prima scelta. $m_2$ ha $w_3$ e preferirebbe $w_1$: ma $w_1$ ha $m_1$ e $m_1\succ_{w_1}m_2$, nessuna obiezione. $m_3$ ha $w_2$, la sua prima scelta. Dunque $f_M$ è stabile.

\emph{Women's courtship.} $k=1$: $w_1\to m_3$, $w_2\to m_1$, $w_3\to m_2$: le tre proposte vanno a uomini distinti, STOP immediato. Matching \emph{women-optimal}:
\[ f_W=\{(w_1,m_3),(w_2,m_1),(w_3,m_2)\}. \]
Anche $f_W$ è stabile (ogni donna ottiene la prima scelta, quindi nessuna donna può far parte di una blocking pair).

\emph{Confronto (Knuth).} $f_M\neq f_W$: $w_1$ e $w_2$ hanno partner diversi. Gli \textbf{uomini} preferiscono $f_M$ ($m_1$: $w_1\succ_{m_1}w_2$; $m_3$: $w_2\succ_{m_3}w_1$), le \textbf{donne} preferiscono $f_W$ ($w_1$: $m_3\succ_{w_1}m_1$; $w_2$: $m_1\succ_{w_2}m_3$). Il conflitto è esattamente quello descritto dal teorema di competitività intrinseca.
\end{exbox}

\section{Ottimalità per i proponenti}

\begin{thmbox}{Proponent optimality}
Il matching $f$ prodotto dal men's courtship algorithm è
\begin{itemize}
  \item \textbf{men-optimal}: ogni altro matching stabile $\hat f$ fornisce a ogni $m\in M$ un partner non migliore, cioè $\hat f^{-1}(m)\preceq_m f^{-1}(m)$;
  \item \textbf{women-disappointing}: ogni altro matching stabile $\hat f$ fornisce a ogni $w\in W$ un partner non peggiore, cioè $\hat f(w)\not\prec_w f(w)$.
\end{itemize}
In particolare il men's courtship mechanism è \textbf{weak Pareto optimal} per gli uomini.
\end{thmbox}

\textit{Idea (sketch).} Poiché gli uomini scendono lungo le proprie liste e le donne salgono, il men's courtship si ferma il prima possibile per gli uomini: si mostra che nessun uomo viene mai rifiutato da una donna con cui potrebbe essere accoppiato in un qualche matching stabile (altrimenti si costruisce una coppia con obiezione in quel matching). La parte \emph{women-disappointing} è l'altra faccia dello stesso enunciato, per il teorema di Knuth qui sotto. \qed

\begin{thmbox}{Competitività intrinseca dei matching stabili (Knuth 1976)}
Siano $f$ e $\hat f$ due matching stabili. Allora
\[ \hat f^{-1}(m)\preceq_m f^{-1}(m)\ \ \forall m\in M \iff \hat f(w)\not\prec_w f(w)\ \ \forall w\in W. \]
\end{thmbox}

\textit{Dimostrazione (sketch).} Se tutti gli uomini stanno debolmente peggio in $\hat f$ che in $f$, si consideri una donna $w$ con $\hat f(w)\prec_w f(w)$: posto $m=f(w)$, si ha $w\succ_m \hat f^{-1}(m)$ (perché $m$ sta peggio in $\hat f$) e $m\succ_w \hat f(w)$: la coppia $(w,m)$ obietta a $\hat f$, contraddicendone la stabilità. L'implicazione inversa è simmetrica. \qed

\begin{remark}
Il teorema esprime l'\textbf{opposizione strutturale} fra i due lati: sull'insieme dei matching stabili, ciò che è meglio per tutti gli uomini è (debolmente) peggio per tutte le donne. Ne segue che il matching del men's courtship è il migliore per gli uomini e il peggiore per le donne fra tutti gli stabili.
\end{remark}

\section{Strategy-proofness e manipolazione}

\begin{thmbox}{Impossibilità di strategy-proofness (Roth 1982)}
Nessun meccanismo che produca sempre matching stabili può essere strategy-proof.
\end{thmbox}

\begin{thmbox}{Coalitional strategy-proofness (Dubins--Freedman 1981)}
Il men's courtship mechanism è \textbf{coalitionally strategy-proof per gli uomini}: nessun sottogruppo $S\subseteq M$ può ottenere
\[ (f')^{-1}(m)\succ_m f^{-1}(m)\qquad\text{per ogni } m\in S, \]
dove $f$ è il matching del men's courtship sotto le preferenze vere $(\succ_S,\succ_{-S})$ e $f'$ un qualunque matching stabile sotto $(\succ_S',\succ_{-S})$.
\end{thmbox}

(Le slide riportano il solo enunciato. Si noti che è la controparte positiva del risultato di impossibilità: la manipolazione resta possibile, ma solo sul lato non proponente.)

\begin{thmbox}{Manipolazione della DA (Gale--Sotomayor 1985)}
Se esistono almeno due matching stabili, allora può esistere una donna $w\in W$ e una preferenza (falsa) $\succ_w'$ tale che
\[ f'(w)\succ_w f(w), \]
dove $f$ è il matching del men's courtship sotto $(\succ_w,\succ_{-w})$ e $f'$ quello sotto $(\succ_w',\succ_{-w})$.
\end{thmbox}

\begin{exbox}{Manipolazione da parte di una donna}
Si riprenda l'istanza $n=3$ dell'esempio precedente. Con le preferenze vere, il men's courtship dà a $w_1$ il partner $m_1$, ma $w_1$ preferirebbe $m_3$. Supponiamo che $w_1$ dichiari la falsa lista $\succ_{w_1}'$: $m_3\succ' m_2\succ' m_1$ (cioè finga di preferire $m_2$ a $m_1$).

$k=1$: $m_1\to w_1$, $m_2\to w_1$, $m_3\to w_2$. Ora $w_1$ trattiene $m_2$ e \emph{rifiuta} $m_1$.

$k=2$: $m_1\to w_2$. $L_2(w_2)=\{m_3,m_1\}$ e $m_1\succ_{w_2}m_3$: $w_2$ trattiene $m_1$ e rifiuta $m_3$.

$k=3$: $m_3\to w_1$. $L_3(w_1)=\{m_2,m_3\}$ e, secondo la lista dichiarata, $m_3$ vince: $w_1$ trattiene $m_3$ e rifiuta $m_2$.

$k=4$: $m_2\to w_3$; tutti singoletti, STOP. Il matching prodotto è
\[ f'=\{(w_1,m_3),(w_2,m_1),(w_3,m_2)\}=f_W. \]
Con le preferenze \emph{vere} $w_1$ ottiene $m_3\succ_{w_1}m_1$: la manipolazione ha successo. Si noti che $f'$ coincide con il matching prodotto dal women's courtship, stabile anche rispetto alle preferenze vere: falsificando le preferenze, $w_1$ ha ottenuto l'esito favorevole al proprio lato.
\end{exbox}

\begin{thmbox}{Numero di matching stabili (Irving--Leather 1986)}
Per ogni $\ell\in\mathbb{N}$ esiste un marriage problem di taglia $n=2^\ell$ con almeno $2^{\,n-1}$ matching stabili.
\end{thmbox}

\section{Altri contesti}

\begin{itemize}
  \item preferenze \textbf{non strette} (indifferenze): tie-breaking arbitrario oppure gestione interna al meccanismo;
  \item numero diverso di uomini e donne;
  \item \textbf{preferenze mancanti}: alternative inaccettabili $\Rightarrow$ si resta single;
  \item \textbf{many-to-one} e \textbf{many-to-many}: college admission (ospedali/specializzandi);
  \item \textbf{pool unico} per i match (\emph{stable roommates}): qui i matching stabili possono \textbf{non esistere};
  \item \textbf{fairness} nella selezione fra i molti matching stabili (né men- né women-optimal).
\end{itemize}

\section{Social choice}

Si passa ora al problema di aggregare le preferenze degli agenti nella scelta di un \emph{unico} bene.

\begin{defbox}{Social choice function}
Dati $N$ (agenti) e $H$ (commodities), sia $P$ la famiglia di tutte le preferenze strette su $H$, con $|P|=|H|!$. Un profilo è $(p_1,\dots,p_n)\in P^N$ (si scrive $\succ_{p_i}$ per la relazione indotta).

Una \textbf{social choice function} è una qualunque
\[ G:P\times\cdots\times P\longrightarrow H. \]
Esempi: maggioranza relativa, metodi di tipo Borda e Condorcet (tutti con regole di tie-breaking).

$G$ è \textbf{dittatoriale} se esiste $i$ (il \emph{dittatore}) con $G(p_1,\dots,p_n)=\max_{p_i}H$ per ogni profilo.
\end{defbox}

\begin{defbox}{Requisiti ragionevoli}
\begin{itemize}
  \item \textbf{unanimità}: $G(p_1,\dots,p_n)=h$ ogni volta che $h=\max_{p_i}H$ per ogni $i\in N$;
  \item \textbf{monotonicità}: se l'implicazione $\bar h\succ_{p_i}h'\Rightarrow \bar h\succ_{q_i}h'$ vale per ogni $h'\in H$ e ogni $i\in N$, allora
  \[ G(p_1,\dots,p_n)=\bar h\ \Longrightarrow\ G(q_1,\dots,q_n)=\bar h. \]
\end{itemize}
\end{defbox}

\begin{thmbox}{Impossibilità di una social choice function ``valida''}
Se $|H|\ge 3$, allora ogni $G$ che soddisfa unanimità e monotonicità è dittatoriale.
\end{thmbox}

\begin{remark}
Con $|H|=2$ la regola di maggioranza semplice soddisfa unanimità e monotonicità senza essere dittatoriale: la soglia $|H|\ge 3$ è essenziale.
\end{remark}

\begin{defbox}{Manipolabilità}
$G$ è \textbf{manipolabile} se esistono un agente $i\in N$, una preferenza stretta $q_i\in P$ e un profilo $(p_1,\dots,p_n)\in P^N$ con
\[ G(q_i,p_{-i})\succ_{p_i}G(p_1,\dots,p_n), \]
cioè un agente può imporre una scelta a lui più gradita dichiarando preferenze false.
\end{defbox}

\begin{thmbox}{Manipulability (Gibbard 1973, Satterthwaite 1975)}
Sia $G$ una social choice function non manipolabile che soddisfa l'unanimità. Se $|H|\ge 3$, allora $G$ è dittatoriale.
\end{thmbox}

\textit{Dimostrazione (sketch).} Il punto chiave è l'implicazione
\[ \text{non manipolabilità}+\text{unanimità}\ \Longrightarrow\ \text{monotonicità}, \]
dopo la quale si applica il teorema di impossibilità precedente. \qed

\begin{remark}
Anche introducendo il caso (meccanismi randomizzati $G:P^N\to\Delta H$) la manipolabilità permane (Gibbard 1977).
\end{remark}

\section*{In sintesi per l'orale}
\addcontentsline{toc}{section}{In sintesi per l'orale (L18)}
\begin{itemize}
  \item \textbf{Matching stabile}: nessuna blocking pair $(w,m)$ che si preferisca reciprocamente ai partner assegnati; equivalentemente, $f$ non è dominato da nessun $\hat f$ tramite una coalizione $S\subseteq M\cup W$ chiusa rispetto al matching.
  \item \textbf{Gale--Shapley}: un matching stabile esiste sempre; la prova è costruttiva (deferred acceptance).
  \item \textbf{Deferred acceptance}: gli uomini scendono lungo le liste, le donne salgono; al più $n(n-1)+1$ iterazioni.
  \item \textbf{Proponent optimality}: il men's courtship è men-optimal e women-disappointing, quindi weak Pareto optimal per gli uomini.
  \item \textbf{Knuth 1976}: sull'insieme dei matching stabili gli interessi dei due lati sono esattamente opposti.
  \item \textbf{Roth 1982}: nessun meccanismo stabile è strategy-proof; ma (Dubins--Freedman) il men's courtship è coalitionally strategy-proof \emph{per i proponenti}; (Gale--Sotomayor) le donne possono manipolare quando esistono almeno due matching stabili.
  \item \textbf{Irving--Leather}: il numero di matching stabili può essere esponenziale ($\ge 2^{n-1}$).
  \item \textbf{Collegamenti}: contrasto con TTC (L17), che \emph{è} strategy-proof perché one-sided; e con Gibbard--Satterthwaite, che è l'analogo risultato di impossibilità per le social choice functions ($|H|\ge3$: non manipolabilità $+$ unanimità $\Rightarrow$ dittatura).
\end{itemize}

% ============================================================
\chapter{Cooperative Maximum Flow e Giochi TU (L19)}
\label{ch:coop-flow}
% ============================================================

\section{Il gioco del massimo flusso cooperativo}

\begin{exbox}{Cooperative maximum flow}
Alcuni trasportatori (\emph{carriers}) gestiscono rotte commerciali \textbf{esclusive} per spedire (gli stessi) beni da una sorgente $s$ a una destinazione $t$; ogni carrier possiede un insieme di archi (nel disegno delle slide, un colore per carrier). Il grafo complessivo descrive la mappa delle rotte con le relative capacità di spedizione.

Il \textbf{worth} (valore) di una coalizione $S$ di carrier è
\[ v(S)=\text{massimo flusso } s\text{--}t \text{ sul sottografo indotto dagli archi posseduti da } S. \]

Nell'istanza delle slide, con $N=\{1,2,3\}$:
\[ v(1)=1,\quad v(2)=2,\quad v(3)=3, \]
\[ v(\{1,2\})=8,\quad v(\{1,3\})=9,\quad v(\{2,3\})=22,\quad v(\{1,2,3\})=29. \]
Si osservi il fortissimo effetto sinergico: mettere in comune gli archi permette di comporre cammini $s$--$t$ che nessun carrier possiede per intero. Il gioco è \textbf{superadditivo}: ad esempio $v(1)+v(\{2,3\})=1+22=23\le 29=v(N)$.
\end{exbox}

Le domande naturali sono: \emph{quali coalizioni si formeranno?} e \emph{come si divide il worth fra i membri?} Per rispondere serve il formalismo dei giochi cooperativi.

\section{Giochi coalizionali con utilità trasferibile}

\begin{defbox}{Gioco coalizionale TU}
Un gioco coalizionale con utilità trasferibile (TU game) è una coppia $(N,v)$ con
\begin{itemize}
  \item $N=\{1,\dots,n\}$: $n$ giocatori che cercano un \emph{accordo vincolante} (binding agreement);
  \item \textbf{coalizione}: un qualunque $S\subseteq N$;
  \item \textbf{funzione caratteristica} $v:\mathcal P(N)\to\mathbb{R}$ con $v(\emptyset)=0$; $v(S)$ è il \emph{worth} della coalizione $S$.
\end{itemize}
Ipotesi implicite del modello:
\begin{itemize}
  \item ogni giocatore può aderire a \emph{una sola} coalizione;
  \item il worth di una coalizione \emph{non dipende} dalle altre coalizioni formate;
  \item l'utilità può essere divisa \emph{arbitrariamente} fra i membri della coalizione (trasferibilità).
\end{itemize}
\end{defbox}

Si scrive $v(i)$ per $v(\{i\})$.

\section{Giochi semplici}

\begin{defbox}{Giochi semplici}
Un TU game $(N,v)$ è \textbf{semplice} se $v(S)\in\{0,1\}$ per ogni $S\subseteq N$. Una coalizione $S$ con $v(S)=1$ si dice \textbf{vincente} (winning coalition).
\end{defbox}

\begin{example}[Maggioranza pesata nei voti di un comitato]
Sia $w_i$ il numero di voti detenuti dal giocatore $i$ e $q$ il quorum richiesto per approvare una proposta:
\[
v_{q,w}(S)=\begin{cases}1 & \text{se }\sum_{i\in S}w_i\ge q,\\[2pt] 0 & \text{altrimenti.}\end{cases}
\]
Nella \emph{procedura bicamerale} una proposta passa solo se approvata in entrambe le camere:
$v(S)=\min\{v_{q,w}(S),\,v_{q',w'}(S)\}$.
\end{example}

\begin{example}[Consiglio di sicurezza ONU: voti con potere di veto]
$n=15$, $P\subseteq N$ membri permanenti con $|P|=5$:
\[
v(S)=\begin{cases}1 & \text{se } P\subseteq S \text{ e } |S|\ge 9,\\[2pt] 0 & \text{altrimenti.}\end{cases}
\]
Ogni membro permanente è un \emph{giocatore con veto}: nessuna coalizione che lo escluda è vincente.
\end{example}

\section{Cost games e savings games}

\begin{exbox}{Alberi ricoprenti minimi}
Alcune città devono collegarsi a una centrale elettrica $r$. I cavi ad alta tensione possono essere costruiti seguendo lo schema e i costi indicati sul grafo. Il \textbf{costo} di una coalizione $S$ di città è
\[ v(S)=\text{minimum spanning tree sul sottografo indotto dai nodi } S\cup\{r\}. \]
Nell'istanza delle slide: $v(1)=40$, $v(2)=5$, $v(\{1,2\})=25$, $v(\{1,2,4\})=35$, $v(\{1,3,4\})=45$, \dots

Trattandosi di \textbf{costi} e non di valori, nell'analisi teorica si \emph{cambiano i segni} (un cost game $(N,c)$ si studia come il TU game $v=-c$, oppure invertendo tutte le disuguaglianze: superadditività diventa subadditività, ecc.).
\end{exbox}

\begin{exbox}{Risparmi aggregati per cooperazione}
Ogni città riceve $40$ milioni di euro di fondi statali per opere pubbliche e vuole collegarsi alla centrale risparmiando il più possibile. Il worth di una coalizione $S$ è
\[ v(S)=\text{fondi complessivi di } S \;-\; \text{costo del MST su } S\cup\{r\}. \]
Con quattro città:
\begin{center}
\begin{tabular}{ccccccccc}
\toprule
$S$ & 1 & 2 & 3 & 4 & 12 & 13 & 14 & 23\\
$v(S)$ & $0$ & $35$ & $35$ & $30$ & $55$ & $45$ & $30$ & $70$\\
\midrule
$S$ & 24 & 34 & 123 & 124 & 134 & 234 & 1234 & \\
$v(S)$ & $65$ & $65$ & $90$ & $85$ & $75$ & $100$ & $120$ & \\
\bottomrule
\end{tabular}
\end{center}
Il gioco è superadditivo (ad es.\ $v(12)+v(34)=55+65=120=v(N)$, $v(1)+v(234)=0+100\le 120$). Le domande sono di nuovo: \emph{quali città sono disposte a cooperare} e \emph{come si riallocano i fondi risparmiati}.
\end{exbox}

\section{Market games}

\begin{defbox}{Market game (vino e bottiglie)}
$N=W\cup B$, dove $W$ è l'insieme dei produttori di vino e $B$ quello dei produttori di bottiglie; $\ell_i$ litri di vino prodotti da $i\in W$ e $b_i$ bottiglie da un litro prodotte da $i\in B$. Il worth di $S$ è
\[ v(S)=\min\Bigl\{\sum_{i\in S\cap W}\ell_i,\ \sum_{i\in S\cap B}b_i\Bigr\} \]
(il vino imbottigliato è l'unico bene vendibile: serve la complementarità fra i due lati).
\end{defbox}

\begin{exbox}{Istanza numerica}
$W=\{1,2,3\}$ con $\ell=(6,5,4)$, $B=\{4,5\}$ con $b=(5,7)$. Alcuni valori:
\[ v(14)=5,\ v(15)=6,\ v(145)=6,\ v(125)=7,\ v(1245)=11,\ v(12345)=\min\{15,12\}=12, \]
mentre $v(S)=0$ per ogni $S$ contenuta interamente in $W$ o in $B$ (una sola ``metà'' del mercato non produce nulla di vendibile).
\end{exbox}

\begin{defbox}{Struttura generale dei market games (Shapley--Shubik 1969)}
\begin{itemize}
  \item $n$ produttori/trader di un paniere di beni;
  \item il paniere posseduto fornisce un'utilità;
  \item lo scambio dei beni è consentito \emph{all'interno} di una coalizione;
  \item worth di una coalizione: il massimo valore della somma delle utilità dei membri ottenibile riallocando i beni fra loro.
\end{itemize}
\end{defbox}

\begin{remark}
Il legame con L17 è evidente: un market game è l'analogo \emph{cardinale e trasferibile} di un'economia di scambio. Là gli agenti avevano solo preferenze ordinali sugli oggetti; qui l'utilità è numerica e liberamente ripartibile fra i membri della coalizione.
\end{remark}

\section{Cooperazione nei giochi non cooperativi}

\begin{defbox}{Valore di sicurezza di una coalizione}
Sia $(N,(\Delta_i)_{i\in N},(u_i)_{i\in N})$ un gioco strategico non cooperativo (in strategie miste). Si supponga ora che i giocatori possano cooperare e condividere le utilità:
\[ \Delta_S=\prod_{i\in S}\Delta_i \]
è l'insieme delle strategie congiunte della coalizione $S$. La \emph{strategia di sicurezza} della coalizione ne determina il worth:
\[ v(S)=\max_{x_S\in\Delta_S}\ \min_{x_{S^c}\in\Delta_{S^c}}\ \sum_{i\in S}u_i(x_S,x_{S^c}),\qquad S^c=N\setminus S. \]
\end{defbox}

\begin{proposition}
Il gioco $(N,v)$ così definito è superadditivo: $S\cap T=\emptyset\Rightarrow v(S)+v(T)\le v(S\cup T)$.
\end{proposition}

\textit{Idea (sketch).} Siano $x_S$ e $x_T$ strategie di sicurezza di $S$ e $T$: la coalizione $S\cup T$ può giocare $(x_S,x_T)$ e garantirsi almeno $v(S)+v(T)$, perché gli avversari residui $(S\cup T)^c$ hanno meno libertà di quanta ne avessero $S^c$ e $T^c$ separatamente. \qed

\section{Superadditività e monotonia}

\begin{defbox}{Superadditività e monotonia}
Un TU game $(N,v)$ è
\begin{itemize}
  \item \textbf{superadditivo} se $S\cap T=\emptyset\ \Rightarrow\ v(S)+v(T)\le v(S\cup T)$;
  \item \textbf{monotono} se $S\subseteq T\ \Rightarrow\ v(S)\le v(T)$.
\end{itemize}
Vale l'implicazione: superadditività $+$ $v\ge 0$ $\Rightarrow$ monotonia.
\end{defbox}

\textit{Dimostrazione.} Se $S\subseteq T$ allora $T=S\cup(T\setminus S)$ con unione disgiunta, quindi $v(T)\ge v(S)+v(T\setminus S)\ge v(S)$, essendo $v\ge0$. \qed

\begin{proposition}
\label{prop:superadd-fam}
Se $(N,v)$ è superadditivo, allora ogni famiglia finita $\{S_i\}_{i\in I}$ di coalizioni a due a due disgiunte soddisfa
\[ \sum_{i\in I}v(S_i)\ \le\ v\Bigl(\bigcup_{i\in I}S_i\Bigr). \]
In particolare ogni $S\subseteq N$ soddisfa $\sum_{i\in S}v(i)\le v(S)$; con $S=N$: $v(1)+\cdots+v(n)\le v(N)$.
\end{proposition}

\textit{Dimostrazione (sketch).} Induzione su $|I|$: il caso $|I|=2$ è la definizione; posto $U=\bigcup_{i\le k}S_i$, da $U\cap S_{k+1}=\emptyset$ segue $\sum_{i\le k+1}v(S_i)\le v(U)+v(S_{k+1})\le v(U\cup S_{k+1})$. \qed

\begin{remark}
La superadditività \textbf{impedisce la formazione di coalizioni ``segrete''}: conviene sempre unirsi piuttosto che agire in sottogruppi separati; ciò giustifica l'assunzione che si formi la \emph{grande coalizione} $N$ e che il problema si riduca alla ripartizione di $v(N)$.
\end{remark}

\section{Superadditive cover}

\begin{defbox}{Copertura superadditiva}
Dato un TU game $(N,v)$, si definisce
\[ v^*(S)=\max_{\mathcal P\in\mathbb P(S)}\ \sum_{T\in\mathcal P}v(T), \]
dove $\mathbb P(S)$ è l'insieme di tutte le partizioni di $S$. Il gioco $(N,v^*)$ è la \textbf{superadditive cover} di $(N,v)$: $v^*(S)$ massimizza l'utilità complessiva ottenibile \emph{spezzando} $S$ in sottocoalizioni.
\end{defbox}

\begin{proposition}
(i) $v^*\ge v$, cioè $v^*(S)\ge v(S)$ per ogni $S\subseteq N$;\quad
(ii) $(N,v^*)$ è superadditivo;\quad
(iii) se $(N,w)$ è superadditivo e $w\ge v$, allora $w\ge v^*$.
\end{proposition}

\textit{Dimostrazione (sketch).} (i) La partizione banale $\mathcal P=\{S\}$ è ammissibile. (ii) Se $S\cap T=\emptyset$ e $\mathcal P_S,\mathcal P_T$ sono partizioni ottime, $\mathcal P_S\cup\mathcal P_T$ è una partizione ammissibile di $S\cup T$, dunque $v^*(S\cup T)\ge v^*(S)+v^*(T)$. (iii) Per ogni partizione $\mathcal P$ di $S$: $\sum_{T\in\mathcal P}v(T)\le\sum_{T\in\mathcal P}w(T)\le w(S)$, dove l'ultima disuguaglianza è la Proposizione~\ref{prop:superadd-fam} applicata a $w$; passando al massimo su $\mathcal P$ si ottiene $v^*(S)\le w(S)$. \qed

Dunque $(N,v^*)$ è il \emph{più piccolo} gioco superadditivo che domina $v$.

\begin{exbox}{Market game con costi di accordo}
Con gli stessi dati dell'esempio precedente (istanza vino/bottiglie), si aggiunga un costo $c$ per ogni accordo bilaterale fra un produttore di vino e uno di bottiglie:
\[ v(S)=\min\Bigl\{\sum_{i\in S\cap W}\ell_i,\ \sum_{i\in S\cap B}b_i\Bigr\}-c\,|S\cap W|\,|S\cap B|. \]
Per $c=1$ alcuni valori sono
\begin{center}
\begin{tabular}{lcccccccc}
\toprule
$S$      & $14$ & $15$ & $25$ & $145$ & $1245$ & $1345$ & $2345$ & $12345$\\
$v(S)$   & $4$  & $5$  & $4$  & $4$   & $7$    & $6$    & $5$    & $6$\\
$v^*(S)$ & $4$  & $5$  & $4$  & $5$   & $9$    & $8$    & $7$    & $9$\\
\bottomrule
\end{tabular}
\end{center}
Il gioco \textbf{non è superadditivo}: $v(14)+v(25)=4+4=8>7=v(1245)$, perché unendo le due coppie si pagano $2\cdot2=4$ costi di accordo invece di $1+1=2$. La copertura ripristina la superadditività: $v^*(1245)=v(15)+v(24)=5+4=9$, ottenuto con la partizione $\{\{1,5\},\{2,4\}\}$; analogamente $v^*(12345)=v(15)+v(24)+v(3)=5+4+0=9$.
\end{exbox}

\section{Equivalenza strategica}

\begin{defbox}{Equivalenza strategica}
Un TU game $(N,w)$ è \textbf{strategicamente equivalente} a $(N,v)$ se esistono $a>0$ e $b\in\mathbb{R}^n$ tali che ogni $S\subseteq N$ soddisfi
\[ w(S)=a\,v(S)+b(S),\qquad b(S)=\sum_{i\in S}b_i. \]
($a$ è un cambio di unità di misura, $b_i$ un valore costante attribuito al giocatore $i$.)
\end{defbox}

\begin{remark}
\begin{itemize}
  \item L'equivalenza strategica è una relazione di equivalenza (riflessiva, simmetrica, transitiva).
  \item La \textbf{superadditività è invariante} per equivalenza strategica (perché $b(S)+b(T)=b(S\cup T)$ per $S,T$ disgiunte).
  \item La \textbf{monotonia non lo è}: se $v(S)=|S|$ il gioco è monotono, ma con $b_i=-2$ si ottiene $w(S)=v(S)+b(S)=-|S|$, che non è monotono.
\end{itemize}
\end{remark}

\begin{defbox}{Classi di equivalenza strategica}
Un TU game $(N,v)$ è
\begin{itemize}
  \item \textbf{0-1 normalizzato} se $v(N)=1$ e $v(i)=0$ per ogni $i\in N$;
  \item \textbf{0-0 normalizzato} se $v(N)=0$ e $v(i)=0$ per ogni $i\in N$;
  \item \textbf{0-(-1) normalizzato} se $v(N)=-1$ e $v(i)=0$ per ogni $i\in N$.
\end{itemize}
\end{defbox}

\begin{thmbox}{Classificazione}
Un TU game $(N,v)$ è strategicamente equivalente a un gioco
\begin{itemize}
  \item[(i)] 0-1 normalizzato se e solo se $\sum_{i\in N}v(i)<v(N)$;
  \item[(ii)] 0-0 normalizzato se e solo se $\sum_{i\in N}v(i)=v(N)$;
  \item[(iii)] 0-(-1) normalizzato se e solo se $\sum_{i\in N}v(i)>v(N)$.
\end{itemize}
\end{thmbox}

\textit{Dimostrazione (sketch).} Si ponga $b_i=-a\,v(i)$, cosicché $w(i)=0$ per ogni $i$ e $w(N)=a\bigl(v(N)-\sum_{i\in N}v(i)\bigr)$. Nei casi (i) e (iii) basta scegliere $a$ pari al reciproco del valore assoluto di tale differenza per ottenere $w(N)=1$ o $w(N)=-1$; nel caso (ii) ogni $a>0$ dà $w(N)=0$. Il viceversa segue dal fatto che il segno di $v(N)-\sum_i v(i)$ non cambia per equivalenza strategica. \qed

\begin{remark}
Il segno di $v(N)-\sum_{i\in N}v(i)$ classifica dunque ogni gioco in una delle tre famiglie: il caso $\sum_i v(i)<v(N)$ è quello in cui la cooperazione crea valore, e la forma 0-1 normalizzata ne è il rappresentante canonico.
\end{remark}

\section*{In sintesi per l'orale}
\addcontentsline{toc}{section}{In sintesi per l'orale (L19)}
\begin{itemize}
  \item \textbf{Cooperative max-flow}: $v(S)=$ massimo flusso $s$--$t$ sul sottografo degli archi di $S$; esempio delle slide $v(1)=1,v(2)=2,v(3)=3,v(12)=8,v(13)=9,v(23)=22,v(N)=29$. Motiva le due domande fondamentali: quali coalizioni si formano e come si divide il worth.
  \item \textbf{TU game}: $(N,v)$, $v(\emptyset)=0$; ipotesi: una sola coalizione per giocatore, indipendenza dalle altre coalizioni, utilità arbitrariamente divisibile.
  \item \textbf{Giochi semplici}: $v(S)\in\{0,1\}$; maggioranza pesata $v_{q,w}$, procedura bicamerale come $\min$, consiglio ONU con giocatori di veto.
  \item \textbf{Cost games}: MST per collegare le città alla centrale; con i costi si invertono i segni. Il savings game associato (fondi meno costo del MST) è superadditivo.
  \item \textbf{Market games} (Shapley--Shubik): $v(S)=\min\{\sum\ell_i,\sum b_i\}$; worth $=$ massimo valore ottenibile riallocando i beni dentro la coalizione.
  \item \textbf{Valore di sicurezza}: $v(S)=\max_{x_S}\min_{x_{S^c}}\sum_{i\in S}u_i$; è sempre superadditivo, quindi ogni gioco non cooperativo induce un TU game superadditivo.
  \item \textbf{Superadditività}: implica $\sum_{i\in S}v(i)\le v(S)$ e giustifica la grande coalizione (niente coalizioni ``segrete''); con $v\ge0$ implica monotonia.
  \item \textbf{Superadditive cover} $v^*(S)=\max_{\mathcal P}\sum_{T\in\mathcal P}v(T)$: è il minimo gioco superadditivo $\ge v$. Esempio dei market games con costo d'accordo $c=1$, dove $v(14)+v(25)>v(1245)$.
  \item \textbf{Equivalenza strategica} $w=av+b$, $a>0$: preserva la superadditività ma non la monotonia; il segno di $v(N)-\sum_i v(i)$ classifica il gioco nelle normalizzazioni 0-1, 0-0, 0-(-1).
  \item \textbf{Collegamenti}: max-flow game, cost/savings game e market game sono i tre esempi con cui le slide introducono i giochi TU; le due domande che pongono (quali coalizioni si formano, come si divide il worth) sono quelle cui rispondono i concetti di soluzione cooperativi. Il parallelo ordinale è la stabilità del matching di L17.
\end{itemize}

% ============================================================
%  Capitoli L20 -- L24: Giochi coalizionali con utilità trasferibile
%  (da inserire nel mainmatter di dispensa_AGT.tex, dopo il cap. L13)
% ============================================================

\chapter{Giochi coalizionali TU: imputazioni, stable sets e core (L20)}
\label{ch:tu-core}
% ============================================================

\section{Il modello}

Fino a qui i giocatori agivano in modo \emph{non cooperativo}: ognuno sceglieva la
propria strategia e l'unica forma di coordinamento possibile era quella
autoimposta dall'equilibrio. Nei \textbf{giochi coalizionali} cambia
radicalmente il punto di vista: i giocatori possono stipulare \emph{accordi
vincolanti} (\textit{binding agreements}) e l'oggetto di studio non è più il
profilo di strategie, ma la \emph{ripartizione del valore} generato dalla
cooperazione.

\begin{defbox}{Gioco coalizionale con utilità trasferibile (TU)}
Sia $N=\{1,\dots,n\}$ l'insieme dei giocatori. Una \textbf{coalizione} è un
sottoinsieme $S\subseteq N$; $\mathcal{P}(N)$ denota l'insieme delle parti di $N$.
Un \textbf{gioco coalizionale TU} è una coppia $(N,v)$ dove
\[
  v:\mathcal{P}(N)\to\mathbb{R},\qquad v(\emptyset)=0
\]
è la \textbf{funzione caratteristica}: $v(S)$ è il \emph{worth} (valore) che la
coalizione $S$ è in grado di garantirsi da sola.
\end{defbox}

L'aggettivo \emph{transferable utility} significa che il valore $v(S)$ è una
quantità liberamente divisibile e trasferibile fra i membri di $S$ (denaro): non
serve modellare le preferenze individuali sui panieri, basta la funzione $v$.

\medskip
\noindent\textbf{La domanda guida.} Supponiamo che si formi la \emph{grande
coalizione} $N$. Come si dividono i giocatori il valore $v(N)$? Un candidato è
un vettore $x\in\mathbb{R}^n$ con
\[
  x(N)\;=\;\sum_{i\in N}x_i\;=\;v(N),
  \qquad\text{dove in generale } x(S)=\sum_{i\in S}x_i .
\]
La teoria serve a rispondere alla domanda: \emph{quali $x$ sono
ragionevoli/accettabili?}

\section{Tre famiglie di esempi}

\begin{exbox}{Cooperative maximum flow}
Alcuni vettori (\textit{carriers}) gestiscono rotte commerciali esclusive per
spedire la stessa merce da $s$ a $t$; ogni arco del grafo appartiene a un vettore
e ha una capacità. Il worth di una coalizione $S$ è il \textbf{massimo flusso
$s$--$t$ nel sottografo indotto dagli archi posseduti dai membri di $S$}.
Sull'istanza vista a lezione:
\[
  v(1)=1,\quad v(2)=2,\quad v(3)=3,
\]
\[
  v(\{1,2\})=8,\quad v(\{1,3\})=9,\quad v(\{2,3\})=22,\quad v(\{1,2,3\})=29 .
\]
Si noti la forte \emph{superadditività}: mettere insieme gli archi apre cammini
che nessun vettore possedeva per intero.
\end{exbox}

\begin{exbox}{Aggregate savings by cooperation (opere pubbliche)}
Quattro città ricevono ciascuna 40 milioni di euro di fondi statali e devono essere
collegate a una centrale elettrica $r$. Il worth di una coalizione $S$ è
\[
  v(S)\;=\;40\,|S|\;-\;\text{costo del minimum spanning tree su } S\cup\{r\},
\]
cioè i fondi complessivi meno il costo della rete: è il \emph{risparmio
aggregato}. I dati dell'istanza:
\begin{center}
\begin{tabular}{cc@{\qquad}cc}
\toprule
$S$ & $v(S)$ & $S$ & $v(S)$\\
\midrule
$1$ & $0$   & $24$   & $65$\\
$2$ & $35$  & $34$   & $65$\\
$3$ & $35$  & $123$  & $90$\\
$4$ & $30$  & $124$  & $85$\\
$12$& $55$  & $134$  & $75$\\
$13$& $45$  & $234$  & $100$\\
$14$& $30$  & $1234$ & $120$\\
$23$& $70$  &        & \\
\bottomrule
\end{tabular}
\end{center}
La domanda diventa: \emph{come si riallocano i fondi risparmiati fra le città?
quali città sono disposte a cooperare?}
\end{exbox}

\begin{exbox}{Market game: vino e bottiglie}
$N=W\cup B$, con $W$ produttori di vino ($\ell_i$ litri per $i\in W$) e $B$
produttori di bottiglie ($b_i$ bottiglie da un litro per $i\in B$). Il vino
sfuso e le bottiglie vuote non valgono nulla: si vende solo il vino
imbottigliato, quindi
\[
  v(S)\;=\;\min\Bigl\{\sum_{i\in S\cap W}\ell_i,\;\sum_{i\in S\cap B}b_i\Bigr\}.
\]
Con $W=\{1,2,3\}$, $\ell=(6,5,4)$ e $B=\{4,5\}$, $b=(5,7)$, si ottiene ad esempio
$v(\{1,4\})=5$, $v(\{1,5\})=6$, $v(\{1,2,4,5\})=11$, $v(N)=\min\{15,12\}=12$;
tutte le coalizioni contenute in $W$ o in $B$ valgono $0$.
\end{exbox}

\section{Imputazioni e dominanza}

\begin{defbox}{Imputazione}
$x\in\mathbb{R}^n$ è un'\textbf{imputazione} per $(N,v)$ se è
\begin{itemize}
  \item \textbf{socialmente razionale} (o \emph{efficiente}): $x(N)=v(N)$;
  \item \textbf{individualmente razionale}: $x_i\ge v(i)$ per ogni $i\in N$.
\end{itemize}
Si scrive $X(N,v)=\{\text{imputazioni}\}$.
\end{defbox}

L'efficienza dice che si distribuisce esattamente tutto il valore prodotto (né
si spreca né si inventa); la razionalità individuale dice che nessuno accetta
meno di quanto otterrebbe da solo.

\begin{defbox}{Dominanza}
Siano $x,y\in X(N,v)$ e $S\subseteq N$. Si dice che $y$ \textbf{domina} $x$
\emph{attraverso} $S$ se
\begin{itemize}
  \item $y_i>x_i$ per ogni $i\in S$ \quad (tutti i membri di $S$ ci guadagnano);
  \item $y(S)\le v(S)$ \quad (la coalizione $S$ è in grado di realizzare $y$ da sola).
\end{itemize}
$x$ è \textbf{non dominata} se non esistono $y\in X(N,v)$ e $S\subseteq N$ tali
che $y$ domina $x$ attraverso $S$.
\end{defbox}

La dominanza è l'analogo cooperativo della deviazione profittevole: una
coalizione che può fare di meglio da sola ha un motivo per rompere l'accordo.

\section{Stable sets di von Neumann--Morgenstern}

\begin{defbox}{Stable set (von Neumann--Morgenstern, 1944)}
$G\subseteq X(N,v)$ è uno \textbf{stable set} per $(N,v)$ se è
\begin{itemize}
  \item \textbf{internamente stabile}: nessun $x\in G$ è dominato da un altro $y\in G$;
  \item \textbf{esternamente stabile}: ogni imputazione $x\notin G$ è dominata da qualche $y\in G$.
\end{itemize}
\end{defbox}

\begin{exbox}{Gloves game}
I giocatori 1 e 2 possiedono un guanto destro ciascuno, il giocatore 3 un guanto
sinistro; solo una coppia completa ha valore:
\[
  v(1)=v(2)=v(3)=v(\{1,2\})=0,\qquad
  v(\{1,3\})=v(\{2,3\})=v(\{1,2,3\})=1 .
\]
Qui $X(N,v)$ è il simplesso unitario di $\mathbb{R}^3$. Sono stable sets
\[
  G_1=\{(t,t,1-2t):0\le t\le \tfrac12\},
\]
\[
  G_2=\{(t,0,1-t):0\le t\le 1\},\qquad
  G_3=\{(0,t,1-t):0\le t\le 1\}.
\]
Già un gioco a tre giocatori ha dunque \emph{più} stable sets, di natura
qualitativamente diversa ($G_1$ ``egualitario'' fra 1 e 2, $G_2,G_3$
``discriminatori'').
\end{exbox}

\begin{remark}[Perché gli stable sets non hanno avuto successo]
Come osserva Aumann, il problema della teoria degli stable sets è che è
estremamente difficile lavorarci: trovare gli stable sets richiede un nuovo
\emph{tour de force} matematico per ogni gioco o classe di giochi considerata; a
parte poche verità elementari, non c'è teoria, non ci sono strumenti e
certamente non ci sono algoritmi. Da qui l'interesse per un concetto più
maneggevole: il \emph{core}.
\end{remark}

\section{Il core}

\begin{defbox}{Core di un gioco (Gillies, 1959)}
\[
  C(N,v)\;=\;\bigl\{x\in\mathbb{R}^n:\;x(N)=v(N),\;\;x(S)\ge v(S)\;\;\forall\,S\subseteq N\bigr\}.
\]
Sono le allocazioni efficienti che risultano \emph{razionali per ogni possibile
coalizione}: nessun gruppo di giocatori ha interesse a staccarsi.
\end{defbox}

\begin{proposition}[Proprietà elementari del core]
\begin{itemize}
  \item[(i)] $C(N,v)\subseteq X(N,v)$: prendendo $S=\{i\}$ si ritrova la razionalità individuale.
  \item[(ii)] $C(N,v)$ giace sull'iperpiano $\{x:x_1+\dots+x_n=v(N)\}$ ed è un \textbf{poliedro compatto}, in particolare \emph{convesso}.
  \item[(iii)] Ogni $x\in C(N,v)$ è non dominata; dunque il core è internamente stabile.
\end{itemize}
\end{proposition}

\textit{Dimostrazione (sketch).} (i) e (ii) sono immediate: il core è definito da
un'equazione lineare e da $2^n-1$ disuguaglianze lineari, quindi è un poliedro
convesso; la limitatezza segue perché $x_i\le v(N)-\sum_{j\ne i}v(j)$ per ogni
$i$. (iii) Se $y$ dominasse $x\in C(N,v)$ attraverso $S$, allora
$v(S)\ge y(S)>x(S)\ge v(S)$, assurdo. \qed

\begin{example}[Core del gloves game]
Per il gioco dei guanti $X(N,v)$ è il simplesso unitario, mentre
\[
  C(N,v)=\{(0,0,1)\}.
\]
Infatti $x_1+x_3\ge 1$ e $x_2+x_3\ge 1$ con $x_1+x_2+x_3=1$ e $x\ge 0$ forzano
$x_1=x_2=0$. Interpretazione: \emph{nessuna coalizione vincente è possibile senza
il giocatore 3}, che quindi si prende tutto. La competizione fra 1 e 2 azzera il
loro potere contrattuale.
\end{example}

\subsection{Un esempio parametrico a tre giocatori}

\begin{exbox}{Il gioco $\rho$}
Sia $n=3$ e
\[
  v(S)=\begin{cases}
   0 & \text{se } |S|=1,\\
   \rho & \text{se } |S|=2,\\
   3 & \text{se } |S|=3.
  \end{cases}
\]
Le imputazioni sono $x\ge 0$ con $x_1+x_2+x_3=3$; le condizioni di core sono
$x_i+x_j\ge\rho$, cioè $3-x_k\ge\rho$, cioè $x_k\le 3-\rho$ per ogni $k$.
Sommando le tre: $3\le 3(3-\rho)$, cioè $\rho\le 2$. Quindi
\begin{itemize}
  \item $C(N,v)\ne\emptyset \iff \rho\le 2$;
  \item $C(N,v)=\{(1,1,1)\}$ per $\rho=2$ (le tre disuguaglianze diventano attive);
  \item $|C(N,v)|=2^{\aleph_0}$ per $\rho<2$ (il core è un triangolo pieno);
  \item $(N,v)$ è superadditivo $\iff 0\le\rho\le 3$;
  \item se $2<\rho\le 3$, l'imputazione ``equa'' $(1,1,1)$ è \emph{dominata} da
        $(\rho/2,\rho/2,3-\rho)$ attraverso $\{1,2\}$: infatti $\rho/2>1$ e
        $\rho/2+\rho/2=\rho\le v(\{1,2\})$.
\end{itemize}
Morale: la superadditività \emph{non} basta a garantire il core non vuoto.
\end{exbox}

\section{Superadditività e covarianza}

\begin{defbox}{Superadditività}
$(N,v)$ è \textbf{superadditivo} se
\[
  S\cap T=\emptyset\;\Longrightarrow\; v(S)+v(T)\le v(S\cup T).
\]
Cioè: fondere due coalizioni disgiunte non fa mai perdere valore.
\end{defbox}

\begin{thmbox}{Proposizione (superadditività e trasformazioni affini)}
\begin{itemize}
  \item[(i)] Se $(N,v)$ è superadditivo, allora $C(N,v)$ coincide con l'insieme
  di tutte le imputazioni non dominate.
  \item[(ii)] Se $w(S)=a\,v(S)+b(S)$ con $a>0$ e $b\in\mathbb{R}^n$ (\emph{equivalenza
  strategica}), allora $C(N,w)=a\,C(N,v)+b$.
\end{itemize}
\end{thmbox}

In (i) l'inclusione $\subseteq$ è il punto (iii) della proposizione precedente;
per (ii) basta osservare che, posto $y=ax+b$, si ha $w(S)-y(S)=a\,(v(S)-x(S))$:
le disuguaglianze del core si trasformano una nell'altra.

\begin{remark}
La (ii) è molto usata nei conti: permette di \emph{normalizzare} un gioco (per
esempio portandolo in forma $0$-$1$) senza cambiare la struttura del core, e vale
in modo identico per il nucleolo (Cap.~\ref{ch:nucleolo}) e per il valore di
Shapley (Cap.~\ref{ch:shapley}).
\end{remark}

\section*{In sintesi per l'orale}
\addcontentsline{toc}{section}{In sintesi per l'orale (L20)}
\begin{itemize}
  \item Un gioco TU è $(N,v)$ con $v(\emptyset)=0$: si abbandonano le strategie e
        si ragiona su \emph{quanto vale ogni coalizione}. Esempi tipo: max flow
        cooperativo, risparmi aggregati (MST), market game vino/bottiglie.
  \item Imputazione = efficienza $x(N)=v(N)$ + razionalità individuale
        $x_i\ge v(i)$. Preimputazione (L22) = solo efficienza.
  \item Dominanza attraverso $S$: tutti in $S$ migliorano e $S$ si può permettere
        $y$ da sola ($y(S)\le v(S)$).
  \item Stable sets (vNM): internamente + esternamente stabili. Esistono in
        genere in numero infinito (gloves game: $G_1,G_2,G_3$), non c'è
        algoritmo generale: da qui il passaggio al core.
  \item Core $=\{x: x(N)=v(N),\ x(S)\ge v(S)\ \forall S\}$: poliedro compatto e
        convesso, contenuto nelle imputazioni, fatto di imputazioni non dominate.
        \emph{Può essere vuoto}.
  \item Se il gioco è superadditivo, core $=$ imputazioni non dominate; e il
        core è covariante per equivalenza strategica ($w=av+b$).
  \item Esempio chiave da ricordare: gioco $\rho$ a 3 giocatori --- core non vuoto
        sse $\rho\le 2$, singoletto $(1,1,1)$ per $\rho=2$; superadditivo per
        $\rho\le 3$. Quindi \emph{superadditivo $\not\Rightarrow$ core $\ne\emptyset$}.
  \item Collegamento: il core sta ai giochi cooperativi come il NE sta ai giochi
        non cooperativi (nessuna deviazione conveniente), ma la deviazione è di
        \emph{gruppo}; ed è un insieme, non un punto --- da cui la necessità di
        raffinamenti a valore singolo (nucleolo, Shapley).
\end{itemize}

% ============================================================
\chapter{Bondareva--Shapley, giochi bilanciati e market games (L21)}
\label{ch:bondareva}
% ============================================================

\section{Il core di un gioco superadditivo: un esempio completo}

Riprendiamo il gioco dei risparmi aggregati (opere pubbliche) del
Cap.~\ref{ch:tu-core}. Il gioco è superadditivo, quindi il suo core coincide con
l'insieme delle imputazioni non dominate.

\begin{exbox}{Core del gioco dei risparmi aggregati}
Cerchiamo $x$ con $x(N)=120$ e $x(S)\ge v(S)$ per ogni $S$. Le disuguaglianze
attive sono poche:
\begin{itemize}
  \item $x_3\ge 35$ e $x_{123}\ge 90$, cioè $x_4\le 30$; con $x_4\ge 30$ segue $x_4=30$;
  \item $x_{124}\ge 85$ dà $x_3\le 35$, dunque $x_3=35$;
  \item $x_{12}\ge 55$ e $x_{34}=65$ danno $x_1+x_2=55$;
  \item $x_{13}\ge 45$ impone $x_1\ge 10$; $x_{234}\ge 100$ impone $x_1\le 20$.
\end{itemize}
Dunque
\[
  C(N,v)=\{(t,\;55-t,\;35,\;30)\;:\;10\le t\le 20\}.
\]
Il core è un segmento: le città 3 e 4 sono ``bloccate'' sul loro valore
individuale, mentre 1 e 2 si spartiscono 55 con un margine di contrattazione di
ampiezza 10. Si noti che il giocatore 1 da solo vale $0$ ma nel core prende
almeno 10: è il valore che porta alle coalizioni $\{1,3\}$ e $\{1,3,4\}$.
\end{exbox}

\section{Condizioni per la non vuotezza del core}

\begin{thmbox}{Proposizione (condizione necessaria)}
\begin{itemize}
  \item[(i)] Se $C(N,v)\ne\emptyset$, allora per ogni partizione $\{S_1,\dots,S_k\}$ di $N$
  \[
    v(S_1)+\dots+v(S_k)\;\le\;v(N). \tag{$*$}
  \]
  \item[(ii)] $C(N,v)=\emptyset$ per ogni gioco strategicamente equivalente a un
  gioco $0$-$(-1)$ normalizzato (cioè $v(i)=0$ per ogni $i$ e $v(N)=-1$).
\end{itemize}
\end{thmbox}

\textit{Dimostrazione.} (i) Sia $x\in C(N,v)$: sommando $x(S_j)\ge v(S_j)$ su
$j=1,\dots,k$ e usando che i $S_j$ partizionano $N$ si ottiene
$\sum_j v(S_j)\le \sum_j x(S_j)=x(N)=v(N)$. (ii) In un gioco $0$-$(-1)$
normalizzato la partizione in singoletti viola $(*)$: $\sum_i v(i)=0>-1=v(N)$;
si conclude con la covarianza del core. \qed

\begin{remark}
Se $(N,v)$ è superadditivo, ogni partizione soddisfa $(*)$ automaticamente
(per induzione sulla fusione dei blocchi). Quindi $(*)$ è \emph{necessaria ma non
sufficiente}: il gioco $\rho$ con $2<\rho\le 3$ è superadditivo e ha core vuoto.
Serve una condizione che consideri anche famiglie di coalizioni \emph{non}
disgiunte: sono i sistemi di pesi bilancianti.
\end{remark}

\begin{thmbox}{Caratterizzazione via programmazione lineare}
$C(N,v)\ne\emptyset$ se e solo se $v_P\le v(N)$, dove $v_P$ è il valore ottimo del
problema di programmazione lineare
\[
  (P)\qquad \min\Bigl\{\sum_{i\in N}x_i\;:\;\sum_{i\in S}x_i\ge v(S)\;\;\forall S\subseteq N,\,S\ne N\Bigr\}
\]
ovvero $\min\{x(N):x(S)\ge v(S)\ \ S\subsetneq N\}$.
\end{thmbox}

\textit{Dimostrazione.} Se $v_P\le v(N)$, sia $\bar x$ ottimo per $(P)$; aggiungendo
$v(N)-v_P\ge 0$ a una qualsiasi componente si ottiene un punto ammissibile con
$x(N)=v(N)$, che sta nel core (le disuguaglianze restano soddisfatte perché si
è solo aumentato $x$). Viceversa, ogni $x$ del core è ammissibile per $(P)$ con
$x(N)=v(N)$, dunque $v_P\le v(N)$. \qed

\section{Pesi bilancianti e giochi bilanciati}

\begin{defbox}{Vettore di pesi bilancianti e gioco bilanciato}
Sia $\mathcal{C}=\mathcal{P}(N)\setminus\{\emptyset,N\}$. La famiglia
$\{\lambda_S\}_{S\in\mathcal{C}}$ è un \textbf{vettore di pesi bilancianti} se
\begin{itemize}
  \item $\lambda_S\ge 0$ per ogni $S\in\mathcal{C}$;
  \item per ogni giocatore $i\in N$: $\displaystyle\sum_{S\in\mathcal{P}(i)}\lambda_S=1$, dove
        $\mathcal{P}(i)=\{S\in\mathcal{C}: i\in S\}$.
\end{itemize}
Il gioco $(N,v)$ è \textbf{bilanciato} (\emph{balanced}) se \emph{ogni} vettore di
pesi bilancianti soddisfa
\[
  \sum_{S\in\mathcal{C}}\lambda_S\,v(S)\;\le\;v(N).
\]
\end{defbox}

\begin{remark}[Interpretazione ``a frazioni di tempo'']
$\lambda_S$ è la frazione di tempo in cui i membri di $S$ lavorano per $S$.
Dedicare tempo pieno a $S$ produce $v(S)$; dedicarne una frazione $\lambda_S$
produce $\lambda_S v(S)$. La condizione $\sum_{S\ni i}\lambda_S=1$ dice che ogni
giocatore è occupato esattamente al 100\%. Bilanciato significa allora: nessuna
organizzazione ``a tempo parziale'' del lavoro fa meglio della grande coalizione.
Le partizioni sono il caso particolare $\lambda_S\in\{0,1\}$: la condizione $(*)$
è quindi un caso speciale della bilanciatezza.
\end{remark}

\begin{remark}
La bilanciatezza è \emph{invariante per equivalenza strategica}: se $w=av+b$ con
$a>0$, allora $\sum_S\lambda_S w(S)=a\sum_S\lambda_S v(S)+b(N)$ (usando
$\sum_{S\ni i}\lambda_S=1$), e la disuguaglianza si conserva.
\end{remark}

\begin{thmbox}{Teorema di Bondareva (1963) -- Shapley (1967)}
\[
  C(N,v)\ne\emptyset \iff (N,v)\ \text{è bilanciato.}
\]
\end{thmbox}

\textit{Dimostrazione (sketch, via dualità lineare).}
\begin{itemize}
  \item La non vuotezza del core equivale a $v_P\le v(N)$, con $v_P$ valore ottimo di
    \[
      (P)\quad \min\Bigl\{\sum_{i\in N}x_i:\sum_{i\in S}x_i\ge v(S),\ S\subsetneq N\Bigr\}.
    \]
  \item Il duale di $(P)$ ha una variabile $\lambda_S\ge 0$ per ogni vincolo
    $S\in\mathcal{C}$ e un vincolo per ogni giocatore $i$ (la colonna di $x_i$):
    \[
      (D)\quad \max\Bigl\{\sum_{S\in\mathcal{C}}\lambda_S v(S)\;:\;\{\lambda_S\}\ \text{è un vettore di pesi bilancianti}\Bigr\},
    \]
    perché la colonna di $x_i$ ha coefficiente $1$ in tutti e soli i vincoli con
    $i\in S$, e il costo di $x_i$ è $1$: il vincolo duale è
    $\sum_{S\ni i}\lambda_S= 1$.
  \item Essere bilanciato equivale quindi a $v_D\le v(N)$.
  \item $(P)$ è ammissibile e limitato inferiormente, quindi per la
    \emph{dualità forte} $v_P=v_D$; dunque $v_P\le v(N)\iff v_D\le v(N)$. \qed
\end{itemize}

\begin{remark}
Il teorema è il ponte fra un oggetto combinatorio (le coperture frazionarie di
$N$ mediante coalizioni) e uno geometrico (la non vuotezza di un poliedro). Si noti
che $(P)$ ha $n$ variabili ma $2^n-2$ vincoli.
\end{remark}

\section{Balanced cover}

\begin{defbox}{Balanced cover}
Dato $(N,v)$, si definisce
\[
  v^{\star}(S)=\begin{cases}
    v(S) & \text{se } S\ne N,\\[2pt]
    \max\Bigl\{\max\bigl\{\sum_{S\in\mathcal{C}}\lambda_S v(S):\{\lambda_S\}\ \text{pesi bilancianti}\bigr\},\;v(N)\Bigr\} & \text{se } S=N.
  \end{cases}
\]
$(N,v^{\star})$ è il \textbf{balanced cover} di $(N,v)$: si aumenta $v(N)$ del
minimo indispensabile per rendere il gioco bilanciato.
\end{defbox}

\begin{proposition}
\begin{itemize}
  \item[(i)] $v^{\star}\ge v$, cioè $v^{\star}(S)\ge v(S)$ per ogni $S\subseteq N$;
  \item[(ii)] $(N,v^{\star})$ è bilanciato;
  \item[(iii)] $(N,v)$ è bilanciato se e solo se $v^{\star}=v$;
  \item[(iv)] se $(N,w)$ è bilanciato e $w\ge v$, allora $w\ge v^{\star}$ (minimalità).
\end{itemize}
\end{proposition}

\textit{Dimostrazione (sketch).} (i) e (iii) seguono dalla definizione (il $\max$
con $v(N)$ garantisce (i); se il gioco è bilanciato il primo termine è
$\le v(N)$). (ii) I pesi bilancianti non coinvolgono $N$, quindi il valore ottimo
di $(D)$ è lo stesso per $v$ e $v^{\star}$, ed è per costruzione
$\le v^{\star}(N)$. (iv) Se $w\ge v$ e $w$ bilanciato, allora per ogni sistema
bilanciato $\sum_S\lambda_S v(S)\le\sum_S\lambda_S w(S)\le w(N)$; passando al max
si ha $v^{\star}(N)\le w(N)$, mentre per $S\ne N$ vale $v^{\star}(S)=v(S)\le w(S)$. \qed

\section{Market games}

\begin{defbox}{Market game (Shapley--Shubik, 1969)}
Si considerino $n$ produttori/commercianti di $p$ merci:
\begin{itemize}
  \item $w^i\in\mathbb{R}^p_+$ è la \emph{dotazione} (endowment) del giocatore $i$;
  \item $u_i:\mathbb{R}^p_+\to\mathbb{R}$ è la sua utilità (funzione del paniere posseduto);
  \item dentro una coalizione è consentito scambiare/riallocare le merci.
\end{itemize}
$(N,v)$ è un \textbf{market game} se ogni $u_i$ è \emph{concava} e il worth di
ogni coalizione $S$ è
\[
  v(S)=\max\Bigl\{\sum_{i\in S}u_i(z^i)\;:\;\sum_{i\in S}z^i=\sum_{i\in S}w^i,\;\;z^i\in\mathbb{R}^p_+\Bigr\}.
\]
In particolare $v(i)=u_i(w^i)$.
\end{defbox}

\begin{remark}
La classe è chiusa per equivalenza strategica: se $w(S)=a v(S)+b(S)$ con $a>0$,
allora $(N,w)$ è il market game con utilità $z^i\mapsto a\,u_i(z^i)+b_i$.
\end{remark}

\begin{exbox}{Vino e bottiglie come market game}
Con $p=2$ merci (vino, bottiglie), $w^i=(\ell_i,b_i)$ e
$u_i(z^i)=\min\{z^i_1,z^i_2\}$ (utilità di Leontief, concava) si ottiene
esattamente
\[
  v(S)=\min\Bigl\{\sum_{i\in S\cap P}\ell_i,\;\sum_{i\in S\cap B}b_i\Bigr\},
\]
perché l'allocazione ottima concentra tutte le coppie vino--bottiglia. Con
$\ell=(6,5,4)$ per $P=\{1,2,3\}$ e $b=(5,7)$ per $B=\{4,5\}$: $v(\{1,4\})=5$,
$v(\{1,2,5\})=7$, $v(\{1,2,4,5\})=11$, $v(N)=12$.
\end{exbox}

\begin{exbox}{Market game con intermediari}
$N=P\cup B\cup R$, con $R$ insieme di \emph{rivenditori}. Produttori e
imbottigliatori non ricavano nulla dal possesso ($u_i\equiv 0$ per $i\in P\cup B$),
mentre un rivenditore $i\in R$ ha
\[
  u_i(z^i)=\alpha_i\min\{z^i_1,z^i_2\}+\sum_{j=1}^{2}\beta^i_j\max\{0,\;z^i_j-\min\{z^i_1,z^i_2\}\},
  \qquad \alpha_i>\beta^i_1+\beta^i_2 ,
\]
cioè valorizza le coppie complete ($\alpha_i$) e, in misura minore, gli avanzi
($\beta^i_j$). L'ipotesi $\alpha_i>\beta^i_1+\beta^i_2$ rende $u_i$ concava e
premia l'abbinamento. Nell'istanza a 6 giocatori vista a lezione si ottengono
per esempio $v(\{1,2,6\})=22$, $v(\{1,3,5\})=20$, $v(N)=45$: senza rivenditori
il worth è $0$, perché nessuno sa monetizzare la merce.
\end{exbox}

\section{Sottogiochi e giochi totalmente bilanciati}

\begin{thmbox}{Teorema}
Ogni market game $(N,v)$ è bilanciato; in particolare il core di un market game
è non vuoto.
\end{thmbox}

\begin{defbox}{Sottogioco e gioco totalmente bilanciato}
$(S,w)$ è un \textbf{sottogioco} di $(N,v)$ se $S\subseteq N$ e $w$ è la
restrizione di $v$ a $\mathcal{P}(S)$, cioè $w(T)=v(T)$ per ogni $T\subseteq S$.

$(N,v)$ è \textbf{totalmente bilanciato} (\emph{totally balanced}) se tutti i suoi
sottogiochi sono bilanciati.
\end{defbox}

Ogni sottogioco di un market game è ancora un market game (basta restringere la
lista dei commercianti), dunque è bilanciato: i market games sono totalmente
bilanciati. Vale anche il viceversa.

\begin{thmbox}{Teorema di caratterizzazione}
\begin{itemize}
  \item[(i)] $(N,v)$ è totalmente bilanciato $\iff$ $(N,v)$ è un market game;
  \item[(ii)] $(N,v)$ è totalmente bilanciato $\iff$ $(N,v)$ è un maximum flow game.
\end{itemize}
\end{thmbox}

\begin{remark}
Corollario pratico: i due esempi motivanti del Cap.~\ref{ch:tu-core} (max flow
cooperativo e vino/bottiglie) hanno \emph{sempre} core non vuoto. Il gioco dei
risparmi aggregati e il market game \emph{con costi di accordo} (Cap.~\ref{ch:coal-struct})
invece no: il secondo ha core vuoto, e infatti non è un market game.
\end{remark}

\section*{In sintesi per l'orale}
\addcontentsline{toc}{section}{In sintesi per l'orale (L21)}
\begin{itemize}
  \item Condizione necessaria: se $C\ne\emptyset$ allora ogni partizione soddisfa
        $\sum_j v(S_j)\le v(N)$. I giochi $0$-$(-1)$ normalizzati (e i loro
        equivalenti strategici) hanno core vuoto.
  \item Core non vuoto $\iff v_P\le v(N)$ per l'LP
        $(P)\ \min\{x(N):x(S)\ge v(S),\,S\subsetneq N\}$.
  \item Pesi bilancianti: $\lambda_S\ge 0$ con $\sum_{S\ni i}\lambda_S=1$ per ogni
        $i$. Interpretazione: frazioni di tempo. Le partizioni sono il caso
        $0/1$.
  \item \textbf{Bondareva--Shapley}: $C(N,v)\ne\emptyset\iff$ il gioco è bilanciato.
        Dimostrazione da saper abbozzare: $(D)$ è il duale di $(P)$, poi
        dualità forte $v_P=v_D$.
  \item Balanced cover $v^{\star}$: alza $v(N)$ al minimo necessario; è il più
        piccolo gioco bilanciato che domina $v$.
  \item Market game (Shapley--Shubik): utilità concave + worth come massimo
        della riallocazione interna. Ogni market game è bilanciato, quindi ha core non vuoto.
  \item Totalmente bilanciato $\iff$ market game $\iff$ maximum flow game. Ogni
        sottogioco di un market game è un market game.
  \item Collegamento: la bilanciatezza è l'analogo cooperativo delle condizioni
        di esistenza via punto fisso nei giochi non cooperativi; qui però lo
        strumento è la dualità lineare, non Kakutani.
\end{itemize}

% ============================================================
\chapter{Strutture coalizionali, quasi core e least core (L22)}
\label{ch:coal-struct}
% ============================================================

\section{Il core per strutture coalizionali}

Finora si è supposto che si formasse la grande coalizione $N$. In molte
situazioni ciò non accade: i giocatori si organizzano in \emph{più coalizioni
disgiunte}. Una \textbf{struttura coalizionale} è una partizione
$\mathcal{P}$ di $N$ (scriviamo $\mathcal{P}\in\mathcal{P}(N)$ con abuso di
notazione: $\mathcal{P}$ è una famiglia di coalizioni disgiunte la cui unione
è $N$).

\begin{defbox}{Imputazioni rispetto a una struttura coalizionale}
$x\in\mathbb{R}^n$ è un'\textbf{imputazione} per $(N,v)$ rispetto a
$\mathcal{P}$ se è
\begin{itemize}
  \item \textbf{socialmente razionale}: $x(S)=v(S)$ per ogni $S\in\mathcal{P}$
        (ogni coalizione formata distribuisce esattamente il proprio worth: non
        ci sono trasferimenti fra coalizioni diverse);
  \item \textbf{individualmente razionale}: $x_i\ge v(i)$ per ogni $i\in N$.
\end{itemize}
Si scrive $X(\mathcal{P},v)=\{\text{imputazioni rispetto a }\mathcal{P}\}$.
\end{defbox}

\begin{defbox}{Core per una struttura coalizionale}
\[
  C(N,v,\mathcal{P})=\bigl\{x\in\mathbb{R}^n:\;x(S)=v(S)\ \ \forall S\in\mathcal{P},
  \quad x(S)\ge v(S)\ \ \forall S\subseteq N\bigr\}.
\]
Sono le imputazioni (rispetto a $\mathcal{P}$) razionali per ogni possibile
coalizione, anche per quelle che attraversano i blocchi di $\mathcal{P}$.
\end{defbox}

Per $\mathcal{P}=\{N\}$ si ritrova $C(N,v)$.

\section{La superadditività guida le strutture coalizionali plausibili}

\begin{defbox}{Superadditive cover}
Il \textbf{superadditive cover} di $(N,v)$ è il gioco $(N,v^{*})$ con
\[
  v^{*}(S)\;=\;\max_{\mathcal{T}\in\mathcal{P}(S)}\ \sum_{T\in\mathcal{T}} v(T),
\]
dove il massimo è su tutte le partizioni $\mathcal{T}$ di $S$: è il valore
migliore che i membri di $S$ possono ottenere \emph{organizzandosi
liberamente in sottocoalizioni}. $(N,v^{*})$ è superadditivo e $v^{*}\ge v$;
inoltre $v^{*}=v$ se e solo se $(N,v)$ è superadditivo.
\end{defbox}

\begin{thmbox}{Teorema}
Per ogni struttura coalizionale $\mathcal{P}$:
\[
  C(N,v,\mathcal{P})\;=\;C(N,v^{*})\,\cap\,X(\mathcal{P},v).
\]
\end{thmbox}

\textit{Dimostrazione (sketch).} Il punto chiave è:
\[
  \bigl[\,x(S)\ge v(S)\ \forall S\subseteq N\,\bigr]
  \iff
  \bigl[\,x(S)\ge v^{*}(S)\ \forall S\subseteq N\,\bigr].
\]
L'implicazione $\Leftarrow$ è ovvia perché $v^{*}\ge v$. Per $\Rightarrow$: data
una partizione $\mathcal{T}$ di $S$, si ha
$x(S)=\sum_{T\in\mathcal{T}}x(T)\ge\sum_{T\in\mathcal{T}}v(T)$; prendendo il
massimo su $\mathcal{T}$ si ottiene $x(S)\ge v^{*}(S)$. A questo punto
$x\in C(N,v,\mathcal{P})$ significa: (a) $x(S)\ge v^{*}(S)$ per ogni $S$, (b)
$x(S)=v(S)$ per ogni $S\in\mathcal{P}$ (da cui $x(N)=\sum_{T\in\mathcal{P}}v(T)$)
e (c) $x_i\ge v(i)$. Le condizioni (a) e $x(N)=v^{*}(N)$ definiscono $C(N,v^{*})$,
le (b)--(c) definiscono $X(\mathcal{P},v)$. \qed

\begin{corollary}
Sia $\mathcal{P}$ una struttura coalizionale.
\begin{itemize}
  \item[(i)] Se $v^{*}(N)=\sum_{T\in\mathcal{P}}v(T)$, allora $C(N,v,\mathcal{P})=C(N,v^{*})$.
  \item[(ii)] Se $v^{*}(N)>\sum_{T\in\mathcal{P}}v(T)$, allora $C(N,v,\mathcal{P})=\emptyset$.
\end{itemize}
\end{corollary}

\textit{Dimostrazione.} (ii) Se $x\in C(N,v,\mathcal{P})$ allora
$x(N)=\sum_{T\in\mathcal{P}}v(T)<v^{*}(N)$, contro $x(N)\ge v^{*}(N)$ (dimostrato
sopra). (i) Sia $x\in C(N,v^{*})$: allora $x(N)=v^{*}(N)=\sum_{T\in\mathcal{P}}v(T)$
e $x(T)\ge v(T)$ per ogni $T\in\mathcal{P}$; sommando, tutte queste disuguaglianze
devono essere uguaglianze, dunque $x\in X(\mathcal{P},v)$. \qed

\begin{remark}
Conclusione operativa: \emph{sopravvivono solo le strutture coalizionali ottime},
cioè quelle che realizzano il massimo nella definizione di $v^{*}(N)$. Per tutte
le altre il core è vuoto: qualche gruppo di giocatori riorganizzerebbe la
partizione.
\end{remark}

\section{Market games con costi di accordo}

\begin{exbox}{Vino e bottiglie con costi di negoziazione}
$N=P\cup B$ come nel Cap.~\ref{ch:bondareva}, ma ogni coppia produttore/
imbottigliatore che deve mettersi d'accordo costa $c$:
\[
  v(S)=\min\Bigl\{\sum_{i\in S\cap P}\ell_i,\ \sum_{i\in S\cap B}b_i\Bigr\}\;-\;c\,|S\cap P|\,|S\cap B| .
\]
Con $c=1$, $P=\{1,2,3\}$, $\ell=(6,5,4)$, $B=\{4,5\}$, $b=(5,7)$:
\begin{center}
\begin{tabular}{ccc@{\qquad}ccc}
\toprule
$S$ & $v(S)$ & $v^{*}(S)$ & $S$ & $v(S)$ & $v^{*}(S)$\\
\midrule
$14$  & $4$ & $4$ & $1245$  & $7$ & $9$\\
$15$  & $5$ & $5$ & $134$   & $3$ & $4$\\
$145$ & $4$ & $5$ & $135$   & $5$ & $5$\\
$24$  & $4$ & $4$ & $1345$  & $6$ & $8$\\
$25$  & $4$ & $4$ & $234$   & $3$ & $4$\\
$245$ & $3$ & $4$ & $235$   & $5$ & $5$\\
$34$  & $3$ & $3$ & $2345$  & $5$ & $7$\\
$35$  & $3$ & $3$ & $1234$  & $2$ & $4$\\
$345$ & $2$ & $3$ & $1235$  & $4$ & $5$\\
$124$ & $3$ & $4$ & $12345$ & $6$ & $9$\\
$125$ & $5$ & $5$ & altre   & $0$ & $0$\\
\bottomrule
\end{tabular}
\end{center}
Il costo di accordo distrugge la struttura di market game (che è totalmente
bilanciato) e infatti
\[
  C(N,v)=\emptyset,\qquad
  C(N,v^{*})=\{(t,\,t,\,0,\,4-t,\,5-t)\;:\;0\le t\le 1\}.
\]
Si osservi $v^{*}(N)=9>6=v(N)$: conviene spezzarsi. Le uniche partizioni che
realizzano $v^{*}(N)=9$ sono
\[
  \mathcal{P}=\bigl\{\{1,3,5\},\{2,4\}\bigr\}
  \quad\text{e}\quad
  \mathcal{P}=\bigl\{\{2,3,5\},\{1,4\}\bigr\},
\]
(infatti $v(\{1,3,5\})=5$, $v(\{2,4\})=4$ e $v(\{2,3,5\})=5$, $v(\{1,4\})=4$).
Per queste due strutture $C(N,v,\mathcal{P})=C(N,v^{*})\ne\emptyset$; per ogni
altra $\mathcal{P}$ (compresa $\mathcal{P}=\{N\}$) il core è vuoto.
\end{exbox}

\section{Tornare alla grande coalizione: preimputazioni e quasi core}

Se si vuole comunque una soluzione per la grande coalizione anche quando
$C(N,v)=\emptyset$, occorre rilassare qualcosa. Primo passo: rinunciare alla
razionalità individuale.

\begin{defbox}{Preimputazione}
$x\in\mathbb{R}^n$ è una \textbf{preimputazione} per $(N,v)$ se è socialmente
razionale, cioè $x(N)=v(N)$. L'insieme
$X^{0}(N,v)=\{\text{preimputazioni}\}$ è un iperpiano, dunque \emph{sempre non
vuoto}. Con questa notazione
\[
  C(N,v)=\{x\in X^{0}(N,v)\;:\;x(S)\ge v(S)\ \ \forall S\subsetneq N\}.
\]
\end{defbox}

Secondo passo: rilassare le disuguaglianze coalizionali.

\begin{defbox}{Strong $\varepsilon$-core (Shapley--Shubik, 1966)}
Per $\varepsilon\in\mathbb{R}$,
\[
  C_{\varepsilon}(N,v)=\{x\in X^{0}(N,v)\;:\;x(S)\ge v(S)-\varepsilon\ \ \forall S\subsetneq N\}.
\]
\end{defbox}

Interpretazione: $\varepsilon$ è il ``costo fisso'' (o la penale) che una
coalizione deve sostenere per staccarsi; se $\varepsilon>0$ si tollera un pò di
insoddisfazione, se $\varepsilon<0$ si chiede un margine di sicurezza.

\begin{proposition}[Proprietà del quasi core]
\begin{itemize}
  \item[(i)] $C_{0}(N,v)=C(N,v)$;
  \item[(ii)] $C_{\varepsilon_2}(N,v)\subseteq C_{\varepsilon_1}(N,v)$ per $\varepsilon_2<\varepsilon_1$ (monotonia);
  \item[(iii)] $C_{\varepsilon_2}(N,v)\subseteq C(N,v)\subseteq C_{\varepsilon_1}(N,v)$ per $\varepsilon_2<0<\varepsilon_1$;
  \item[(iv)] $C_{\varepsilon}(N,v)\ne\emptyset$ per $\varepsilon$ sufficientemente grande;
  \item[(v)] $C_{\varepsilon}(N,v)=\emptyset$ per $\varepsilon$ sufficientemente piccolo (eventualmente negativo).
\end{itemize}
\end{proposition}

\textit{Dimostrazione (sketch).} (i)--(iii) sono immediate dalla definizione.
(iv) Si prenda una qualsiasi preimputazione $\bar x$ e
$\varepsilon\ge\max_{S\subsetneq N}\{v(S)-\bar x(S)\}$: allora
$\bar x\in C_{\varepsilon}$. (v) Sommando i vincoli relativi a una partizione
$\{S_1,\dots,S_k\}$ si ottiene $v(N)=x(N)\ge\sum_j v(S_j)-k\varepsilon$, che è
falsa per $\varepsilon$ abbastanza negativo. \qed

\begin{defbox}{Weak $\varepsilon$-core (Shapley--Shubik, 1966)}
\[
  \widetilde{C}_{\varepsilon}(N,v)=\{x\in X^{0}(N,v)\;:\;x(S)\ge v(S)-|S|\,\varepsilon\ \ \forall S\subsetneq N\}.
\]
Qui la penale è \emph{pro capite}: $\varepsilon$ per ciascun membro di $S$. Vale
\[
  C_{\varepsilon}(N,v)\subseteq\widetilde{C}_{\varepsilon}(N,v)\subseteq C_{n\varepsilon}(N,v)
  \qquad(\varepsilon\ge 0),
\]
con le inclusioni invertite per $\varepsilon\le 0$ (perché $1\le|S|\le n$).
\end{defbox}

\section{Il least core}

\begin{defbox}{Least core (Maschler--Peleg--Shapley, 1979)}
Sia $I(N,v)=\{\varepsilon\in\mathbb{R}: C_{\varepsilon}(N,v)\ne\emptyset\}$. Per la
monotonia e per i punti (iv)--(v) della proposizione precedente, $I(N,v)$ è una
semiretta chiusa: $I(N,v)=[\bar\varepsilon,+\infty)$ con $\bar\varepsilon>-\infty$.
Il \textbf{least core} è
\[
  LC(N,v)=\bigcap_{\varepsilon\in I(N,v)}C_{\varepsilon}(N,v)=C_{\bar\varepsilon}(N,v).
\]
\end{defbox}

\begin{proposition}
\begin{itemize}
  \item[(i)] $LC(N,v)\ne\emptyset$ (ed è un poliedro compatto convesso);
  \item[(ii)] $C(N,v)\ne\emptyset\implies\bar\varepsilon\le 0$, equivalentemente $LC(N,v)\subseteq C(N,v)$;
  \item[(iii)] $LC(N,v)=\argmin\bigl\{\max\{v(S)-x(S)\,:\,S\subsetneq N\}\;:\;x\in X^{0}(N,v)\bigr\}$.
\end{itemize}
\end{proposition}

\textit{Dimostrazione (sketch).} Si consideri il problema di programmazione
lineare
\[
  \min\ \varepsilon \quad\text{s.t.}\quad x(S)+\varepsilon\ge v(S)\ \ (S\subsetneq N),\qquad x(N)=v(N).
\]
È ammissibile (per (iv)) e limitato inferiormente (per (v)), dunque ha soluzione
ottima: il valore ottimo è $\bar\varepsilon$ e la proiezione dell'insieme delle
soluzioni ottime sulle $x$ è esattamente $C_{\bar\varepsilon}(N,v)$, che quindi è
non vuoto e coincide con l'insieme dei minimizzatori della funzione
$x\mapsto\max_{S\subsetneq N}\{v(S)-x(S)\}$ (funzione convessa, essendo massimo di
funzioni affini). Questo prova (i) e (iii). (ii) Se il core è non vuoto allora
$0\in I(N,v)$, quindi $\bar\varepsilon\le 0$ e
$LC=C_{\bar\varepsilon}\subseteq C_{0}=C$. \qed

\begin{remark}
Il least core è quindi \emph{sempre} definito: se il core è non vuoto ne
seleziona un sottoinsieme (quello ``più robusto'', con il massimo margine di
sicurezza); se è vuoto fornisce comunque le allocazioni che minimizzano la
massima insoddisfazione.
\end{remark}

\begin{exbox}{Least core del market game con costi di accordo}
Riprendiamo la tabella con $c=1$. Il minimo $\varepsilon$ per cui
$C_{\varepsilon}(N,v)\ne\emptyset$ è $\bar\varepsilon=1.5$: nella colonna
$v(S)-\bar\varepsilon$ compaiono i valori $2.5,3.5,2.5,\dots$ e il sistema
$x(S)\ge v(S)-1.5$ con $x(N)=6$ ammette soluzione, mentre per $\varepsilon<1.5$ no
(le coalizioni $\{1,4\},\{1,5\},\{2,4\}$ e $\{1,2,3,4\}$ diventano incompatibili).
Si ottiene
\[
  LC(N,v)=C_{1.5}(N,v)=\{(t,\,t,\,0,\,2.5-t,\,3.5-t)\;:\;0\le t\le 1\}.
\]
Confronto con la sezione precedente: $LC$ è la traslazione del core del
superadditive cover, ``compresso'' di $\bar\varepsilon$ sulle due coalizioni
critiche. Il giocatore 3 (il produttore piccolo) riceve $0$ in tutto il least
core: è quello che aggiunge meno valore.
\end{exbox}

\section{Eccessi e insoddisfazione}

\begin{defbox}{Eccesso di una coalizione}
L'\textbf{eccesso} di $S\subseteq N$ in $x\in\mathbb{R}^n$ è
\[
  e(S,x)=v(S)-x(S).
\]
\begin{itemize}
  \item $e(S,x)\le 0$: la coalizione $S$ è \emph{soddisfatta} (riceve almeno quanto vale);
  \item $e(S,x)>0$: il valore fornisce una \emph{misura di insoddisfazione}.
\end{itemize}
\end{defbox}

Con questa notazione i quasi core si riscrivono in modo trasparente:
\[
  C_{\varepsilon}(N,v)=\{x\in X^{0}(N,v)\;:\;e(S,x)\le\varepsilon\ \ \forall S\subsetneq N\},
\]
e il least core è l'insieme delle preimputazioni che realizzano la
\emph{minima massima insoddisfazione}.

\begin{remark}
Il least core in generale \emph{non} è un punto: nell'esempio precedente è un
segmento. Un arbitro che debba scegliere una singola ripartizione ha bisogno di
un criterio ulteriore: iterare la minimizzazione sui livelli successivi di
insoddisfazione. È l'idea del nucleolo (Cap.~\ref{ch:nucleolo}).
\end{remark}

\begin{defbox}{Vettore ordinato di insoddisfazione}
Sia $\mathcal{P}(N)=\{S_1,\dots,S_{2^n}\}$. Il vettore
\[
  \theta(x)=\bigl(e(S_1,x),\dots,e(S_{2^n},x)\bigr)
\]
è il \textbf{vettore ordinato di insoddisfazione} in $x$ se
$e(S_i,x)\ge e(S_{i+1},x)$ per $i=1,\dots,2^n-1$, cioè se gli eccessi sono
elencati in ordine non crescente.
\end{defbox}

\begin{defbox}{Ordine lessicografico}
Dati $a,b\in\mathbb{R}^m$, si pone $a\ge_{\mathrm{lex}}b$ se
\begin{itemize}
  \item esiste $k$ tale che $a_i=b_i$ per ogni $i<k$ e $a_k>b_k$ (in tal caso $a>_{\mathrm{lex}}b$), oppure
  \item $a=b$.
\end{itemize}
\begin{itemize}
  \item $\ge_{\mathrm{lex}}$ è un \emph{ordine totale} (riflessivo, antisimmetrico, transitivo e connesso);
  \item $>_{\mathrm{lex}}$ \emph{non} è continuo: da $a^h>_{\mathrm{lex}}b$ e $a^h\to\bar a$ non segue $\bar a>_{\mathrm{lex}}b$.
        Esempio: $a^h=(1/h,-1)\to\bar a=(0,-1)$ e $b=(0,0)$: si ha $a^h>_{\mathrm{lex}}b$ per ogni $h$
        ma $\bar a<_{\mathrm{lex}}b$.
\end{itemize}
\end{defbox}

L'obiettivo diventa: \emph{minimizzare lessicograficamente} $\theta(x)$, cioè
abbassare prima la massima insoddisfazione, poi (a parità di questa) la
seconda, e così via.

\section*{In sintesi per l'orale}
\addcontentsline{toc}{section}{In sintesi per l'orale (L22)}
\begin{itemize}
  \item Struttura coalizionale $\mathcal{P}$ = partizione di $N$. Imputazione
        rispetto a $\mathcal{P}$: $x(S)=v(S)$ per ogni blocco $S\in\mathcal{P}$ e
        $x_i\ge v(i)$. Core: in più $x(S)\ge v(S)$ per \emph{tutte} le $S$.
  \item Superadditive cover $v^{*}(S)=\max_{\mathcal{T}}\sum_{T\in\mathcal{T}}v(T)$.
        Teorema: $C(N,v,\mathcal{P})=C(N,v^{*})\cap X(\mathcal{P},v)$. Trucco della
        dimostrazione: $x(S)\ge v(S)\ \forall S\iff x(S)\ge v^{*}(S)\ \forall S$.
  \item Corollario: solo le partizioni \emph{ottime} ($\sum_{T\in\mathcal{P}}v(T)=v^{*}(N)$)
        hanno core non vuoto, e in tal caso $C(N,v,\mathcal{P})=C(N,v^{*})$.
  \item Esempio guida: market game con costi di accordo $c=1$ --- $C(N,v)=\emptyset$,
        $v^{*}(N)=9>6=v(N)$, partizioni ottime $\{135\},\{24\}$ e $\{235\},\{14\}$.
  \item Preimputazione = solo efficienza $x(N)=v(N)$: l'iperpiano $X^{0}$ non è
        mai vuoto (a differenza di $X$ e di $C$).
  \item Strong $\varepsilon$-core: penale fissa $\varepsilon$; weak $\varepsilon$-core:
        penale $|S|\varepsilon$ pro capite. $C_\varepsilon\subseteq\widetilde C_\varepsilon\subseteq C_{n\varepsilon}$ per $\varepsilon\ge0$.
  \item Least core $LC=C_{\bar\varepsilon}$ con $\bar\varepsilon=\min\{\varepsilon: C_\varepsilon\ne\emptyset\}$:
        sempre non vuoto, contenuto nel core quando questo è non vuoto, e uguale a
        $\argmin_x\max_{S}e(S,x)$ (un LP con $2^n-2$ vincoli).
  \item Eccesso $e(S,x)=v(S)-x(S)$: misura di insoddisfazione. $\theta(x)$ =
        eccessi in ordine non crescente; l'ordine lessicografico è totale ma non
        continuo. Minimizzare $\theta$ lessicograficamente porta al nucleolo.
\end{itemize}

% ============================================================
\chapter{Nucleolo e prenucleolo (L23)}
\label{ch:nucleolo}
% ============================================================

\section{Minimizzare l'insoddisfazione: i nucleoli}

Il least core minimizza il \emph{primo} livello di insoddisfazione, ma in
generale non seleziona un unico punto. L'idea di Schmeidler è di iterare:
a parità di massima insoddisfazione si minimizza la seconda, poi la terza, e
così via. Formalmente: si minimizza $\theta(x)$ rispetto all'ordine
lessicografico.

\begin{defbox}{Nucleoli di un gioco (Schmeidler, 1969)}
Dato $K\subseteq\mathbb{R}^n$, il \textbf{nucleolo di $(N,v)$ relativo a $K$} è
\[
  \mathcal{N}(N,v,K)=\{x\in K\;:\;\theta(x)\le_{\mathrm{lex}}\theta(y)\ \ \forall\,y\in K\},
\]
cioè l'insieme dei minimi lessicografici di $\theta$ su $K$.
\end{defbox}

\begin{proposition}[Proprietà generali]
\begin{itemize}
  \item[(i)] $K_2\subseteq K_1\implies \mathcal{N}(N,v,K_1)\cap K_2\subseteq \mathcal{N}(N,v,K_2)$;
  \item[(ii)] se $K_2\subseteq K_1$ e $\emptyset\ne\mathcal{N}(N,v,K_1)\subseteq K_2$, allora
        $\mathcal{N}(N,v,K_1)=\mathcal{N}(N,v,K_2)$;
  \item[(iii)] (covarianza) $\mathcal{N}(N,av+b,aK+b)=a\,\mathcal{N}(N,v,K)+b$ per $a>0$, $b\in\mathbb{R}^n$.
\end{itemize}
\end{proposition}

\textit{Dimostrazione (sketch).} (i) Se $x$ è minimo lessicografico su $K_1$ e
$x\in K_2\subseteq K_1$, a maggior ragione lo è su $K_2$. (ii) Da (i) segue
$\subseteq$; l'altra inclusione perché un minimo su $K_2$ ha lo stesso valore di
$\theta$ del minimo su $K_1$, che sta in $K_2$. (iii) Ponendo $y=ax+b$ e
$w=av+b$ si ha
\[
  e_w(S,y)=w(S)-y(S)=a\,v(S)+b(S)-a\,x(S)-b(S)=a\,e_v(S,x),
\]
quindi $\theta_w(y)=a\,\theta_v(x)$ con $a>0$: l'ordinamento (e quindi l'ordine
lessicografico) è preservato. \qed

\begin{defbox}{Nucleolo e prenucleolo}
\begin{itemize}
  \item $\mathcal{N}(N,v)=\mathcal{N}\bigl(N,v,X(N,v)\bigr)$ è il \textbf{nucleolo} di $(N,v)$
        (minimo lessicografico sulle \emph{imputazioni});
  \item $\mathcal{PN}(N,v)=\mathcal{N}\bigl(N,v,X^{0}(N,v)\bigr)$ è il \textbf{prenucleolo}
        (minimo lessicografico sulle \emph{preimputazioni}).
\end{itemize}
\end{defbox}

\begin{remark}
Poiché $X(N,v)\subseteq X^{0}(N,v)$, dalla proprietà (i) segue: se
$x\in\mathcal{PN}(N,v)$ è individualmente razionale, allora
$x\in\mathcal{N}(N,v)$. Il viceversa è falso: si veda l'Esempio~\ref{ex:pn-non-imp}.
Inoltre né $\mathcal{N}$ né $\mathcal{PN}$ sono \emph{additivi}:
in generale
$[\mathcal{P}]\mathcal{N}(N,v+w)\ne[\mathcal{P}]\mathcal{N}(N,v)+[\mathcal{P}]\mathcal{N}(N,w)$.
È la differenza fondamentale rispetto al valore di Shapley (Cap.~\ref{ch:shapley}).
\end{remark}

\section{Esistenza e unicità}

\begin{thmbox}{Teorema (Schmeidler)}
\begin{itemize}
  \item[(i)] Se $X(N,v)\ne\emptyset$, allora $|\mathcal{N}(N,v)|=1$;
  \item[(ii)] $|\mathcal{PN}(N,v)|=1$ (sempre, perché $X^{0}(N,v)\ne\emptyset$).
\end{itemize}
Nucleolo e prenucleolo sono dunque concetti di soluzione \emph{single-valued}.
\end{thmbox}

\begin{thmbox}{Scelta singola accettabile}
Se $C(N,v)\ne\emptyset$, allora
\[
  \mathcal{N}(N,v)=\mathcal{PN}(N,v)\in C(N,v).
\]
In generale, $\mathcal{PN}(N,v)\in LC(N,v)$ e $LC(N,v)=C_{\bar\varepsilon}(N,v)$ con
$\bar\varepsilon=\theta_1\bigl(\mathcal{PN}(N,v)\bigr)$.
\end{thmbox}

\textit{Dimostrazione (sketch).} Il prenucleolo minimizza in particolare
$\theta_1(x)=\max_{S\subsetneq N}e(S,x)$, quindi appartiene al least core (che è
esattamente l'insieme dei minimizzatori di $\theta_1$, per la Proposizione del
Cap.~\ref{ch:coal-struct}), e il valore ottimo $\bar\varepsilon$ è proprio
$\theta_1(\mathcal{PN})$. Se il core è non vuoto, allora
$LC\subseteq C\subseteq X(N,v)$, quindi il prenucleolo è individualmente
razionale e per la proprietà (ii) dei nucleoli coincide con il nucleolo. \qed

\begin{proposition}[Proprietà aggiuntive]
Valgono per $\mathcal{PN}$ (e per $\mathcal{N}$ quando non vuoto):
\begin{itemize}
  \item (\emph{null player}) se $v(S\cup\{i\})=v(S)$ per ogni $S\subseteq N$, allora
        $\bigl(\mathcal{PN}(N,v)\bigr)_i=0$;
  \item (\emph{symmetry}) se $v(S\cup\{i\})=v(S\cup\{j\})$ per ogni $S\subseteq N$ con
        $i,j\notin S$, allora
        $\bigl(\mathcal{PN}(N,v)\bigr)_i=\bigl(\mathcal{PN}(N,v)\bigr)_j$.
\end{itemize}
\end{proposition}

\section{Come si calcola: lo schema di Maschler}

Il calcolo del prenucleolo si riduce a una \emph{successione di problemi di
programmazione lineare}. Il punto di partenza è la riscrittura
\[
  \min_{x\in K}\ \max\{e(S,x):S\in\mathcal{F}\}
  \qquad\Longleftrightarrow\qquad
  \begin{array}{ll}
    \min & \delta\\
    \text{s.t.} & \delta\ge e(S,x)=v(S)-x(S)\quad S\in\mathcal{F}\\
                & x\in K,
  \end{array}
\]
dove $K$ è un poliedro. Il numero di vincoli è esponenziale in $n$.

\begin{defbox}{Schema di Maschler}
Si pone $\mathcal{F}_1=\mathcal{P}(N)\setminus\{\emptyset,N\}$ e $K_1=X^{0}(N,v)$.
Al passo $k=1,2,\dots$:
\begin{enumerate}
  \item si risolve $\displaystyle \delta_k=\min_{x\in K_k}\max\{e(S,x):S\in\mathcal{F}_k\}$
        (LP come sopra) e si pone
        $K_{k+1}=\{x\in K_k: e(S,x)\le\delta_k\ \forall S\in\mathcal{F}_k\}$;
  \item si individuano le coalizioni \emph{fissate}
        $\mathcal{D}_k=\{S\in\mathcal{F}_k: e(S,x)=\delta_k\ \ \forall x\in K_{k+1}\}$,
        cioè quelle il cui eccesso non è più migliorabile, e si pone
        $\mathcal{F}_{k+1}=\mathcal{F}_k\setminus\mathcal{D}_k$;
  \item ci si ferma quando $K_{k+1}$ è un singoletto (o $\mathcal{F}_{k+1}=\emptyset$).
\end{enumerate}
Il punto finale è $\mathcal{PN}(N,v)$. La dimensione di $K_k$ cala di almeno 1 a
ogni passo, quindi lo schema termina in al più $n$ iterazioni. Il primo passo
produce $\delta_1=\bar\varepsilon$ e $K_2=LC(N,v)$.
\end{defbox}

\section{Esempi svolti}

\begin{exbox}{Un caso semplice di bancarotta}
Un'impresa fallisce; l'attivo è $p$ e ci sono due creditori con crediti
$d_1,d_2$ tali che $d_1+d_2>p$ (l'attivo non basta). Il gioco associato è
\[
  v(S)=\begin{cases}
    \max\{p-d_{-i},\,0\} & \text{se } S=\{i\},\\
    p & \text{se } S=N=\{1,2\},
  \end{cases}
\]
cioè un creditore da solo può contare su ciò che resta \emph{dopo} aver
soddisfatto interamente l'altro. Allora
\[
  \mathcal{N}(N,v)=\mathcal{PN}(N,v)
  =\Bigl(\frac{v(\{1,2\})+v(1)-v(2)}{2},\ \frac{v(\{1,2\})+v(2)-v(1)}{2}\Bigr)\in C(N,v).
\]
\textit{Derivazione (covarianza + simmetria).} Si ponga $b=(v(1),v(2))$ e
$w(S)=v(S)-b(S)$: allora $w(1)=w(2)=0$ e $w(N)=v(N)-v(1)-v(2)\ge 0$. Nel gioco
$w$ i due giocatori sono \emph{simmetrici}, quindi per la proprietà di simmetria
$\mathcal{PN}(w)=\bigl(\tfrac{w(N)}{2},\tfrac{w(N)}{2}\bigr)$; per covarianza
$\mathcal{PN}(v)=\mathcal{PN}(w)+b$, che è la formula sopra. Se
$v(1)+v(2)<v(\{1,2\})$ il core ha la cardinalità del continuo (è un segmento):
il nucleolo ne seleziona il punto medio. In pratica ogni creditore riceve la
parte non contesa $\max\{p-d_{-i},0\}$ e il resto viene diviso a metà.
\end{exbox}

\begin{exbox}{Nucleolo del gioco simmetrico a tre giocatori}
Sia $n=3$, $v(i)=0$, $v(S)=1$ per $|S|=2$ e $v(N)=1$. Il core è vuoto (da
$x_i+x_j\ge 1$ segue $x_k\le 0$ per ogni $k$, incompatibile con $x(N)=1$).
Gli eccessi in $x$ sono $e(\{i\},x)=-x_i$, $e(\{i,j\},x)=1-(1-x_k)=x_k$,
$e(\emptyset,x)=e(N,x)=0$. Dunque
\[
  \theta_1(x)=\max\{0,\ \max_i x_i,\ \max_i(-x_i)\}\ \ge\ \max_i x_i\ \ge\ \tfrac13,
\]
con uguaglianza solo per $x=(\tfrac13,\tfrac13,\tfrac13)$. Quindi
$\mathcal{N}(N,v)=\mathcal{PN}(N,v)=\{(\tfrac13,\tfrac13,\tfrac13)\}$, e
$\bar\varepsilon=\tfrac13>0$ (coerente con il core vuoto).
\end{exbox}

\begin{example}[Il prenucleolo può non essere un'imputazione]
\label{ex:pn-non-imp}
Sia $n=3$ con $v(\{1,2\})=1$ e $v(S)=0$ per ogni altra $S$ (in particolare
$v(N)=0$). Allora $C(N,v)=\emptyset$ e $X(N,v)=\{(0,0,0)\}$, dunque
$\mathcal{N}(N,v)=\{(0,0,0)\}$. Sulle preimputazioni
$X^{0}=\{x:x_1+x_2+x_3=0\}$ gli eccessi sono
\[
  e(\{1\},x)=-x_1,\quad e(\{2\},x)=-x_2,\quad e(\{3\},x)=x_1+x_2,
\]
\[
  e(\{1,2\},x)=1-(x_1+x_2),\quad e(\{1,3\},x)=x_2,\quad e(\{2,3\},x)=x_1,\quad
  e(\emptyset,x)=e(N,x)=0 .
\]
Posto $s=x_1+x_2$: $\theta_1(x)\ge\max\{s,1-s\}\ge\tfrac12$, con uguaglianza sse
$s=\tfrac12$ e $|x_1|,|x_2|\le\tfrac12$. Al secondo livello resta da minimizzare
$\max\{0,|x_1|,|x_2|\}$ con $x_1+x_2=\tfrac12$: il minimo è $\tfrac14$ per
$x_1=x_2=\tfrac14$. Dunque
\[
  \mathcal{PN}(N,v)=\bigl(\tfrac14,\tfrac14,-\tfrac12\bigr)\ \ne\ (0,0,0)=\mathcal{N}(N,v),
\]
e il prenucleolo \emph{non è un'imputazione} (viola $x_3\ge v(3)=0$). Questo
mostra che l'ipotesi ``$\mathcal{PN}$ individualmente razionale'' è essenziale
per l'identificazione dei due concetti.
\end{example}

\begin{exbox}{Nucleolo del market game con costi di accordo}
Riprendiamo il gioco con $c=1$ del Cap.~\ref{ch:coal-struct}. \emph{Primo LP}:
$\bar\varepsilon=1.5$ e
\[
  LC(N,v)=C_{1.5}(N,v)=\{x_t=(t,t,0,2.5-t,3.5-t)\;:\;0\le t\le 1\}.
\]
\emph{Secondo LP}: si calcolano gli eccessi in $x_t$ (tabella delle slide). Le
coalizioni $\{1,4\},\{1,5\},\{2,4\},\{1,3,5\},\{2,3,5\}$ hanno eccesso
costantemente $1.5$: sono \emph{fissate} e vengono rimosse. Fra le rimanenti,
quelle rilevanti hanno eccessi
\[
  e(\{3,4\},x_t)=t+0.5,\qquad e(\{1,2,5\},x_t)=1.5-t,\qquad e(\{1,3,4,5\},x_t)=t,
\]
(le altre sono dominate da queste su $[0,1]$). Si risolve
\[
  \min\bigl\{\max\{t+0.5,\ 1.5-t,\ t\}\;:\;0\le t\le 1\bigr\},
\]
il cui minimo si ha in $t=0.5$ (dove $t+0.5=1.5-t=1$). Dunque
\[
  \mathcal{PN}(N,v)=\mathcal{N}(N,v)=\{(0.5,\;0.5,\;0,\;2,\;3)\}.
\]
Si noti che il nucleolo esiste ed è unico anche se $C(N,v)=\emptyset$: è
proprio la situazione in cui serve.
\end{exbox}

\begin{exbox}{Nucleolo del gioco dei risparmi aggregati}
Per il gioco a 4 giocatori del Cap.~\ref{ch:tu-core}, gli eccessi lungo il core
$x_t=(t,55-t,35,30)$ sono
\[
  e(\{1\})=-t,\quad e(\{2\})=t-20,\quad e(\{1,3\})=10-t,\quad e(\{1,4\})=-t,
\]
\[
  e(\{2,3\})=e(\{2,4\})=e(\{2,3,4\})=t-20,\quad e(\{1,3,4\})=10-t,
\]
mentre $\{3\},\{4\},\{1,2\},\{3,4\},\{1,2,3\},\{1,2,4\}$ hanno eccesso
identicamente $0$.
\emph{Primo LP}: $\bar\varepsilon=0$ (gli eccessi nulli non si possono abbassare) e
$LC(N,v)=C_{0}(N,v)=C(N,v)=\{x_t:10\le t\le 20\}$.
\emph{Secondo LP}: rimosse le coalizioni fissate a $0$, si minimizza
$\max\{-t,\ t-20,\ 10-t\}$ su $[10,20]$; il minimo è $-5$, raggiunto in $t=15$
(dove $t-20=10-t=-5$). Dunque
\[
  \mathcal{PN}(N,v)=\mathcal{N}(N,v)=\{(15,\,40,\,35,\,30)\}\in C(N,v).
\]
Il nucleolo sceglie il \emph{punto medio} del segmento-core, esattamente come
suggerisce l'intuizione di equità dell'arbitro.
\end{exbox}

\section*{In sintesi per l'orale}
\addcontentsline{toc}{section}{In sintesi per l'orale (L23)}
\begin{itemize}
  \item $\mathcal{N}(N,v,K)$ = minimi lessicografici di $\theta(x)$ su $K$.
        Nucleolo: $K=X(N,v)$ (imputazioni); prenucleolo: $K=X^{0}(N,v)$
        (preimputazioni).
  \item \textbf{Teorema (Schmeidler)}: $|\mathcal{PN}|=1$ sempre; $|\mathcal{N}|=1$
        se $X(N,v)\ne\emptyset$. Sono quindi concetti single-valued: è ciò che
        serve a un arbitro che deve scegliere una sola ripartizione.
  \item Se $C(N,v)\ne\emptyset$: $\mathcal{N}=\mathcal{PN}\in C$. Sempre:
        $\mathcal{PN}\in LC$ e $\bar\varepsilon=\theta_1(\mathcal{PN})$.
  \item Il prenucleolo soddisfa null player, simmetria e covarianza, ma
        \emph{non} l'additività: è la differenza chiave con Shapley.
  \item Se il prenucleolo è individualmente razionale allora coincide con il
        nucleolo; non vale il viceversa (esempio $v(\{1,2\})=1$, tutto il resto
        $0$: $\mathcal{PN}=(\tfrac14,\tfrac14,-\tfrac12)$, $\mathcal{N}=(0,0,0)$).
  \item Calcolo: schema di Maschler = successione di LP $\min\delta$ s.t.
        $\delta\ge e(S,x)$, fissando a ogni passo le coalizioni il cui eccesso non
        è più migliorabile. Al più $n$ iterazioni, ma $O(2^n)$ vincoli per LP.
  \item Esempi da saper rifare: bancarotta a 2 creditori (covarianza $+$
        simmetria); risparmi aggregati
        $\to(15,40,35,30)$; market game con costi $\to(0.5,0.5,0,2,3)$.
  \item Collegamento con il bargaining: il nucleolo è un arbitrato che
        ``pareggia l'insoddisfazione'', concettualmente vicino a Kalai--Smorodinsky
        (equalizzare rapporti) più che a Nash (prodotto dei guadagni); e come la
        soluzione di Nash è \emph{un punto} e non un insieme, a differenza del
        core.
\end{itemize}

% ============================================================
\chapter{Assiomatica, valore di Shapley e indici di potere (L24)}
\label{ch:shapley}
% ============================================================

\section{Concetti di soluzione: il punto di vista assiomatico}

Fin qui abbiamo costruito concetti di soluzione ``dal basso'', partendo da
ragionamenti di stabilità. Il punto di vista assiomatico è opposto: si elencano
le \emph{proprietà desiderabili} e ci si chiede quale (se esiste) sia il concetto
di soluzione che le soddisfa.

\begin{defbox}{Concetti di soluzione}
Sia $N$ finito, $V_N=\{v:\mathcal{P}(N)\to\mathbb{R}\;:\;v(\emptyset)=0\}$ l'insieme
di tutte le funzioni caratteristiche su $N$ e
$\mathcal{G}=\{(N,v): N\text{ finito},\,v\in V_N\}$ la classe dei giochi
coalizionali TU. Si distinguono
\begin{itemize}
  \item $\Phi:\mathcal{G}\rightrightarrows\mathbb{R}^N$ concetto di soluzione
        \textbf{set-valued} (es.\ il core, il least core, gli stable sets);
  \item $\varphi:\mathcal{G}\to\mathbb{R}^N$ concetto di soluzione
        \textbf{single-valued} (es.\ il prenucleolo).
\end{itemize}
Fissato $N=\{1,\dots,n\}$ si scrive semplicemente $\varphi:V_N\to\mathbb{R}^n$.
\end{defbox}

\begin{defbox}{Proprietà desiderabili per $\varphi$}
\begin{description}
  \item[Efficienza (razionalità sociale)] $\displaystyle\sum_{i\in N}\varphi_i(v)=v(N)$
        (cioè $\varphi(v)$ è una preimputazione).
  \item[Simmetria] se $v(S\cup\{i\})=v(S\cup\{j\})$ per ogni $S\subseteq N$ con
        $i,j\notin S$, allora $\varphi_i(v)=\varphi_j(v)$.
  \item[Null player] se $v(S\cup\{i\})=v(S)$ per ogni $S\subseteq N$, allora
        $\varphi_i(v)=0$.
  \item[Additività] $\varphi(v+w)=\varphi(v)+\varphi(w)$.
  \item[Covarianza per equivalenza strategica] $\varphi(av+b)=a\varphi(v)+b$ per
        ogni $a>0$, $b\in\mathbb{R}^n$.
\end{description}
\end{defbox}

Simmetria e null player sono requisiti di \emph{equità} (giocatori
indistinguibili ricevono lo stesso; chi non contribuisce mai non riceve nulla);
l'additività è un requisito di \emph{coerenza}: se due progetti indipendenti
vengono valutati insieme, i compensi si sommano.

\section{Tentativi preliminari}

\begin{example}[Tre candidati naturali, tutti insufficienti]
\begin{itemize}
  \item $\varphi_i(v)=v(i)$: soddisfa simmetria, null player, additività e
        covarianza; \emph{manca l'efficienza} (a meno che il gioco sia
        strategicamente equivalente a un gioco $0$-$0$ normalizzato).
  \item $\varphi_i(v)=\max\{v(S\cup\{i\})-v(S)\;:\;S\subseteq N,\ i\notin S\}$
        (il contributo marginale massimo): soddisfa simmetria, null player e
        covarianza; \emph{mancano efficienza e additività} (il massimo di una
        somma non è la somma dei massimi).
  \item $\varphi_i(v)=v(\{1,2,\dots,i\})-v(\{1,2,\dots,i-1\})$ (il contributo
        marginale rispetto a un ordine \emph{fissato}): soddisfa efficienza (la
        somma telescopica dà $v(N)$), null player, additività e covarianza;
        \emph{manca la simmetria}, perché privilegia arbitrariamente l'ordine
        $1,2,\dots,n$.
\end{itemize}
\end{example}

L'ultimo tentativo suggerisce il rimedio: \emph{mediare su tutti gli ordini
possibili}, ristabilendo così la simmetria senza perdere le altre proprietà.

\section{Il valore di Shapley}

\begin{thmbox}{Valore di Shapley (Shapley, 1953)}
Esiste un unico concetto di soluzione single-valued che soddisfa
\emph{efficienza}, \emph{simmetria}, \emph{null player} e \emph{additività},
ed è
\[
  \Sh_i(v)=\frac{1}{n!}\sum_{\pi\in\Pi_N}\bigl[v\bigl(P_i(\pi)\cup\{i\}\bigr)-v\bigl(P_i(\pi)\bigr)\bigr],
\]
dove $\Pi_N$ è l'insieme delle permutazioni di $N$ e
$P_i(\pi)=\{j\in N:\pi(j)<\pi(i)\}$ è l'insieme dei predecessori di $i$ in $\pi$.
\end{thmbox}

Interpretazione: i giocatori entrano nella stanza in ordine casuale (uniforme);
ciascuno riceve il valore che \emph{aggiunge} al gruppo già presente; $\Sh_i(v)$
è il contributo marginale atteso di $i$.

\begin{defbox}{Formulazione alternativa}
Raggruppando le permutazioni con lo stesso insieme di predecessori, e usando
\[
  \bigl|\{\pi\in\Pi_N: P_i(\pi)=S\}\bigr| = |S|!\,(n-|S|-1)!,
\]
si ottiene
\[
  \Sh_i(v)=\frac{1}{n!}\sum_{S\subseteq N\setminus\{i\}}|S|!\,(n-|S|-1)!\,\bigl[v(S\cup\{i\})-v(S)\bigr].
\]
\end{defbox}

\textit{Dimostrazione (sketch del teorema).}
\emph{La formula soddisfa gli assiomi.} Efficienza: fissata $\pi$, la somma su $i$
dei contributi marginali è telescopica e vale $v(N)$; si media su $\pi$.
Additività: la formula è lineare in $v$. Null player e simmetria: immediati dalla
formula (nel secondo caso si accoppiano le permutazioni scambiando $i$ e $j$).

La dimostrazione dell'unicità non è stata svolta a lezione: all'orale basta
sapere che i quattro assiomi \emph{determinano} $\varphi$ e che $\Sh$ li soddisfa.
Il valore di Shapley è, inoltre, covariante per equivalenza strategica.

\section{Caratterizzazione alternativa: marginalità (Young)}

\begin{defbox}{Monotonicità e marginalità}
\begin{description}
  \item[Monotonicità dei contributi marginali] se
        $v(S\cup\{i\})-v(S)\ge w(S\cup\{i\})-w(S)$ per ogni $S\subseteq N$ con
        $i\notin S$, allora $\varphi_i(v)\ge\varphi_i(w)$.
  \item[Marginalità] se $v(S\cup\{i\})-v(S)= w(S\cup\{i\})-w(S)$ per ogni
        $S\subseteq N$ con $i\notin S$, allora $\varphi_i(v)=\varphi_i(w)$.
\end{description}
La marginalità dice che il compenso di $i$ dipende \emph{solo} dai suoi
contributi marginali. Ovviamente monotonicità $\Rightarrow$ marginalità, e il
valore di Shapley soddisfa entrambe.
\end{defbox}

\begin{proposition}
Efficienza $+$ simmetria $+$ marginalità $\implies$ null player.
\end{proposition}

\begin{thmbox}{Caratterizzazione di Young (1985)}
Il valore di Shapley è l'unico concetto di soluzione single-valued che soddisfa
\emph{efficienza}, \emph{simmetria} e \emph{marginalità}.
\end{thmbox}

\begin{remark}[Shapley vs.\ nucleolo]
Il valore di Shapley è additivo ma \emph{non} garantisce di stare nel core,
neppure quando il core è non vuoto; il nucleolo sta nel core ogni volta che
questo è non vuoto ma non è additivo. Sono due risposte diverse alla stessa
domanda: Shapley misura il \emph{contributo}, il nucleolo la
\emph{stabilita'/insoddisfazione}.
\end{remark}

\section{Esempi di calcolo}

\begin{exbox}{Gloves game: Shapley a mano}
$v(\{1,3\})=v(\{2,3\})=v(N)=1$, tutto il resto $0$. Le $3!=6$ permutazioni e i
contributi marginali:
\begin{center}
\begin{tabular}{lccc}
\toprule
$\pi$ & giocatore 1 & giocatore 2 & giocatore 3\\
\midrule
$123$ & $0$ & $0$ & $1$\\
$132$ & $0$ & $0$ & $1$\\
$213$ & $0$ & $0$ & $1$\\
$231$ & $0$ & $0$ & $1$\\
$312$ & $1$ & $0$ & $0$\\
$321$ & $0$ & $1$ & $0$\\
\bottomrule
\end{tabular}
\end{center}
Dunque $\Sh(v)=(\tfrac16,\tfrac16,\tfrac23)$. Confronto con il core, che era
$\{(0,0,1)\}$: il valore di Shapley \emph{non} appartiene al core e riconosce ai
possessori del guanto destro un valore positivo (in un quarto dei casi arrivano
prima del giocatore 3). È la differenza fra ``potere contrattuale'' (core) e
``contributo medio'' (Shapley).
\end{exbox}

\begin{exbox}{Shapley per il gioco dei risparmi aggregati}
Per il gioco a 4 giocatori del Cap.~\ref{ch:tu-core}:
\[
  \Sh(v)=\bigl(\tfrac{35}{3},\ \tfrac{125}{3},\ \tfrac{110}{3},\ 30\bigr)
  \approx(11.\overline{6},\ 41.\overline{6},\ 36.\overline{6},\ 30),
\]
da confrontare con
$\mathcal{N}(N,v)=(15,40,35,30)$ e $C(N,v)=\{(t,55-t,35,30):10\le t\le 20\}$.
Il valore di Shapley \emph{non} sta nel core (assegna $36.\overline{6}$ al
giocatore 3, mentre nel core deve essere esattamente $35$): assegna al giocatore 3
un premio per i suoi contributi marginali elevati, che però la coalizione
$\{1,2,4\}$ non è disposta a concedere.
\end{exbox}

\begin{exbox}{Shapley per il market game con costi di accordo}
Per il gioco a 5 giocatori con $c=1$:
\[
  \Sh(v)=(1.15,\ 0.7\overline{3},\ 0.2\overline{3},\ 1.4,\ 2.48\overline{3}),
  \qquad \mathcal{N}(N,v)=(0.5,\ 0.5,\ 0,\ 2,\ 3).
\]
Entrambi sommano a $v(N)=6$. Qui il core è vuoto: il valore di Shapley è comunque
definito, e assegna un valore positivo anche al giocatore 3, che il nucleolo
lascia a zero. Il grande imbottigliatore (giocatore 5) resta il più pagato con
entrambi i criteri.
\end{exbox}

\section{Indici di potere}

\begin{defbox}{Giochi semplici e monotoni}
$(N,v)$ è \textbf{semplice} se $v(S)\in\{0,1\}$ per ogni $S\subseteq N$.
Un gioco semplice è \textbf{monotono} se $v(S)=1$ implica $v(T)=1$ per ogni
$T\supseteq S$. Una coalizione $S$ è \textbf{vincente} se $v(S)=1$,
\textbf{perdente} se $v(S)=0$.
\end{defbox}

Sono il modello standard delle \emph{procedure di voto}: $v(S)=1$ significa che
$S$ ha i voti per far passare una delibera. Per questi giochi il problema della
ripartizione diventa il problema di misurare il \emph{potere} di ciascun votante.
Si definisce l'insieme degli \emph{swing} di $i$:
\[
  W(i)=\{S\subseteq N\;:\;S\ \text{perdente},\ S\cup\{i\}\ \text{vincente}\}.
\]

\begin{defbox}{Shapley--Shubik power index (1954)}
Per un gioco semplice monotono,
\[
  \mathrm{SSPI}_i(v)=\sum_{S\in W(i)}\frac{|S|!\,(n-|S|-1)!}{n!}.
\]
Cioè: \emph{lo Shapley--Shubik power index è il valore di Shapley di un gioco
semplice monotono}. Interpretazione: la probabilità che $i$ sia il votante
\emph{pivotale} in un ordine casuale.
\end{defbox}

\begin{defbox}{Banzhaf power index (Penrose 1946, Banzhaf 1965, Coleman 1971)}
Per un gioco semplice,
\[
  \mathrm{BPI}_i(v)=\frac{|W(i)|}{2^{n-1}} .
\]
Cioè la frazione delle coalizioni di $N\setminus\{i\}$ per cui $i$ è pivotale,
\emph{senza} pesare per la lunghezza della coalizione. Il BPI soddisfa la
proprietà di null player, ma non è efficiente ($\sum_i\mathrm{BPI}_i$ in
generale $\ne 1$): per confrontare i giocatori lo si normalizza dividendo per la
somma.
\end{defbox}

\begin{exbox}{Procedure di voto elementari}
\textbf{Unanimità}: $v_u(S)=1$ sse $S=N$. Allora $W(i)=\{N\setminus\{i\}\}$ e
\[
  \mathrm{SSPI}_i(v_u)=\frac1n,\qquad \mathrm{BPI}_i(v_u)=\frac{1}{2^{n-1}}
\]
(stessa percentuale relativa: entrambi attribuiscono a tutti lo stesso potere).

\textbf{Dittatura}: $v_d(S)=1$ sse $k\in S$. Allora
\[
  \mathrm{SSPI}_k(v_d)=\mathrm{BPI}_k(v_d)=1,\qquad
  \mathrm{SSPI}_i(v_d)=\mathrm{BPI}_i(v_d)=0\quad (i\ne k).
\]
\end{exbox}

\begin{exbox}{Consiglio di Sicurezza dell'ONU}
$n=15$, $P\subseteq N$ è l'insieme dei 5 membri permanenti (diritto di veto):
\[
  v_{sc}(S)=\begin{cases}1 & \text{se } P\subseteq S \text{ e } |S|\ge 9,\\ 0&\text{altrimenti.}\end{cases}
\]
\emph{Conteggio degli swing per un membro non permanente $i$}: $S$ deve contenere
tutti e 5 i permanenti ed essere perdente con $S\cup\{i\}$ vincente, cioè
$|S|=8$; si scelgono 3 non permanenti fra i 9 diversi da $i$:
$|W(i)|=\binom93=84$. Quindi
\[
  \mathrm{SSPI}_i(v_{sc})=84\cdot\frac{8!\,6!}{15!}\approx 0.00186\quad (i\notin P),
\]
\[
  \mathrm{SSPI}_i(v_{sc})=\frac{1-10\cdot 0.00186}{5}\approx 0.19628\quad (i\in P),
\]
dove per i permanenti si è usata l'efficienza. Per il Banzhaf, le slide riportano
\[
  \mathrm{BPI}_i(v_{sc})\approx 0.00513\ (i\notin P),\qquad
  \mathrm{BPI}_i(v_{sc})\approx 0.01282\ (i\in P),
\]
corrispondenti (dopo normalizzazione) a circa il $4.5\%$ e all'$11\%$ del potere
complessivo per giocatore. Il messaggio è comunque netto: un membro permanente
vale circa 100 volte un membro non permanente secondo Shapley--Shubik.
\end{exbox}

\begin{remark}[L'ultima slide del corso]
Il corso si ferma qui, ma la teoria dei giochi comprende molto altro: giochi
bayesiani, aste, mechanism design, giochi dinamici e differenziali, evolutivi,
mean-field, combinatori, razionalità limitata, cake cutting e fair division,
raffinamenti del NE, giochi coalizionali NTU e con esternalità, giochi ripetuti
e di segnalazione.
\end{remark}

\section*{In sintesi per l'orale}
\addcontentsline{toc}{section}{In sintesi per l'orale (L24)}
\begin{itemize}
  \item Cinque assiomi: efficienza, simmetria, null player, additività,
        covarianza. I tre tentativi ingenui ($v(i)$; contributo marginale massimo;
        contributo marginale rispetto a un ordine fisso) falliscono
        rispettivamente su efficienza; efficienza $+$ additività; simmetria.
  \item \textbf{Teorema di Shapley (1953)}: unico $\varphi$ con efficienza $+$
        simmetria $+$ null player $+$ additività, cioè il contributo marginale
        medio su tutte le $n!$ permutazioni; formula equivalente con pesi
        $|S|!(n-|S|-1)!/n!$.
  \item \textbf{Young (1985)}: efficienza $+$ simmetria $+$ marginalità
        caratterizzano ugualmente $\Sh$; e efficienza $+$ simmetria $+$
        marginalità $\Rightarrow$ null player.
  \item Il valore di Shapley è covariante ma \emph{può stare fuori dal core}
        anche quando il core è non vuoto (esempi: gloves game
        $\Sh=(\tfrac16,\tfrac16,\tfrac23)$ contro $C=\{(0,0,1)\}$; risparmi
        aggregati $\Sh=(11.\overline6,41.\overline6,36.\overline6,30)\notin C$).
  \item Giochi semplici monotoni = procedure di voto. Swing:
        $W(i)=\{S$ perdente$:S\cup\{i\}$ vincente$\}$.
        $\mathrm{SSPI}_i=\sum_{S\in W(i)}|S|!(n-|S|-1)!/n!$ (= Shapley del gioco
        semplice), $\mathrm{BPI}_i=|W(i)|/2^{n-1}$.
  \item Unanimità: $\mathrm{SSPI}_i=1/n$, $\mathrm{BPI}_i=1/2^{n-1}$; dittatura:
        entrambi $(0,\dots,1,\dots,0)$; Consiglio di Sicurezza ONU: permanente
        $\approx 0.196$ contro non permanente $\approx 0.00186$ (SSPI).
  \item \textbf{Quando usare cosa}: il \emph{core} risponde a ``quali accordi sono
        stabili?'' (insieme, spesso vuoto o troppo grande); il \emph{least core}
        e il \emph{nucleolo} rispondono a ``come arbitrare minimizzando
        l'insoddisfazione?'' (sempre definiti, nel core se questo è non vuoto);
        il \emph{valore di Shapley} risponde a ``quanto ha contribuito ciascuno?''
        (sempre definito, additivo, ma non necessariamente stabile). Rispetto al
        bargaining: Shapley è l'analogo assiomatico della soluzione di Nash
        (assiomi $+$ unicità), il nucleolo è più vicino a un arbitrato
        egualitario; rispetto al NE, il core generalizza la non-deviazione da
        singoli a coalizioni.
\end{itemize}


% ============================================================
\appendix
% ============================================================
\chapter{Trappole formali: come si dice giusto alla lavagna}

Elenco delle formulazioni che a voce si sbagliano più spesso. A sinistra la versione
imprecisa, a destra quella da scrivere.

\begin{enumerate}[leftmargin=*]
  \item \textbf{Dominio dell'utilità.}
        \emph{Sbagliato}: ``$u_i$ è definita sulle strategie di $i$''.
        \emph{Giusto}: $u_i:S_1\times\cdots\times S_n\to\mathbb{R}$, dipende dal
        profilo completo; è proprio questo che rende il gioco interdipendente.
  \item \textbf{Strategie pure in forma estesa.}
        \emph{Sbagliato}: ``sono gli archi che il giocatore controlla''.
        \emph{Giusto}: una strategia è una funzione dagli insiemi di informazione
        del giocatore alle azioni, quindi il loro numero è il \emph{prodotto} delle
        azioni disponibili nei suoi insiemi di informazione.
  \item \textbf{Equilibrio di Nash.}
        \emph{Sbagliato}: ``nessuno può migliorare'' con disuguaglianza stretta.
        \emph{Giusto}: $u_i(s_i^*,s_{-i}^*)\ge u_i(s_i,s_{-i}^*)$ per ogni $i$ e
        ogni $s_i$: disuguaglianza debole, deviazione unilaterale, gli altri fermi.
  \item \textbf{Utilità attesa vNM.}
        \emph{Sbagliato}: ``l'utilità è unica''.
        \emph{Giusto}: unica a meno di trasformazioni affini positive $au+b$ con
        $a>0$; serve per dare senso alle lotterie generate dalle strategie miste.
  \item \textbf{Indifferenza nel calcolo del misto.}
        \emph{Sbagliato}: ``impongo che io sia indifferente e trovo la mia mista''.
        \emph{Giusto}: la mia indifferenza determina la mista dell'\emph{avversario};
        la mia esce dall'indifferenza dell'avversario.
  \item \textbf{IESDS contro IEWDS.}
        \emph{Giusto}: IESDS preserva tutti i NE ed è indipendente dall'ordine;
        IEWDS dipende dall'ordine e può eliminare equilibri.
  \item \textbf{Never best response.}
        \emph{Giusto}: non esiste \emph{nessuna} credenza mista sulle strategie
        altrui a cui $s_i$ sia best response. Il quantificatore è su tutte le
        credenze, non solo sulle pure.
  \item \textbf{Potenziale.}
        \emph{Giusto}: esatto $=$ le differenze di utilità \emph{coincidono} con
        quelle del potenziale; ordinale $=$ hanno solo lo stesso segno.
        Monderer--Shapley: basta verificare i cicli di lunghezza 4.
  \item \textbf{Sottogioco.}
        \emph{Sbagliato}: ``un qualunque sottoalbero''.
        \emph{Giusto}: sottoalbero radicato in un nodo singoletto che contiene
        \emph{interi} gli insiemi di informazione che incontra.
  \item \textbf{Alpha-beta.}
        \emph{Giusto}: $\alpha$ è il minimo garantito a MAX, $\beta$ il massimo
        concesso a MIN, si taglia quando $\alpha\ge\beta$; il valore resta esatto e
        la complessità scende a $O(b^{d/2})$ con ordinamento ottimo.
  \item \textbf{Problema di bargaining.}
        \emph{Giusto}: coppia $(U,u^*)$ con $U\subseteq\mathbb{R}^2$ compatto,
        convesso e con almeno un punto che domina il disaccordo $u^*$; punto utopia
        $m(U)=(\max\{u_1\},\max\{u_2\})$ sui punti individualmente razionali.
        Nash massimizza $(u_1-u^*_1)(u_2-u^*_2)$; Kalai--Smorodinsky prende
        l'intersezione della frontiera con il segmento $u^*\to m(U)$ e sostituisce
        l'IIA con la monotonia debole.
  \item \textbf{Rubinstein.}
        \emph{Giusto}: SPE unico stazionario, quota del proponente
        $x_1=(1-\delta_2)/(1-\delta_1\delta_2)$, ricavata dalle due condizioni di
        indifferenza del rispondente; per $\delta\to 1$ tende a metà, cioè alla
        soluzione di Nash (Nash program).
  \item \textbf{Core e imputazioni.}
        \emph{Giusto}: imputazione $=$ efficienza $x(N)=v(N)$ \emph{più} razionalità
        individuale $x_i\ge v(i)$; core $=$ imputazioni con $x(S)\ge v(S)$ per ogni
        $S$. Bondareva--Shapley: core non vuoto $\iff$ gioco bilanciato (duale dell'LP).
  \item \textbf{Nucleolo.}
        \emph{Sbagliato}: ``minimizza l'eccesso massimo''. Quello è il \emph{least core}.
        \emph{Giusto}: il nucleolo minimizza in ordine \emph{lessicografico} l'intero
        vettore degli eccessi ordinato in modo decrescente; è unico e, se il core non
        è vuoto, vi appartiene.
  \item \textbf{Valore di Shapley.}
        \emph{Giusto}: contributo marginale medio su tutte le $n!$ permutazioni,
        $\phi_i=\sum_{S\not\ni i}\frac{|S|!\,(n-|S|-1)!}{n!}[v(S\cup i)-v(S)]$;
        caratterizzato da efficienza, simmetria, dummy player e additività.
        Shapley--Shubik conta le permutazioni in cui si è pivot, Banzhaf conta le
        coalizioni critiche su $2^{n-1}$.
\end{enumerate}

\backmatter
\end{document}
